% 2026-10-25 draft: where the straddle strategy's profits come from (Section 7). Not \input by paper_rehaul.tex.
% Numbers: every number is a macro -- \rhed... from writeup/generated/rehaul_econ_decomp_numbers.tex (generator
% writeup/make_rehaul_econ_decomp.py, run with REHAUL_FINAL_RUN=1), the others from the number files the preamble of
% paper_rehaul.tex already loads (\rheSzDays, \rheShortSh, \rheShortShX, \rheHACS, \rhePoolWorthlessShareMed, ...).
% To place it: add % AUTO-GENERATED by writeup/make_rehaul_econ_decomp.py -- do not edit; re-run the script.
% \rhed... macros: where the straddle strategy's profits come from (Section 7 decomposition).
\providecommand{\rhedNDays}{}\renewcommand{\rhedNDays}{614}  % post-warmup option days (= rheSzDays)
\providecommand{\rhedNFull}{}\renewcommand{\rhedNFull}{6}  % pooled intraday forecasts on the full predictor set
\providecommand{\rhedNTop}{}\renewcommand{\rhedNTop}{10}  % best / worst days per forecast
\providecommand{\rhedUncRVAboveIV}{}\renewcommand{\rhedUncRVAboveIV}{28.2}  % % of the 614 days with RV_t > IV_t^2
\providecommand{\rhedUncRVIVMean}{}\renewcommand{\rhedUncRVIVMean}{0.87}  % mean RV_t / IV_t^2, 614 days
\providecommand{\rhedUncRVIVMed}{}\renewcommand{\rhedUncRVIVMed}{0.77}  % median RV_t / IV_t^2, 614 days
\providecommand{\rhedUncShortMean}{}\renewcommand{\rhedUncShortMean}{2.4}  % mean mid return of the unconditional short straddle, % of premium
\providecommand{\rhedUncShortHit}{}\renewcommand{\rhedUncShortHit}{71.8}  % % of days on which q = -1 matches sign(RV - IV^2) (always short)
\providecommand{\rhedLongNMin}{}\renewcommand{\rhedLongNMin}{117}  % lowest number of long days across the six full-set pooled forecasts, 614 days
\providecommand{\rhedLongNMax}{}\renewcommand{\rhedLongNMax}{146}  % highest number of long days across the six full-set pooled forecasts, 614 days
\providecommand{\rhedLongPctMin}{}\renewcommand{\rhedLongPctMin}{19}  % lowest % long days across the six full-set pooled forecasts, 614 days
\providecommand{\rhedLongPctMax}{}\renewcommand{\rhedLongPctMax}{24}  % highest % long days across the six full-set pooled forecasts, 614 days
\providecommand{\rhedHarLongNMin}{}\renewcommand{\rhedHarLongNMin}{150}  % lowest number of long days across the six HAR-benchmark pooled forecasts, 614 days
\providecommand{\rhedHarLongNMax}{}\renewcommand{\rhedHarLongNMax}{193}  % highest number of long days across the six HAR-benchmark pooled forecasts, 614 days
\providecommand{\rhedLongMidMin}{}\renewcommand{\rhedLongMidMin}{19}  % lowest mean mid return on long days, % across the six full-set pooled forecasts, 614 days
\providecommand{\rhedLongMidMax}{}\renewcommand{\rhedLongMidMax}{30}  % highest mean mid return on long days, % across the six full-set pooled forecasts, 614 days
\providecommand{\rhedLongBaMin}{}\renewcommand{\rhedLongBaMin}{16}  % lowest mean bid/ask return on long days, % across the six full-set pooled forecasts, 614 days
\providecommand{\rhedLongBaMax}{}\renewcommand{\rhedLongBaMax}{26}  % highest mean bid/ask return on long days, % across the six full-set pooled forecasts, 614 days
\providecommand{\rhedShortMidMin}{}\renewcommand{\rhedShortMidMin}{9}  % lowest mean mid return on short days, % across the six full-set pooled forecasts, 614 days
\providecommand{\rhedShortMidMax}{}\renewcommand{\rhedShortMidMax}{11}  % highest mean mid return on short days, % across the six full-set pooled forecasts, 614 days
\providecommand{\rhedShortBaMin}{}\renewcommand{\rhedShortBaMin}{7}  % lowest mean bid/ask return on short days, % across the six full-set pooled forecasts, 614 days
\providecommand{\rhedShortBaMax}{}\renewcommand{\rhedShortBaMax}{8}  % highest mean bid/ask return on short days, % across the six full-set pooled forecasts, 614 days
\providecommand{\rhedHarLongMidMin}{}\renewcommand{\rhedHarLongMidMin}{0}  % lowest mean mid return on long days, % across the six HAR-benchmark pooled forecasts, 614 days
\providecommand{\rhedHarLongMidMax}{}\renewcommand{\rhedHarLongMidMax}{10}  % highest mean mid return on long days, % across the six HAR-benchmark pooled forecasts, 614 days
\providecommand{\rhedHarShortMidMin}{}\renewcommand{\rhedHarShortMidMin}{3}  % lowest mean mid return on short days, % across the six HAR-benchmark pooled forecasts, 614 days
\providecommand{\rhedHarShortMidMax}{}\renewcommand{\rhedHarShortMidMax}{7}  % highest mean mid return on short days, % across the six HAR-benchmark pooled forecasts, 614 days
\providecommand{\rhedLongHitMin}{}\renewcommand{\rhedLongHitMin}{47}  % lowest % of long days with RV > IV^2 across the six full-set pooled forecasts, 614 days
\providecommand{\rhedLongHitMax}{}\renewcommand{\rhedLongHitMax}{53}  % highest % of long days with RV > IV^2 across the six full-set pooled forecasts, 614 days
\providecommand{\rhedHarLongHitMin}{}\renewcommand{\rhedHarLongHitMin}{46}  % lowest % of long days with RV > IV^2 across the six HAR-benchmark pooled forecasts, 614 days
\providecommand{\rhedHarLongHitMax}{}\renewcommand{\rhedHarLongHitMax}{48}  % highest % of long days with RV > IV^2 across the six HAR-benchmark pooled forecasts, 614 days
\providecommand{\rhedLongPosMin}{}\renewcommand{\rhedLongPosMin}{42}  % lowest % of long days with a positive return across the six full-set pooled forecasts, 614 days
\providecommand{\rhedLongPosMax}{}\renewcommand{\rhedLongPosMax}{45}  % highest % of long days with a positive return across the six full-set pooled forecasts, 614 days
\providecommand{\rhedHitMin}{}\renewcommand{\rhedHitMin}{70.7}  % lowest % of days on which q matches sign(RV - IV^2) across the six full-set pooled forecasts, 614 days
\providecommand{\rhedHitMax}{}\renewcommand{\rhedHitMax}{73.0}  % highest % of days on which q matches sign(RV - IV^2) across the six full-set pooled forecasts, 614 days
\providecommand{\rhedPosMin}{}\renewcommand{\rhedPosMin}{58.6}  % lowest % of days with a positive mid return across the six full-set pooled forecasts, 614 days
\providecommand{\rhedPosMax}{}\renewcommand{\rhedPosMax}{60.4}  % highest % of days with a positive mid return across the six full-set pooled forecasts, 614 days
\providecommand{\rhedLongShareMin}{}\renewcommand{\rhedLongShareMin}{39}  % lowest % of the summed mid return earned on long days across the six full-set pooled forecasts, 614 days
\providecommand{\rhedLongShareMax}{}\renewcommand{\rhedLongShareMax}{41}  % highest % of the summed mid return earned on long days across the six full-set pooled forecasts, 614 days
\providecommand{\rhedLongShareXMin}{}\renewcommand{\rhedLongShareXMin}{42}  % lowest % of the summed bid/ask return earned on long days across the six full-set pooled forecasts, 614 days
\providecommand{\rhedLongShareXMax}{}\renewcommand{\rhedLongShareXMax}{45}  % highest % of the summed bid/ask return earned on long days across the six full-set pooled forecasts, 614 days
\providecommand{\rhedMeanMin}{}\renewcommand{\rhedMeanMin}{11.5}  % lowest mean mid return, % across the six full-set pooled forecasts, 614 days
\providecommand{\rhedMeanMax}{}\renewcommand{\rhedMeanMax}{14.1}  % highest mean mid return, % across the six full-set pooled forecasts, 614 days
\providecommand{\rhedGainMin}{}\renewcommand{\rhedGainMin}{9.1}  % lowest mean gain over the short straddle (2 R 1(long)), % across the six full-set pooled forecasts, 614 days
\providecommand{\rhedGainMax}{}\renewcommand{\rhedGainMax}{11.6}  % highest mean gain over the short straddle (2 R 1(long)), % across the six full-set pooled forecasts, 614 days
\providecommand{\rhedGainShareMin}{}\renewcommand{\rhedGainShareMin}{79}  % lowest gain over short / mean return, % across the six full-set pooled forecasts, 614 days
\providecommand{\rhedGainShareMax}{}\renewcommand{\rhedGainShareMax}{83}  % highest gain over short / mean return, % across the six full-set pooled forecasts, 614 days
\providecommand{\rhedHarGainMin}{}\renewcommand{\rhedHarGainMin}{0.0}  % lowest mean gain over the short straddle, % across the six HAR-benchmark pooled forecasts, 614 days
\providecommand{\rhedHarGainMax}{}\renewcommand{\rhedHarGainMax}{5.0}  % highest mean gain over the short straddle, % across the six HAR-benchmark pooled forecasts, 614 days
\providecommand{\rhedOlsLongN}{}\renewcommand{\rhedOlsLongN}{198}  % long days of the closing-interval HAR forecast (sub_ols_baseline)
\providecommand{\rhedOlsLongMid}{}\renewcommand{\rhedOlsLongMid}{9}  % ... its mean mid return on long days, %
\providecommand{\rhedOlsShortMid}{}\renewcommand{\rhedOlsShortMid}{8}  % ... its mean mid return on short days, %
\providecommand{\rhedOlsLongHit}{}\renewcommand{\rhedOlsLongHit}{46}  % ... % of its long days with RV > IV^2
\providecommand{\rhedOlsSh}{}\renewcommand{\rhedOlsSh}{1.11}  % ... its Sharpe ratio, mid (= rhrOlsSharpe)
\providecommand{\rhedCalLevelMin}{}\renewcommand{\rhedCalLevelMin}{0.95}  % lowest mean forecast / mean RV across the six full-set pooled forecasts, 614 days
\providecommand{\rhedCalLevelMax}{}\renewcommand{\rhedCalLevelMax}{1.01}  % highest mean forecast / mean RV across the six full-set pooled forecasts, 614 days
\providecommand{\rhedCalScaleMin}{}\renewcommand{\rhedCalScaleMin}{1.09}  % lowest mean of RV / forecast (QLIKE-optimal scale) across the six full-set pooled forecasts, 614 days
\providecommand{\rhedCalScaleMax}{}\renewcommand{\rhedCalScaleMax}{1.20}  % highest mean of RV / forecast (QLIKE-optimal scale) across the six full-set pooled forecasts, 614 days
\providecommand{\rhedHarCalScaleMin}{}\renewcommand{\rhedHarCalScaleMin}{1.04}  % lowest mean of RV / forecast across the six HAR-benchmark pooled forecasts, 614 days
\providecommand{\rhedHarCalScaleMax}{}\renewcommand{\rhedHarCalScaleMax}{1.10}  % highest mean of RV / forecast across the six HAR-benchmark pooled forecasts, 614 days
\providecommand{\rhedMedFcIVMin}{}\renewcommand{\rhedMedFcIVMin}{0.75}  % lowest median forecast / IV^2 across the six full-set pooled forecasts, 614 days
\providecommand{\rhedMedFcIVMax}{}\renewcommand{\rhedMedFcIVMax}{0.82}  % highest median forecast / IV^2 across the six full-set pooled forecasts, 614 days
\providecommand{\rhedBcScaleLoMin}{}\renewcommand{\rhedBcScaleLoMin}{1.06}  % lowest lowest causal scale c_t over the 614 days across the six full-set pooled forecasts, 614 days
\providecommand{\rhedBcScaleLoMax}{}\renewcommand{\rhedBcScaleLoMax}{1.15}  % highest lowest causal scale c_t over the 614 days across the six full-set pooled forecasts, 614 days
\providecommand{\rhedBcScaleHiMin}{}\renewcommand{\rhedBcScaleHiMin}{1.10}  % lowest highest causal scale c_t over the 614 days across the six full-set pooled forecasts, 614 days
\providecommand{\rhedBcScaleHiMax}{}\renewcommand{\rhedBcScaleHiMax}{1.22}  % highest highest causal scale c_t over the 614 days across the six full-set pooled forecasts, 614 days
\providecommand{\rhedBcFlipMin}{}\renewcommand{\rhedBcFlipMin}{55}  % lowest positions changed by the bias correction across the six full-set pooled forecasts, 614 days
\providecommand{\rhedBcFlipMax}{}\renewcommand{\rhedBcFlipMax}{120}  % highest positions changed by the bias correction across the six full-set pooled forecasts, 614 days
\providecommand{\rhedBcLongMin}{}\renewcommand{\rhedBcLongMin}{201}  % lowest long days under the bias correction across the six full-set pooled forecasts, 614 days
\providecommand{\rhedBcLongMax}{}\renewcommand{\rhedBcLongMax}{241}  % highest long days under the bias correction across the six full-set pooled forecasts, 614 days
\providecommand{\rhedBcFlipRMin}{}\renewcommand{\rhedBcFlipRMin}{\ensuremath{-}24}  % lowest mean mid straddle return R on the changed days, % across the six full-set pooled forecasts, 614 days
\providecommand{\rhedBcFlipRMax}{}\renewcommand{\rhedBcFlipRMax}{\ensuremath{-}8}  % highest mean mid straddle return R on the changed days, % across the six full-set pooled forecasts, 614 days
\providecommand{\rhedBcShMin}{}\renewcommand{\rhedBcShMin}{0.92}  % lowest Sharpe ratio mid under the bias correction (diagnostic) across the six full-set pooled forecasts, 614 days
\providecommand{\rhedBcShMax}{}\renewcommand{\rhedBcShMax}{1.32}  % highest Sharpe ratio mid under the bias correction (diagnostic) across the six full-set pooled forecasts, 614 days
\providecommand{\rhedBcShXMin}{}\renewcommand{\rhedBcShXMin}{0.57}  % lowest Sharpe ratio bid/ask under the bias correction across the six full-set pooled forecasts, 614 days
\providecommand{\rhedBcShXMax}{}\renewcommand{\rhedBcShXMax}{0.97}  % highest Sharpe ratio bid/ask under the bias correction across the six full-set pooled forecasts, 614 days
\providecommand{\rhedLevFlipMin}{}\renewcommand{\rhedLevFlipMin}{4}  % lowest positions changed by the level correction (sum RV / sum fc) across the six full-set pooled forecasts, 614 days
\providecommand{\rhedLevFlipMax}{}\renewcommand{\rhedLevFlipMax}{19}  % highest positions changed by the level correction (sum RV / sum fc) across the six full-set pooled forecasts, 614 days
\providecommand{\rhedLevShMin}{}\renewcommand{\rhedLevShMin}{1.61}  % lowest Sharpe ratio mid under the level correction across the six full-set pooled forecasts, 614 days
\providecommand{\rhedLevShMax}{}\renewcommand{\rhedLevShMax}{2.01}  % highest Sharpe ratio mid under the level correction across the six full-set pooled forecasts, 614 days
\providecommand{\rhedHarBcShMin}{}\renewcommand{\rhedHarBcShMin}{0.76}  % lowest Sharpe ratio mid under the bias correction across the six HAR-benchmark pooled forecasts, 614 days
\providecommand{\rhedHarBcShMax}{}\renewcommand{\rhedHarBcShMax}{1.05}  % highest Sharpe ratio mid under the bias correction across the six HAR-benchmark pooled forecasts, 614 days
\providecommand{\rhedBcShUpN}{}\renewcommand{\rhedBcShUpN}{0}  % full-set forecasts whose mid SR rises under the bias correction
\providecommand{\rhedBcShDownN}{}\renewcommand{\rhedBcShDownN}{6}  % full-set forecasts whose mid SR falls under the bias correction
\providecommand{\rhedHarBcShUpN}{}\renewcommand{\rhedHarBcShUpN}{2}  % HAR-benchmark forecasts whose mid SR rises under the bias correction
\providecommand{\rhedMoveSqIV}{}\renewcommand{\rhedMoveSqIV}{1.01}  % mean of the squared S&P 500 half-hour log return over IV_t^2, 614 days
\providecommand{\rhedMoveSqRV}{}\renewcommand{\rhedMoveSqRV}{1.15}  % mean squared S&P 500 half-hour log return over mean RV_t, 614 days
\providecommand{\rhedOlsCalScale}{}\renewcommand{\rhedOlsCalScale}{1.01}  % mean of RV / forecast, closing-interval HAR forecast
\providecommand{\rhedOlsBcSh}{}\renewcommand{\rhedOlsBcSh}{1.12}  % Sharpe ratio mid under the bias correction, closing-interval HAR
\providecommand{\rhedBcFlipRidge}{}\renewcommand{\rhedBcFlipRidge}{105}  % positions changed by the bias correction, poolpt13_ridge_all_features
\providecommand{\rhedBcShRidge}{}\renewcommand{\rhedBcShRidge}{1.17}  % Sharpe ratio mid under the bias correction, poolpt13_ridge_all_features
\providecommand{\rhedBcFlipLasso}{}\renewcommand{\rhedBcFlipLasso}{120}  % positions changed by the bias correction, poolpt13_lasso_all_features
\providecommand{\rhedBcShLasso}{}\renewcommand{\rhedBcShLasso}{1.31}  % Sharpe ratio mid under the bias correction, poolpt13_lasso_all_features
\providecommand{\rhedBcFlipEnet}{}\renewcommand{\rhedBcFlipEnet}{108}  % positions changed by the bias correction, poolpt13_enet_all_features
\providecommand{\rhedBcShEnet}{}\renewcommand{\rhedBcShEnet}{0.99}  % Sharpe ratio mid under the bias correction, poolpt13_enet_all_features
\providecommand{\rhedBcFlipLgbm}{}\renewcommand{\rhedBcFlipLgbm}{55}  % positions changed by the bias correction, pooltree_lgbm_all_features
\providecommand{\rhedBcShLgbm}{}\renewcommand{\rhedBcShLgbm}{1.32}  % Sharpe ratio mid under the bias correction, pooltree_lgbm_all_features
\providecommand{\rhedBcFlipXgb}{}\renewcommand{\rhedBcFlipXgb}{61}  % positions changed by the bias correction, pooltree_xgb_all_features
\providecommand{\rhedBcShXgb}{}\renewcommand{\rhedBcShXgb}{1.15}  % Sharpe ratio mid under the bias correction, pooltree_xgb_all_features
\providecommand{\rhedBcFlipRf}{}\renewcommand{\rhedBcFlipRf}{90}  % positions changed by the bias correction, pooltree_rf_all_features
\providecommand{\rhedBcShRf}{}\renewcommand{\rhedBcShRf}{0.92}  % Sharpe ratio mid under the bias correction, pooltree_rf_all_features
\providecommand{\rhedQMonoN}{}\renewcommand{\rhedQMonoN}{6}  % full-set forecasts whose mean RV/IV^2 rises across all five quintiles
\providecommand{\rhedHarQMonoN}{}\renewcommand{\rhedHarQMonoN}{4}  % HAR-benchmark forecasts whose mean RV/IV^2 rises across all five quintiles
\providecommand{\rhedQRPatternN}{}\renewcommand{\rhedQRPatternN}{6}  % full-set forecasts with mean R < 0 in quintiles 1-4 and > 0 in 5
\providecommand{\rhedHarQRPatternN}{}\renewcommand{\rhedHarQRPatternN}{0}  % HAR-benchmark forecasts with that pattern
\providecommand{\rhedQOneRVIV}{}\renewcommand{\rhedQOneRVIV}{0.73}  % mean RV/IV^2 in the lowest quintile, mean of the six full-set
\providecommand{\rhedQFiveRVIV}{}\renewcommand{\rhedQFiveRVIV}{1.05}  % mean RV/IV^2 in the highest quintile, mean of the six full-set
\providecommand{\rhedQFiveR}{}\renewcommand{\rhedQFiveR}{26}  % mean straddle return R in the highest quintile, %, mean of the six full-set
\providecommand{\rhedQFiveFc}{}\renewcommand{\rhedQFiveFc}{1.18}  % geometric-mean forecast / IV^2 in the highest quintile, mean of the six full-set
\providecommand{\rhedQFourFc}{}\renewcommand{\rhedQFourFc}{0.93}  % ... in the fourth quintile
\providecommand{\rhedQFiveLong}{}\renewcommand{\rhedQFiveLong}{99}  % % long days in the highest quintile, mean of the six full-set
\providecommand{\rhedQFiveHit}{}\renewcommand{\rhedQFiveHit}{52}  % hit rate in the highest quintile, %, mean of the six full-set
\providecommand{\rhedQOneHit}{}\renewcommand{\rhedQOneHit}{88}  % hit rate in the lowest quintile, %, mean of the six full-set
\providecommand{\rhedQFiveShare}{}\renewcommand{\rhedQFiveShare}{43}  % % of the summed mid return earned in the highest quintile, six full-set forecasts
\providecommand{\rhedQLowRMin}{}\renewcommand{\rhedQLowRMin}{\ensuremath{-}13}  % most negative mean straddle return R in quintiles 1-4, %, mean of the six
\providecommand{\rhedQLowRMax}{}\renewcommand{\rhedQLowRMax}{\ensuremath{-}6}  % least negative mean straddle return R in quintiles 1-4, %, mean of the six
\providecommand{\rhedHarQOneRVIV}{}\renewcommand{\rhedHarQOneRVIV}{0.74}  % mean RV/IV^2 in the lowest quintile, mean of the six Har
\providecommand{\rhedHarQFiveRVIV}{}\renewcommand{\rhedHarQFiveRVIV}{1.02}  % mean RV/IV^2 in the highest quintile, mean of the six Har
\providecommand{\rhedHarQFiveR}{}\renewcommand{\rhedHarQFiveR}{7}  % mean straddle return R in the highest quintile, %, mean of the six Har
\providecommand{\rhedHarQFiveFc}{}\renewcommand{\rhedHarQFiveFc}{1.23}  % geometric-mean forecast / IV^2 in the highest quintile, mean of the six Har
\providecommand{\rhedHarQFourFc}{}\renewcommand{\rhedHarQFourFc}{0.98}  % ... in the fourth quintile
\providecommand{\rhedHarQFiveLong}{}\renewcommand{\rhedHarQFiveLong}{100}  % % long days in the highest quintile, mean of the six Har
\providecommand{\rhedHarQFiveHit}{}\renewcommand{\rhedHarQFiveHit}{47}  % hit rate in the highest quintile, %, mean of the six Har
\providecommand{\rhedHarQOneHit}{}\renewcommand{\rhedHarQOneHit}{89}  % hit rate in the lowest quintile, %, mean of the six Har
\providecommand{\rhedHarQFiveShare}{}\renewcommand{\rhedHarQFiveShare}{23}  % % of the summed mid return earned in the highest quintile, six Har forecasts
\providecommand{\rhedHarQLowRMin}{}\renewcommand{\rhedHarQLowRMin}{\ensuremath{-}13}  % most negative mean straddle return R in quintiles 1-4, %, mean of the six
\providecommand{\rhedHarQLowRMax}{}\renewcommand{\rhedHarQLowRMax}{5}  % least negative mean straddle return R in quintiles 1-4, %, mean of the six
\providecommand{\rhedQFiveShareRngMin}{}\renewcommand{\rhedQFiveShareRngMin}{40}  % lowest % of the summed mid return in quintile 5 across the six full-set pooled forecasts, 614 days
\providecommand{\rhedQFiveShareRngMax}{}\renewcommand{\rhedQFiveShareRngMax}{48}  % highest % of the summed mid return in quintile 5 across the six full-set pooled forecasts, 614 days
\providecommand{\rhedBestUnion}{}\renewcommand{\rhedBestUnion}{11}  % days in the union of the six full-set forecasts' ten best days
\providecommand{\rhedWorstUnion}{}\renewcommand{\rhedWorstUnion}{11}  % days in the union of the six full-set forecasts' ten worst days
\providecommand{\rhedBestAllSix}{}\renewcommand{\rhedBestAllSix}{9}  % days among the ten best of all six forecasts
\providecommand{\rhedWorstAllSix}{}\renewcommand{\rhedWorstAllSix}{9}  % days among the ten worst of all six forecasts
\providecommand{\rhedTopDate}{}\renewcommand{\rhedTopDate}{22~November~2021}  % the day of the largest straddle return among the best days
\providecommand{\rhedTopR}{}\renewcommand{\rhedTopR}{1{,}032}  % ... the straddle's mid return that day, % of premium
\providecommand{\rhedTopRVIV}{}\renewcommand{\rhedTopRVIV}{2.79}  % ... RV / IV^2 that day
\providecommand{\rhedTopMove}{}\renewcommand{\rhedTopMove}{\ensuremath{-}0.83}  % ... S&P 500 15:30-16:00 log return that day, %
\providecommand{\rhedTopShareMin}{}\renewcommand{\rhedTopShareMin}{12}  % lowest % of a full-set forecast's summed mid return earned on that one day
\providecommand{\rhedTopShareMax}{}\renewcommand{\rhedTopShareMax}{15}  % highest % of a full-set forecast's summed mid return earned on that one day
\providecommand{\rhedBestFomcN}{}\renewcommand{\rhedBestFomcN}{1}  % best-day union: FOMC statement days
\providecommand{\rhedWorstFomcN}{}\renewcommand{\rhedWorstFomcN}{1}  % worst-day union: FOMC statement days
\providecommand{\rhedWorstMeN}{}\renewcommand{\rhedWorstMeN}{3}  % worst-day union: month ends
\providecommand{\rhedBestMeN}{}\renewcommand{\rhedBestMeN}{0}  % best-day union: month ends
\providecommand{\rhedBestOpexN}{}\renewcommand{\rhedBestOpexN}{0}  % best-day union: monthly expirations
\providecommand{\rhedWorstOpexN}{}\renewcommand{\rhedWorstOpexN}{0}  % worst-day union: monthly expirations
\providecommand{\rhedBestNoTagN}{}\renewcommand{\rhedBestNoTagN}{8}  % best-day union: days without any event tag (calendar as covered)
\providecommand{\rhedBestAfterRelN}{}\renewcommand{\rhedBestAfterRelN}{5}  % best-day union: days after the release calendar ends (dagger)
\providecommand{\rhedWorstAfterRelN}{}\renewcommand{\rhedWorstAfterRelN}{4}  % worst-day union: days after the release calendar ends
\providecommand{\rhedDShortAllMin}{}\renewcommand{\rhedDShortAllMin}{1.21}  % lowest Delta SR mid vs the short straddle, all days across the six full-set pooled forecasts, 614 days
\providecommand{\rhedDShortAllMax}{}\renewcommand{\rhedDShortAllMax}{1.56}  % highest Delta SR mid vs the short straddle, all days across the six full-set pooled forecasts, 614 days
\providecommand{\rhedDShortExBestMin}{}\renewcommand{\rhedDShortExBestMin}{\ensuremath{-}0.54}  % lowest Delta SR mid vs the short straddle without the forecast's ten best days across the six full-set pooled forecasts, 614 days
\providecommand{\rhedDShortExBestMax}{}\renewcommand{\rhedDShortExBestMax}{\ensuremath{-}0.12}  % highest Delta SR mid vs the short straddle without the forecast's ten best days across the six full-set pooled forecasts, 614 days
\providecommand{\rhedDShortExBothMin}{}\renewcommand{\rhedDShortExBothMin}{\ensuremath{-}0.64}  % lowest ... without its ten best and ten worst days across the six full-set pooled forecasts, 614 days
\providecommand{\rhedDShortExBothMax}{}\renewcommand{\rhedDShortExBothMax}{\ensuremath{-}0.14}  % highest ... without its ten best and ten worst days across the six full-set pooled forecasts, 614 days
\providecommand{\rhedBestShareMin}{}\renewcommand{\rhedBestShareMin}{107}  % lowest % of the summed gain over the short straddle earned on the ten best days across the six full-set pooled forecasts, 614 days
\providecommand{\rhedBestShareMax}{}\renewcommand{\rhedBestShareMax}{138}  % highest % of the summed gain over the short straddle earned on the ten best days across the six full-set pooled forecasts, 614 days
\providecommand{\rhedDHarAllMin}{}\renewcommand{\rhedDHarAllMin}{0.56}  % lowest Delta SR mid vs the HAR twin, all days across the six full-set pooled forecasts, 614 days
\providecommand{\rhedDHarAllMax}{}\renewcommand{\rhedDHarAllMax}{1.56}  % highest Delta SR mid vs the HAR twin, all days across the six full-set pooled forecasts, 614 days
\providecommand{\rhedDHarExBestMin}{}\renewcommand{\rhedDHarExBestMin}{0.30}  % lowest Delta SR mid vs the HAR twin without the forecast's ten best days across the six full-set pooled forecasts, 614 days
\providecommand{\rhedDHarExBestMax}{}\renewcommand{\rhedDHarExBestMax}{0.83}  % highest Delta SR mid vs the HAR twin without the forecast's ten best days across the six full-set pooled forecasts, 614 days
\providecommand{\rhedDHarExBothMin}{}\renewcommand{\rhedDHarExBothMin}{0.36}  % lowest ... without its ten best and ten worst days across the six full-set pooled forecasts, 614 days
\providecommand{\rhedDHarExBothMax}{}\renewcommand{\rhedDHarExBothMax}{0.97}  % highest ... without its ten best and ten worst days across the six full-set pooled forecasts, 614 days
\providecommand{\rhedHarBestShareMin}{}\renewcommand{\rhedHarBestShareMin}{47}  % lowest % of the summed gain over the HAR twin earned on the ten best days across the six full-set pooled forecasts, 614 days
\providecommand{\rhedHarBestShareMax}{}\renewcommand{\rhedHarBestShareMax}{61}  % highest % of the summed gain over the HAR twin earned on the ten best days across the six full-set pooled forecasts, 614 days
\providecommand{\rhedDShortExBestNeg}{}\renewcommand{\rhedDShortExBestNeg}{6}  % full-set forecasts whose Delta SR vs short is negative without the ten best days
\providecommand{\rhedDHarExBestNeg}{}\renewcommand{\rhedDHarExBestNeg}{0}  % full-set forecasts whose Delta SR vs the HAR twin is negative without the ten best days
\providecommand{\rhedFomcN}{}\renewcommand{\rhedFomcN}{21}  % FOMC days among the 614 post-warmup days
\providecommand{\rhedFomcR}{}\renewcommand{\rhedFomcR}{25}  % mean mid straddle return R on FOMC days, %
\providecommand{\rhedFomcRT}{}\renewcommand{\rhedFomcRT}{1.02}  % ... its t (i.i.d.; the days are weeks apart)
\providecommand{\rhedFomcLong}{}\renewcommand{\rhedFomcLong}{25}  % % of FOMC days long, mean of the six full-set forecasts
\providecommand{\rhedFomcStrat}{}\renewcommand{\rhedFomcStrat}{\ensuremath{-}6}  % mean mid strategy return on FOMC days, %, mean of the six full-set
\providecommand{\rhedFomcOtherR}{}\renewcommand{\rhedFomcOtherR}{\ensuremath{-}3}  % mean straddle R on the other (covered) days, %
\providecommand{\rhedFomcOtherStrat}{}\renewcommand{\rhedFomcOtherStrat}{13}  % mean strategy return on the other (covered) days, %
\providecommand{\rhedFomcRVIV}{}\renewcommand{\rhedFomcRVIV}{0.76}  % mean RV_t / IV_t^2 on FOMC days
\providecommand{\rhedFomcMoveSqIV}{}\renewcommand{\rhedFomcMoveSqIV}{1.43}  % mean squared S&P 500 half-hour log return / IV_t^2 on FOMC days
\providecommand{\rhedFomcShare}{}\renewcommand{\rhedFomcShare}{\ensuremath{-}2}  % % of the summed mid return earned on FOMC days, mean of the six full-set forecasts
\providecommand{\rhedCpiN}{}\renewcommand{\rhedCpiN}{25}  % CPI days among the 614 post-warmup days
\providecommand{\rhedCpiR}{}\renewcommand{\rhedCpiR}{\ensuremath{-}12}  % mean mid straddle return R on CPI days, %
\providecommand{\rhedCpiRT}{}\renewcommand{\rhedCpiRT}{\ensuremath{-}0.55}  % ... its t (i.i.d.; the days are weeks apart)
\providecommand{\rhedCpiLong}{}\renewcommand{\rhedCpiLong}{16}  % % of CPI days long, mean of the six full-set forecasts
\providecommand{\rhedCpiStrat}{}\renewcommand{\rhedCpiStrat}{22}  % mean mid strategy return on CPI days, %, mean of the six full-set
\providecommand{\rhedCpiOtherR}{}\renewcommand{\rhedCpiOtherR}{\ensuremath{-}6}  % mean straddle R on the other (covered) days, %
\providecommand{\rhedCpiOtherStrat}{}\renewcommand{\rhedCpiOtherStrat}{11}  % mean strategy return on the other (covered) days, %
\providecommand{\rhedCpiRVIV}{}\renewcommand{\rhedCpiRVIV}{0.80}  % mean RV_t / IV_t^2 on CPI days
\providecommand{\rhedCpiMoveSqIV}{}\renewcommand{\rhedCpiMoveSqIV}{0.86}  % mean squared S&P 500 half-hour log return / IV_t^2 on CPI days
\providecommand{\rhedCpiShare}{}\renewcommand{\rhedCpiShare}{7}  % % of the summed mid return earned on CPI days, mean of the six full-set forecasts
\providecommand{\rhedNfpN}{}\renewcommand{\rhedNfpN}{27}  % NFP days among the 614 post-warmup days
\providecommand{\rhedNfpR}{}\renewcommand{\rhedNfpR}{1}  % mean mid straddle return R on NFP days, %
\providecommand{\rhedNfpRT}{}\renewcommand{\rhedNfpRT}{0.06}  % ... its t (i.i.d.; the days are weeks apart)
\providecommand{\rhedNfpLong}{}\renewcommand{\rhedNfpLong}{12}  % % of NFP days long, mean of the six full-set forecasts
\providecommand{\rhedNfpStrat}{}\renewcommand{\rhedNfpStrat}{20}  % mean mid strategy return on NFP days, %, mean of the six full-set
\providecommand{\rhedNfpOtherR}{}\renewcommand{\rhedNfpOtherR}{\ensuremath{-}6}  % mean straddle R on the other (covered) days, %
\providecommand{\rhedNfpOtherStrat}{}\renewcommand{\rhedNfpOtherStrat}{11}  % mean strategy return on the other (covered) days, %
\providecommand{\rhedNfpRVIV}{}\renewcommand{\rhedNfpRVIV}{0.73}  % mean RV_t / IV_t^2 on NFP days
\providecommand{\rhedNfpMoveSqIV}{}\renewcommand{\rhedNfpMoveSqIV}{0.93}  % mean squared S&P 500 half-hour log return / IV_t^2 on NFP days
\providecommand{\rhedNfpShare}{}\renewcommand{\rhedNfpShare}{7}  % % of the summed mid return earned on NFP days, mean of the six full-set forecasts
\providecommand{\rhedOpexN}{}\renewcommand{\rhedOpexN}{33}  % OpEx days among the 614 post-warmup days
\providecommand{\rhedOpexR}{}\renewcommand{\rhedOpexR}{\ensuremath{-}45}  % mean mid straddle return R on OpEx days, %
\providecommand{\rhedOpexRT}{}\renewcommand{\rhedOpexRT}{\ensuremath{-}4.01}  % ... its t (i.i.d.; the days are weeks apart)
\providecommand{\rhedOpexLong}{}\renewcommand{\rhedOpexLong}{9}  % % of OpEx days long, mean of the six full-set forecasts
\providecommand{\rhedOpexStrat}{}\renewcommand{\rhedOpexStrat}{36}  % mean mid strategy return on OpEx days, %, mean of the six full-set
\providecommand{\rhedOpexOtherR}{}\renewcommand{\rhedOpexOtherR}{0}  % mean straddle R on the other (covered) days, %
\providecommand{\rhedOpexOtherStrat}{}\renewcommand{\rhedOpexOtherStrat}{11}  % mean strategy return on the other (covered) days, %
\providecommand{\rhedOpexRVIV}{}\renewcommand{\rhedOpexRVIV}{0.83}  % mean RV_t / IV_t^2 on OpEx days
\providecommand{\rhedOpexMoveSqIV}{}\renewcommand{\rhedOpexMoveSqIV}{0.40}  % mean squared S&P 500 half-hour log return / IV_t^2 on OpEx days
\providecommand{\rhedOpexShare}{}\renewcommand{\rhedOpexShare}{15}  % % of the summed mid return earned on OpEx days, mean of the six full-set forecasts
\providecommand{\rhedTwN}{}\renewcommand{\rhedTwN}{11}  % TW days among the 614 post-warmup days
\providecommand{\rhedTwR}{}\renewcommand{\rhedTwR}{\ensuremath{-}52}  % mean mid straddle return R on TW days, %
\providecommand{\rhedTwRT}{}\renewcommand{\rhedTwRT}{\ensuremath{-}3.98}  % ... its t (i.i.d.; the days are weeks apart)
\providecommand{\rhedTwLong}{}\renewcommand{\rhedTwLong}{8}  % % of TW days long, mean of the six full-set forecasts
\providecommand{\rhedTwStrat}{}\renewcommand{\rhedTwStrat}{37}  % mean mid strategy return on TW days, %, mean of the six full-set
\providecommand{\rhedTwOtherR}{}\renewcommand{\rhedTwOtherR}{\ensuremath{-}2}  % mean straddle R on the other (covered) days, %
\providecommand{\rhedTwOtherStrat}{}\renewcommand{\rhedTwOtherStrat}{12}  % mean strategy return on the other (covered) days, %
\providecommand{\rhedTwRVIV}{}\renewcommand{\rhedTwRVIV}{0.89}  % mean RV_t / IV_t^2 on TW days
\providecommand{\rhedTwMoveSqIV}{}\renewcommand{\rhedTwMoveSqIV}{0.31}  % mean squared S&P 500 half-hour log return / IV_t^2 on TW days
\providecommand{\rhedTwShare}{}\renewcommand{\rhedTwShare}{5}  % % of the summed mid return earned on TW days, mean of the six full-set forecasts
\providecommand{\rhedMeN}{}\renewcommand{\rhedMeN}{33}  % ME days among the 614 post-warmup days
\providecommand{\rhedMeR}{}\renewcommand{\rhedMeR}{70}  % mean mid straddle return R on ME days, %
\providecommand{\rhedMeRT}{}\renewcommand{\rhedMeRT}{2.79}  % ... its t (i.i.d.; the days are weeks apart)
\providecommand{\rhedMeLong}{}\renewcommand{\rhedMeLong}{6}  % % of ME days long, mean of the six full-set forecasts
\providecommand{\rhedMeStrat}{}\renewcommand{\rhedMeStrat}{\ensuremath{-}58}  % mean mid strategy return on ME days, %, mean of the six full-set
\providecommand{\rhedMeOtherR}{}\renewcommand{\rhedMeOtherR}{\ensuremath{-}7}  % mean straddle R on the other (covered) days, %
\providecommand{\rhedMeOtherStrat}{}\renewcommand{\rhedMeOtherStrat}{17}  % mean strategy return on the other (covered) days, %
\providecommand{\rhedMeRVIV}{}\renewcommand{\rhedMeRVIV}{1.04}  % mean RV_t / IV_t^2 on ME days
\providecommand{\rhedMeMoveSqIV}{}\renewcommand{\rhedMeMoveSqIV}{1.93}  % mean squared S&P 500 half-hour log return / IV_t^2 on ME days
\providecommand{\rhedMeShare}{}\renewcommand{\rhedMeShare}{\ensuremath{-}24}  % % of the summed mid return earned on ME days, mean of the six full-set forecasts
\providecommand{\rhedRelEndText}{}\renewcommand{\rhedRelEndText}{26~October~2023}  % earliest last flag of the six macro-release series in data/releases.parquet
\providecommand{\rhedRelFomcEndText}{}\renewcommand{\rhedRelFomcEndText}{1~November~2023}  % last FOMC flag in data/releases.parquet
\providecommand{\rhedFomcSuppN}{}\renewcommand{\rhedFomcSuppN}{3}  % FOMC statement days of the sample taken from live/close_signal/state/fomc_statement_dates.csv
\providecommand{\rhedRelAfterN}{}\renewcommand{\rhedRelAfterN}{126}  % post-warmup days after the release calendar ends
\providecommand{\rhedWorthN}{}\renewcommand{\rhedWorthN}{137}  % post-warmup days with R_t = -1 (= rheWorthlessDays)
\providecommand{\rhedWorthPct}{}\renewcommand{\rhedWorthPct}{22.3}  % ... % of the 614 days
\providecommand{\rhedWorthExpIV}{}\renewcommand{\rhedWorthExpIV}{128}  % expected R = -1 days under a normal half-hour log move with volatility IV_t
\providecommand{\rhedWorthExpIVPct}{}\renewcommand{\rhedWorthExpIVPct}{20.9}  % ... as % of the 614 days
\providecommand{\rhedWorthZ}{}\renewcommand{\rhedWorthZ}{0.90}  % (observed - expected) / sd of the count under that normal (independent days)
\providecommand{\rhedWorthExpEmp}{}\renewcommand{\rhedWorthExpEmp}{132}  % expected R = -1 days under the empirical distribution of the earlier days' moves standardized by IV
\providecommand{\rhedWorthExpRaw}{}\renewcommand{\rhedWorthExpRaw}{118}  % expected R = -1 days under the earlier days' raw half-hour moves
\providecommand{\rhedGapFiveN}{}\renewcommand{\rhedGapFiveN}{611}  % post-warmup days with K_C - K_P = 5 points
\providecommand{\rhedGapOtherN}{}\renewcommand{\rhedGapOtherN}{3}  % post-warmup days with another strike gap
\providecommand{\rhedMoveSdPts}{}\renewcommand{\rhedMoveSdPts}{9.6}  % median one-sd implied half-hour move, index points
\providecommand{\rhedStrikeDist}{}\renewcommand{\rhedStrikeDist}{1.23}  % mean distance of the close to the nearest 5-point strike, points
\providecommand{\rhedStrikeDistUnif}{}\renewcommand{\rhedStrikeDistUnif}{1.25}  % ... its mean under a uniform close
\providecommand{\rhedKsD}{}\renewcommand{\rhedKsD}{0.026}  % Kolmogorov-Smirnov D of the close's position within its 5-point strike interval against a uniform
\providecommand{\rhedKsCV}{}\renewcommand{\rhedKsCV}{0.055}  % ... asymptotic 5% critical value
\providecommand{\rhedZSd}{}\renewcommand{\rhedZSd}{1.00}  % sd of the half-hour log move standardized by IV_t, 614 days
\providecommand{\rhedZKurt}{}\renewcommand{\rhedZKurt}{2.0}  % excess kurtosis of that standardized move
\providecommand{\rhedWorthShortMin}{}\renewcommand{\rhedWorthShortMin}{110}  % lowest R = -1 days on which the forecast is short across the six full-set pooled forecasts, 614 days
\providecommand{\rhedWorthShortMax}{}\renewcommand{\rhedWorthShortMax}{115}  % highest R = -1 days on which the forecast is short across the six full-set pooled forecasts, 614 days
\providecommand{\rhedWorthOpexN}{}\renewcommand{\rhedWorthOpexN}{10}  % R = -1 days among the monthly-expiration days
\providecommand{\rhedWorthOpexExp}{}\renewcommand{\rhedWorthOpexExp}{6.1}  % ... expected under the normal with IV_t
\providecommand{\rhedWorthOtherPct}{}\renewcommand{\rhedWorthOtherPct}{21.9}  % % of the other days with R = -1
\providecommand{\rhedWorthOtherExpPct}{}\renewcommand{\rhedWorthOtherExpPct}{21.0}  % ... expected under the normal with IV_t, %
\providecommand{\rhedWorthOpexZ}{}\renewcommand{\rhedWorthOpexZ}{1.77}  % (observed - expected) / sd of the count on the monthly-expiration days
\providecommand{\rhedRegAlphaAMin}{}\renewcommand{\rhedRegAlphaAMin}{11.0}  % lowest alpha, (a) on the short straddle, mid, % per day across the six full-set pooled forecasts, 614 days
\providecommand{\rhedRegAlphaAMax}{}\renewcommand{\rhedRegAlphaAMax}{13.4}  % highest alpha, (a) on the short straddle, mid, % per day across the six full-set pooled forecasts, 614 days
\providecommand{\rhedRegAlphaAXMin}{}\renewcommand{\rhedRegAlphaAXMin}{8.3}  % lowest alpha, (a) on the short straddle, bid/ask, % per day across the six full-set pooled forecasts, 614 days
\providecommand{\rhedRegAlphaAXMax}{}\renewcommand{\rhedRegAlphaAXMax}{10.7}  % highest alpha, (a) on the short straddle, bid/ask, % per day across the six full-set pooled forecasts, 614 days
\providecommand{\rhedRegBetaShortMin}{}\renewcommand{\rhedRegBetaShortMin}{0.23}  % lowest beta on the short-straddle return, (a), mid across the six full-set pooled forecasts, 614 days
\providecommand{\rhedRegBetaShortMax}{}\renewcommand{\rhedRegBetaShortMax}{0.28}  % highest beta on the short-straddle return, (a), mid across the six full-set pooled forecasts, 614 days
\providecommand{\rhedRegAlphaAPMin}{}\renewcommand{\rhedRegAlphaAPMin}{0.005}  % lowest fixed-b p of alpha, (a), mid across the six full-set pooled forecasts, 614 days
\providecommand{\rhedRegAlphaAPMax}{}\renewcommand{\rhedRegAlphaAPMax}{0.025}  % highest fixed-b p of alpha, (a), mid across the six full-set pooled forecasts, 614 days
\providecommand{\rhedRegAlphaAXPMin}{}\renewcommand{\rhedRegAlphaAXPMin}{0.021}  % lowest fixed-b p of alpha, (a), bid/ask across the six full-set pooled forecasts, 614 days
\providecommand{\rhedRegAlphaAXPMax}{}\renewcommand{\rhedRegAlphaAXPMax}{0.077}  % highest fixed-b p of alpha, (a), bid/ask across the six full-set pooled forecasts, 614 days
\providecommand{\rhedRegBetaShortTMin}{}\renewcommand{\rhedRegBetaShortTMin}{1.22}  % lowest t of the beta on the short straddle, (a), mid across the six full-set pooled forecasts, 614 days
\providecommand{\rhedRegBetaShortTMax}{}\renewcommand{\rhedRegBetaShortTMax}{1.42}  % highest t of the beta on the short straddle, (a), mid across the six full-set pooled forecasts, 614 days
\providecommand{\rhedRegRsqAMin}{}\renewcommand{\rhedRegRsqAMin}{0.06}  % lowest R^2 of (a), mid across the six full-set pooled forecasts, 614 days
\providecommand{\rhedRegRsqAMax}{}\renewcommand{\rhedRegRsqAMax}{0.08}  % highest R^2 of (a), mid across the six full-set pooled forecasts, 614 days
\providecommand{\rhedRegBetaAbsMin}{}\renewcommand{\rhedRegBetaAbsMin}{\ensuremath{-}1.87}  % lowest beta on |r30| (bp), (b), mid, % of premium per bp across the six full-set pooled forecasts, 614 days
\providecommand{\rhedRegBetaAbsMax}{}\renewcommand{\rhedRegBetaAbsMax}{\ensuremath{-}1.40}  % highest beta on |r30| (bp), (b), mid, % of premium per bp across the six full-set pooled forecasts, 614 days
\providecommand{\rhedRegBetaAbsTMin}{}\renewcommand{\rhedRegBetaAbsTMin}{\ensuremath{-}3.49}  % lowest t of the beta on |r30|, (b), mid across the six full-set pooled forecasts, 614 days
\providecommand{\rhedRegBetaAbsTMax}{}\renewcommand{\rhedRegBetaAbsTMax}{\ensuremath{-}2.31}  % highest t of the beta on |r30|, (b), mid across the six full-set pooled forecasts, 614 days
\providecommand{\rhedRegBetaMoveMin}{}\renewcommand{\rhedRegBetaMoveMin}{\ensuremath{-}0.01}  % lowest beta on r30 (bp), (b), mid across the six full-set pooled forecasts, 614 days
\providecommand{\rhedRegBetaMoveMax}{}\renewcommand{\rhedRegBetaMoveMax}{0.14}  % highest beta on r30 (bp), (b), mid across the six full-set pooled forecasts, 614 days
\providecommand{\rhedRegMeanMin}{}\renewcommand{\rhedRegMeanMin}{11.5}  % lowest mean return, mid, % per day across the six full-set pooled forecasts, 614 days
\providecommand{\rhedRegMeanMax}{}\renewcommand{\rhedRegMeanMax}{14.1}  % highest mean return, mid, % per day across the six full-set pooled forecasts, 614 days
\providecommand{\rhedRegMeanXMin}{}\renewcommand{\rhedRegMeanXMin}{8.9}  % lowest mean return, bid/ask, % per day across the six full-set pooled forecasts, 614 days
\providecommand{\rhedRegMeanXMax}{}\renewcommand{\rhedRegMeanXMax}{11.4}  % highest mean return, bid/ask, % per day across the six full-set pooled forecasts, 614 days
\providecommand{\rhedRegRsqCMin}{}\renewcommand{\rhedRegRsqCMin}{0.00}  % lowest R^2 of (c), mid across the six full-set pooled forecasts, 614 days
\providecommand{\rhedRegRsqCMax}{}\renewcommand{\rhedRegRsqCMax}{0.00}  % highest R^2 of (c), mid across the six full-set pooled forecasts, 614 days
\providecommand{\rhedRegCPMaxMin}{}\renewcommand{\rhedRegCPMaxMin}{0.74}  % lowest larger fixed-b p of the two betas of (c), mid across the six full-set pooled forecasts, 614 days
\providecommand{\rhedRegCPMaxMax}{}\renewcommand{\rhedRegCPMaxMax}{0.97}  % highest larger fixed-b p of the two betas of (c), mid across the six full-set pooled forecasts, 614 days
\providecommand{\rhedRegCPMinMin}{}\renewcommand{\rhedRegCPMinMin}{0.56}  % lowest smaller fixed-b p of the two betas of (c), mid across the six full-set pooled forecasts, 614 days
\providecommand{\rhedRegCPMinMax}{}\renewcommand{\rhedRegCPMinMax}{0.82}  % highest smaller fixed-b p of the two betas of (c), mid across the six full-set pooled forecasts, 614 days
\providecommand{\rhedRegAlphaDMin}{}\renewcommand{\rhedRegAlphaDMin}{11.4}  % lowest alpha, (d) joint, mid, % per day across the six full-set pooled forecasts, 614 days
\providecommand{\rhedRegAlphaDMax}{}\renewcommand{\rhedRegAlphaDMax}{13.6}  % highest alpha, (d) joint, mid, % per day across the six full-set pooled forecasts, 614 days
\providecommand{\rhedRegAlphaDXMin}{}\renewcommand{\rhedRegAlphaDXMin}{8.7}  % lowest alpha, (d) joint, bid/ask, % per day across the six full-set pooled forecasts, 614 days
\providecommand{\rhedRegAlphaDXMax}{}\renewcommand{\rhedRegAlphaDXMax}{10.9}  % highest alpha, (d) joint, bid/ask, % per day across the six full-set pooled forecasts, 614 days
\providecommand{\rhedRegAlphaDPMin}{}\renewcommand{\rhedRegAlphaDPMin}{0.004}  % lowest fixed-b p of alpha, (d), mid across the six full-set pooled forecasts, 614 days
\providecommand{\rhedRegAlphaDPMax}{}\renewcommand{\rhedRegAlphaDPMax}{0.026}  % highest fixed-b p of alpha, (d), mid across the six full-set pooled forecasts, 614 days
\providecommand{\rhedRegAlphaDXPMin}{}\renewcommand{\rhedRegAlphaDXPMin}{0.017}  % lowest fixed-b p of alpha, (d), bid/ask across the six full-set pooled forecasts, 614 days
\providecommand{\rhedRegAlphaDXPMax}{}\renewcommand{\rhedRegAlphaDXPMax}{0.076}  % highest fixed-b p of alpha, (d), bid/ask across the six full-set pooled forecasts, 614 days
\providecommand{\rhedRegAvgAlphaM}{}\renewcommand{\rhedRegAvgAlphaM}{12.7}  % alpha of the mean of the six full-set returns, (m), mid, % per day
\providecommand{\rhedRegAvgAlphaMT}{}\renewcommand{\rhedRegAvgAlphaMT}{3.48}  % ... its t
\providecommand{\rhedRegAvgAlphaMP}{}\renewcommand{\rhedRegAvgAlphaMP}{0.002}  % ... its fixed-b p
\providecommand{\rhedRegAvgAlphaA}{}\renewcommand{\rhedRegAvgAlphaA}{12.1}  % alpha of the mean of the six full-set returns, (a), mid, % per day
\providecommand{\rhedRegAvgAlphaAT}{}\renewcommand{\rhedRegAvgAlphaAT}{2.83}  % ... its t
\providecommand{\rhedRegAvgAlphaAP}{}\renewcommand{\rhedRegAvgAlphaAP}{0.0100}  % ... its fixed-b p
\providecommand{\rhedRegAvgAlphaD}{}\renewcommand{\rhedRegAvgAlphaD}{12.5}  % alpha of the mean of the six full-set returns, (d), mid, % per day
\providecommand{\rhedRegAvgAlphaDT}{}\renewcommand{\rhedRegAvgAlphaDT}{2.86}  % ... its t
\providecommand{\rhedRegAvgAlphaDP}{}\renewcommand{\rhedRegAvgAlphaDP}{0.009}  % ... its fixed-b p
\providecommand{\rhedRegAvgAlphaMX}{}\renewcommand{\rhedRegAvgAlphaMX}{10.1}  % alpha of the mean of the six full-set returns, (m), ba, % per day
\providecommand{\rhedRegAvgAlphaMXT}{}\renewcommand{\rhedRegAvgAlphaMXT}{2.78}  % ... its t
\providecommand{\rhedRegAvgAlphaMXP}{}\renewcommand{\rhedRegAvgAlphaMXP}{0.011}  % ... its fixed-b p
\providecommand{\rhedRegAvgAlphaAX}{}\renewcommand{\rhedRegAvgAlphaAX}{9.4}  % alpha of the mean of the six full-set returns, (a), ba, % per day
\providecommand{\rhedRegAvgAlphaAXT}{}\renewcommand{\rhedRegAvgAlphaAXT}{2.21}  % ... its t
\providecommand{\rhedRegAvgAlphaAXP}{}\renewcommand{\rhedRegAvgAlphaAXP}{0.040}  % ... its fixed-b p
\providecommand{\rhedRegAvgAlphaDX}{}\renewcommand{\rhedRegAvgAlphaDX}{9.8}  % alpha of the mean of the six full-set returns, (d), ba, % per day
\providecommand{\rhedRegAvgAlphaDXT}{}\renewcommand{\rhedRegAvgAlphaDXT}{2.26}  % ... its t
\providecommand{\rhedRegAvgAlphaDXP}{}\renewcommand{\rhedRegAvgAlphaDXP}{0.036}  % ... its fixed-b p
\providecommand{\rhedRegAvgBetaShort}{}\renewcommand{\rhedRegAvgBetaShort}{0.26}  % beta of the mean full-set return on the short straddle, (a), mid
\providecommand{\rhedRegAvgBetaShortT}{}\renewcommand{\rhedRegAvgBetaShortT}{1.34}  % ... its t
\providecommand{\rhedRegAvgBetaShortP}{}\renewcommand{\rhedRegAvgBetaShortP}{0.20}  % ... its fixed-b p
\providecommand{\rhedRegAvgRsqA}{}\renewcommand{\rhedRegAvgRsqA}{0.07}  % R^2 of (a) for the mean full-set return
\providecommand{\rhedRegAvgBetaAbs}{}\renewcommand{\rhedRegAvgBetaAbs}{\ensuremath{-}1.68}  % beta of the mean full-set return on |r30| (bp), (b), mid
\providecommand{\rhedRegAvgBetaAbsT}{}\renewcommand{\rhedRegAvgBetaAbsT}{\ensuremath{-}3.26}  % ... its t
\providecommand{\rhedRegShortBetaAbs}{}\renewcommand{\rhedRegShortBetaAbs}{\ensuremath{-}3.96}  % beta of the short straddle on |r30| (bp), (b), mid
\providecommand{\rhedRegShortRsqB}{}\renewcommand{\rhedRegShortRsqB}{0.48}  % R^2 of (b) for the short straddle
\providecommand{\rhedRegAbsRatio}{}\renewcommand{\rhedRegAbsRatio}{43}  % beta on |r30| of the mean full-set return as % of the short straddle's
\providecommand{\rhedRegAvgBetaAbsJ}{}\renewcommand{\rhedRegAvgBetaAbsJ}{\ensuremath{-}1.30}  % beta of the mean full-set return on |r30|, (d) joint, mid
\providecommand{\rhedRegAvgBetaAbsJT}{}\renewcommand{\rhedRegAvgBetaAbsJT}{\ensuremath{-}1.90}  % ... its t
\providecommand{\rhedHarRegAlphaAMin}{}\renewcommand{\rhedHarRegAlphaAMin}{1.1}  % lowest alpha, (a), mid, % per day across the six HAR-benchmark pooled forecasts, 614 days
\providecommand{\rhedHarRegAlphaAMax}{}\renewcommand{\rhedHarRegAlphaAMax}{7.0}  % highest alpha, (a), mid, % per day across the six HAR-benchmark pooled forecasts, 614 days
\providecommand{\rhedHarRegAlphaAPMin}{}\renewcommand{\rhedHarRegAlphaAPMin}{0.082}  % lowest fixed-b p of alpha, (a), mid across the six HAR-benchmark pooled forecasts, 614 days
\providecommand{\rhedHarRegAlphaAPMax}{}\renewcommand{\rhedHarRegAlphaAPMax}{0.745}  % highest fixed-b p of alpha, (a), mid across the six HAR-benchmark pooled forecasts, 614 days
\providecommand{\rhedOlsRegAlphaA}{}\renewcommand{\rhedOlsRegAlphaA}{8.1}  % alpha, (a), mid, closing-interval HAR forecast, % per day
\providecommand{\rhedOlsRegAlphaAP}{}\renewcommand{\rhedOlsRegAlphaAP}{0.051}  % ... its fixed-b p
\providecommand{\rhedRegAlphaAExBestMin}{}\renewcommand{\rhedRegAlphaAExBestMin}{\ensuremath{-}0.5}  % lowest alpha, (a), mid, without the forecast's ten best days, % per day across the six full-set pooled forecasts, 614 days
\providecommand{\rhedRegAlphaAExBestMax}{}\renewcommand{\rhedRegAlphaAExBestMax}{2.1}  % highest alpha, (a), mid, without the forecast's ten best days, % per day across the six full-set pooled forecasts, 614 days
\providecommand{\rhedRegAlphaAExBestPMin}{}\renewcommand{\rhedRegAlphaAExBestPMin}{0.43}  % lowest ... its fixed-b p across the six full-set pooled forecasts, 614 days
\providecommand{\rhedRegAlphaAExBestPMax}{}\renewcommand{\rhedRegAlphaAExBestPMax}{0.97}  % highest ... its fixed-b p across the six full-set pooled forecasts, 614 days
 to the preamble of paper_rehaul.tex, put the
% subsection below after "Results" (sec:rh_econ_results) in writeup/rehaul/07_economics.tex, and add the four floats
% to writeup/rehaul/07_economics_floats.tex:
%   % AUTO-GENERATED by writeup/make_rehaul_econ_decomp.py -- do not edit; re-run the script.
% Table tab:rh_econ_longshort: long and short days of the sign strategy, the gain over the short straddle, and
% the calibration of the forecasts (writeup/make_rehaul_econ_decomp.py, statistics 1 and 3).
\ifdefined\rhedNDays\else% AUTO-GENERATED by writeup/make_rehaul_econ_decomp.py -- do not edit; re-run the script.
% \rhed... macros: where the straddle strategy's profits come from (Section 7 decomposition).
\providecommand{\rhedNDays}{}\renewcommand{\rhedNDays}{614}  % post-warmup option days (= rheSzDays)
\providecommand{\rhedNFull}{}\renewcommand{\rhedNFull}{6}  % pooled intraday forecasts on the full predictor set
\providecommand{\rhedNTop}{}\renewcommand{\rhedNTop}{10}  % best / worst days per forecast
\providecommand{\rhedUncRVAboveIV}{}\renewcommand{\rhedUncRVAboveIV}{28.2}  % % of the 614 days with RV_t > IV_t^2
\providecommand{\rhedUncRVIVMean}{}\renewcommand{\rhedUncRVIVMean}{0.87}  % mean RV_t / IV_t^2, 614 days
\providecommand{\rhedUncRVIVMed}{}\renewcommand{\rhedUncRVIVMed}{0.77}  % median RV_t / IV_t^2, 614 days
\providecommand{\rhedUncShortMean}{}\renewcommand{\rhedUncShortMean}{2.4}  % mean mid return of the unconditional short straddle, % of premium
\providecommand{\rhedUncShortHit}{}\renewcommand{\rhedUncShortHit}{71.8}  % % of days on which q = -1 matches sign(RV - IV^2) (always short)
\providecommand{\rhedLongNMin}{}\renewcommand{\rhedLongNMin}{117}  % lowest number of long days across the six full-set pooled forecasts, 614 days
\providecommand{\rhedLongNMax}{}\renewcommand{\rhedLongNMax}{146}  % highest number of long days across the six full-set pooled forecasts, 614 days
\providecommand{\rhedLongPctMin}{}\renewcommand{\rhedLongPctMin}{19}  % lowest % long days across the six full-set pooled forecasts, 614 days
\providecommand{\rhedLongPctMax}{}\renewcommand{\rhedLongPctMax}{24}  % highest % long days across the six full-set pooled forecasts, 614 days
\providecommand{\rhedHarLongNMin}{}\renewcommand{\rhedHarLongNMin}{150}  % lowest number of long days across the six HAR-benchmark pooled forecasts, 614 days
\providecommand{\rhedHarLongNMax}{}\renewcommand{\rhedHarLongNMax}{193}  % highest number of long days across the six HAR-benchmark pooled forecasts, 614 days
\providecommand{\rhedLongMidMin}{}\renewcommand{\rhedLongMidMin}{19}  % lowest mean mid return on long days, % across the six full-set pooled forecasts, 614 days
\providecommand{\rhedLongMidMax}{}\renewcommand{\rhedLongMidMax}{30}  % highest mean mid return on long days, % across the six full-set pooled forecasts, 614 days
\providecommand{\rhedLongBaMin}{}\renewcommand{\rhedLongBaMin}{16}  % lowest mean bid/ask return on long days, % across the six full-set pooled forecasts, 614 days
\providecommand{\rhedLongBaMax}{}\renewcommand{\rhedLongBaMax}{26}  % highest mean bid/ask return on long days, % across the six full-set pooled forecasts, 614 days
\providecommand{\rhedShortMidMin}{}\renewcommand{\rhedShortMidMin}{9}  % lowest mean mid return on short days, % across the six full-set pooled forecasts, 614 days
\providecommand{\rhedShortMidMax}{}\renewcommand{\rhedShortMidMax}{11}  % highest mean mid return on short days, % across the six full-set pooled forecasts, 614 days
\providecommand{\rhedShortBaMin}{}\renewcommand{\rhedShortBaMin}{7}  % lowest mean bid/ask return on short days, % across the six full-set pooled forecasts, 614 days
\providecommand{\rhedShortBaMax}{}\renewcommand{\rhedShortBaMax}{8}  % highest mean bid/ask return on short days, % across the six full-set pooled forecasts, 614 days
\providecommand{\rhedHarLongMidMin}{}\renewcommand{\rhedHarLongMidMin}{0}  % lowest mean mid return on long days, % across the six HAR-benchmark pooled forecasts, 614 days
\providecommand{\rhedHarLongMidMax}{}\renewcommand{\rhedHarLongMidMax}{10}  % highest mean mid return on long days, % across the six HAR-benchmark pooled forecasts, 614 days
\providecommand{\rhedHarShortMidMin}{}\renewcommand{\rhedHarShortMidMin}{3}  % lowest mean mid return on short days, % across the six HAR-benchmark pooled forecasts, 614 days
\providecommand{\rhedHarShortMidMax}{}\renewcommand{\rhedHarShortMidMax}{7}  % highest mean mid return on short days, % across the six HAR-benchmark pooled forecasts, 614 days
\providecommand{\rhedLongHitMin}{}\renewcommand{\rhedLongHitMin}{47}  % lowest % of long days with RV > IV^2 across the six full-set pooled forecasts, 614 days
\providecommand{\rhedLongHitMax}{}\renewcommand{\rhedLongHitMax}{53}  % highest % of long days with RV > IV^2 across the six full-set pooled forecasts, 614 days
\providecommand{\rhedHarLongHitMin}{}\renewcommand{\rhedHarLongHitMin}{46}  % lowest % of long days with RV > IV^2 across the six HAR-benchmark pooled forecasts, 614 days
\providecommand{\rhedHarLongHitMax}{}\renewcommand{\rhedHarLongHitMax}{48}  % highest % of long days with RV > IV^2 across the six HAR-benchmark pooled forecasts, 614 days
\providecommand{\rhedLongPosMin}{}\renewcommand{\rhedLongPosMin}{42}  % lowest % of long days with a positive return across the six full-set pooled forecasts, 614 days
\providecommand{\rhedLongPosMax}{}\renewcommand{\rhedLongPosMax}{45}  % highest % of long days with a positive return across the six full-set pooled forecasts, 614 days
\providecommand{\rhedHitMin}{}\renewcommand{\rhedHitMin}{70.7}  % lowest % of days on which q matches sign(RV - IV^2) across the six full-set pooled forecasts, 614 days
\providecommand{\rhedHitMax}{}\renewcommand{\rhedHitMax}{73.0}  % highest % of days on which q matches sign(RV - IV^2) across the six full-set pooled forecasts, 614 days
\providecommand{\rhedPosMin}{}\renewcommand{\rhedPosMin}{58.6}  % lowest % of days with a positive mid return across the six full-set pooled forecasts, 614 days
\providecommand{\rhedPosMax}{}\renewcommand{\rhedPosMax}{60.4}  % highest % of days with a positive mid return across the six full-set pooled forecasts, 614 days
\providecommand{\rhedLongShareMin}{}\renewcommand{\rhedLongShareMin}{39}  % lowest % of the summed mid return earned on long days across the six full-set pooled forecasts, 614 days
\providecommand{\rhedLongShareMax}{}\renewcommand{\rhedLongShareMax}{41}  % highest % of the summed mid return earned on long days across the six full-set pooled forecasts, 614 days
\providecommand{\rhedLongShareXMin}{}\renewcommand{\rhedLongShareXMin}{42}  % lowest % of the summed bid/ask return earned on long days across the six full-set pooled forecasts, 614 days
\providecommand{\rhedLongShareXMax}{}\renewcommand{\rhedLongShareXMax}{45}  % highest % of the summed bid/ask return earned on long days across the six full-set pooled forecasts, 614 days
\providecommand{\rhedMeanMin}{}\renewcommand{\rhedMeanMin}{11.5}  % lowest mean mid return, % across the six full-set pooled forecasts, 614 days
\providecommand{\rhedMeanMax}{}\renewcommand{\rhedMeanMax}{14.1}  % highest mean mid return, % across the six full-set pooled forecasts, 614 days
\providecommand{\rhedGainMin}{}\renewcommand{\rhedGainMin}{9.1}  % lowest mean gain over the short straddle (2 R 1(long)), % across the six full-set pooled forecasts, 614 days
\providecommand{\rhedGainMax}{}\renewcommand{\rhedGainMax}{11.6}  % highest mean gain over the short straddle (2 R 1(long)), % across the six full-set pooled forecasts, 614 days
\providecommand{\rhedGainShareMin}{}\renewcommand{\rhedGainShareMin}{79}  % lowest gain over short / mean return, % across the six full-set pooled forecasts, 614 days
\providecommand{\rhedGainShareMax}{}\renewcommand{\rhedGainShareMax}{83}  % highest gain over short / mean return, % across the six full-set pooled forecasts, 614 days
\providecommand{\rhedHarGainMin}{}\renewcommand{\rhedHarGainMin}{0.0}  % lowest mean gain over the short straddle, % across the six HAR-benchmark pooled forecasts, 614 days
\providecommand{\rhedHarGainMax}{}\renewcommand{\rhedHarGainMax}{5.0}  % highest mean gain over the short straddle, % across the six HAR-benchmark pooled forecasts, 614 days
\providecommand{\rhedOlsLongN}{}\renewcommand{\rhedOlsLongN}{198}  % long days of the closing-interval HAR forecast (sub_ols_baseline)
\providecommand{\rhedOlsLongMid}{}\renewcommand{\rhedOlsLongMid}{9}  % ... its mean mid return on long days, %
\providecommand{\rhedOlsShortMid}{}\renewcommand{\rhedOlsShortMid}{8}  % ... its mean mid return on short days, %
\providecommand{\rhedOlsLongHit}{}\renewcommand{\rhedOlsLongHit}{46}  % ... % of its long days with RV > IV^2
\providecommand{\rhedOlsSh}{}\renewcommand{\rhedOlsSh}{1.11}  % ... its Sharpe ratio, mid (= rhrOlsSharpe)
\providecommand{\rhedCalLevelMin}{}\renewcommand{\rhedCalLevelMin}{0.95}  % lowest mean forecast / mean RV across the six full-set pooled forecasts, 614 days
\providecommand{\rhedCalLevelMax}{}\renewcommand{\rhedCalLevelMax}{1.01}  % highest mean forecast / mean RV across the six full-set pooled forecasts, 614 days
\providecommand{\rhedCalScaleMin}{}\renewcommand{\rhedCalScaleMin}{1.09}  % lowest mean of RV / forecast (QLIKE-optimal scale) across the six full-set pooled forecasts, 614 days
\providecommand{\rhedCalScaleMax}{}\renewcommand{\rhedCalScaleMax}{1.20}  % highest mean of RV / forecast (QLIKE-optimal scale) across the six full-set pooled forecasts, 614 days
\providecommand{\rhedHarCalScaleMin}{}\renewcommand{\rhedHarCalScaleMin}{1.04}  % lowest mean of RV / forecast across the six HAR-benchmark pooled forecasts, 614 days
\providecommand{\rhedHarCalScaleMax}{}\renewcommand{\rhedHarCalScaleMax}{1.10}  % highest mean of RV / forecast across the six HAR-benchmark pooled forecasts, 614 days
\providecommand{\rhedMedFcIVMin}{}\renewcommand{\rhedMedFcIVMin}{0.75}  % lowest median forecast / IV^2 across the six full-set pooled forecasts, 614 days
\providecommand{\rhedMedFcIVMax}{}\renewcommand{\rhedMedFcIVMax}{0.82}  % highest median forecast / IV^2 across the six full-set pooled forecasts, 614 days
\providecommand{\rhedBcScaleLoMin}{}\renewcommand{\rhedBcScaleLoMin}{1.06}  % lowest lowest causal scale c_t over the 614 days across the six full-set pooled forecasts, 614 days
\providecommand{\rhedBcScaleLoMax}{}\renewcommand{\rhedBcScaleLoMax}{1.15}  % highest lowest causal scale c_t over the 614 days across the six full-set pooled forecasts, 614 days
\providecommand{\rhedBcScaleHiMin}{}\renewcommand{\rhedBcScaleHiMin}{1.10}  % lowest highest causal scale c_t over the 614 days across the six full-set pooled forecasts, 614 days
\providecommand{\rhedBcScaleHiMax}{}\renewcommand{\rhedBcScaleHiMax}{1.22}  % highest highest causal scale c_t over the 614 days across the six full-set pooled forecasts, 614 days
\providecommand{\rhedBcFlipMin}{}\renewcommand{\rhedBcFlipMin}{55}  % lowest positions changed by the bias correction across the six full-set pooled forecasts, 614 days
\providecommand{\rhedBcFlipMax}{}\renewcommand{\rhedBcFlipMax}{120}  % highest positions changed by the bias correction across the six full-set pooled forecasts, 614 days
\providecommand{\rhedBcLongMin}{}\renewcommand{\rhedBcLongMin}{201}  % lowest long days under the bias correction across the six full-set pooled forecasts, 614 days
\providecommand{\rhedBcLongMax}{}\renewcommand{\rhedBcLongMax}{241}  % highest long days under the bias correction across the six full-set pooled forecasts, 614 days
\providecommand{\rhedBcFlipRMin}{}\renewcommand{\rhedBcFlipRMin}{\ensuremath{-}24}  % lowest mean mid straddle return R on the changed days, % across the six full-set pooled forecasts, 614 days
\providecommand{\rhedBcFlipRMax}{}\renewcommand{\rhedBcFlipRMax}{\ensuremath{-}8}  % highest mean mid straddle return R on the changed days, % across the six full-set pooled forecasts, 614 days
\providecommand{\rhedBcShMin}{}\renewcommand{\rhedBcShMin}{0.92}  % lowest Sharpe ratio mid under the bias correction (diagnostic) across the six full-set pooled forecasts, 614 days
\providecommand{\rhedBcShMax}{}\renewcommand{\rhedBcShMax}{1.32}  % highest Sharpe ratio mid under the bias correction (diagnostic) across the six full-set pooled forecasts, 614 days
\providecommand{\rhedBcShXMin}{}\renewcommand{\rhedBcShXMin}{0.57}  % lowest Sharpe ratio bid/ask under the bias correction across the six full-set pooled forecasts, 614 days
\providecommand{\rhedBcShXMax}{}\renewcommand{\rhedBcShXMax}{0.97}  % highest Sharpe ratio bid/ask under the bias correction across the six full-set pooled forecasts, 614 days
\providecommand{\rhedLevFlipMin}{}\renewcommand{\rhedLevFlipMin}{4}  % lowest positions changed by the level correction (sum RV / sum fc) across the six full-set pooled forecasts, 614 days
\providecommand{\rhedLevFlipMax}{}\renewcommand{\rhedLevFlipMax}{19}  % highest positions changed by the level correction (sum RV / sum fc) across the six full-set pooled forecasts, 614 days
\providecommand{\rhedLevShMin}{}\renewcommand{\rhedLevShMin}{1.61}  % lowest Sharpe ratio mid under the level correction across the six full-set pooled forecasts, 614 days
\providecommand{\rhedLevShMax}{}\renewcommand{\rhedLevShMax}{2.01}  % highest Sharpe ratio mid under the level correction across the six full-set pooled forecasts, 614 days
\providecommand{\rhedHarBcShMin}{}\renewcommand{\rhedHarBcShMin}{0.76}  % lowest Sharpe ratio mid under the bias correction across the six HAR-benchmark pooled forecasts, 614 days
\providecommand{\rhedHarBcShMax}{}\renewcommand{\rhedHarBcShMax}{1.05}  % highest Sharpe ratio mid under the bias correction across the six HAR-benchmark pooled forecasts, 614 days
\providecommand{\rhedBcShUpN}{}\renewcommand{\rhedBcShUpN}{0}  % full-set forecasts whose mid SR rises under the bias correction
\providecommand{\rhedBcShDownN}{}\renewcommand{\rhedBcShDownN}{6}  % full-set forecasts whose mid SR falls under the bias correction
\providecommand{\rhedHarBcShUpN}{}\renewcommand{\rhedHarBcShUpN}{2}  % HAR-benchmark forecasts whose mid SR rises under the bias correction
\providecommand{\rhedMoveSqIV}{}\renewcommand{\rhedMoveSqIV}{1.01}  % mean of the squared S&P 500 half-hour log return over IV_t^2, 614 days
\providecommand{\rhedMoveSqRV}{}\renewcommand{\rhedMoveSqRV}{1.15}  % mean squared S&P 500 half-hour log return over mean RV_t, 614 days
\providecommand{\rhedOlsCalScale}{}\renewcommand{\rhedOlsCalScale}{1.01}  % mean of RV / forecast, closing-interval HAR forecast
\providecommand{\rhedOlsBcSh}{}\renewcommand{\rhedOlsBcSh}{1.12}  % Sharpe ratio mid under the bias correction, closing-interval HAR
\providecommand{\rhedBcFlipRidge}{}\renewcommand{\rhedBcFlipRidge}{105}  % positions changed by the bias correction, poolpt13_ridge_all_features
\providecommand{\rhedBcShRidge}{}\renewcommand{\rhedBcShRidge}{1.17}  % Sharpe ratio mid under the bias correction, poolpt13_ridge_all_features
\providecommand{\rhedBcFlipLasso}{}\renewcommand{\rhedBcFlipLasso}{120}  % positions changed by the bias correction, poolpt13_lasso_all_features
\providecommand{\rhedBcShLasso}{}\renewcommand{\rhedBcShLasso}{1.31}  % Sharpe ratio mid under the bias correction, poolpt13_lasso_all_features
\providecommand{\rhedBcFlipEnet}{}\renewcommand{\rhedBcFlipEnet}{108}  % positions changed by the bias correction, poolpt13_enet_all_features
\providecommand{\rhedBcShEnet}{}\renewcommand{\rhedBcShEnet}{0.99}  % Sharpe ratio mid under the bias correction, poolpt13_enet_all_features
\providecommand{\rhedBcFlipLgbm}{}\renewcommand{\rhedBcFlipLgbm}{55}  % positions changed by the bias correction, pooltree_lgbm_all_features
\providecommand{\rhedBcShLgbm}{}\renewcommand{\rhedBcShLgbm}{1.32}  % Sharpe ratio mid under the bias correction, pooltree_lgbm_all_features
\providecommand{\rhedBcFlipXgb}{}\renewcommand{\rhedBcFlipXgb}{61}  % positions changed by the bias correction, pooltree_xgb_all_features
\providecommand{\rhedBcShXgb}{}\renewcommand{\rhedBcShXgb}{1.15}  % Sharpe ratio mid under the bias correction, pooltree_xgb_all_features
\providecommand{\rhedBcFlipRf}{}\renewcommand{\rhedBcFlipRf}{90}  % positions changed by the bias correction, pooltree_rf_all_features
\providecommand{\rhedBcShRf}{}\renewcommand{\rhedBcShRf}{0.92}  % Sharpe ratio mid under the bias correction, pooltree_rf_all_features
\providecommand{\rhedQMonoN}{}\renewcommand{\rhedQMonoN}{6}  % full-set forecasts whose mean RV/IV^2 rises across all five quintiles
\providecommand{\rhedHarQMonoN}{}\renewcommand{\rhedHarQMonoN}{4}  % HAR-benchmark forecasts whose mean RV/IV^2 rises across all five quintiles
\providecommand{\rhedQRPatternN}{}\renewcommand{\rhedQRPatternN}{6}  % full-set forecasts with mean R < 0 in quintiles 1-4 and > 0 in 5
\providecommand{\rhedHarQRPatternN}{}\renewcommand{\rhedHarQRPatternN}{0}  % HAR-benchmark forecasts with that pattern
\providecommand{\rhedQOneRVIV}{}\renewcommand{\rhedQOneRVIV}{0.73}  % mean RV/IV^2 in the lowest quintile, mean of the six full-set
\providecommand{\rhedQFiveRVIV}{}\renewcommand{\rhedQFiveRVIV}{1.05}  % mean RV/IV^2 in the highest quintile, mean of the six full-set
\providecommand{\rhedQFiveR}{}\renewcommand{\rhedQFiveR}{26}  % mean straddle return R in the highest quintile, %, mean of the six full-set
\providecommand{\rhedQFiveFc}{}\renewcommand{\rhedQFiveFc}{1.18}  % geometric-mean forecast / IV^2 in the highest quintile, mean of the six full-set
\providecommand{\rhedQFourFc}{}\renewcommand{\rhedQFourFc}{0.93}  % ... in the fourth quintile
\providecommand{\rhedQFiveLong}{}\renewcommand{\rhedQFiveLong}{99}  % % long days in the highest quintile, mean of the six full-set
\providecommand{\rhedQFiveHit}{}\renewcommand{\rhedQFiveHit}{52}  % hit rate in the highest quintile, %, mean of the six full-set
\providecommand{\rhedQOneHit}{}\renewcommand{\rhedQOneHit}{88}  % hit rate in the lowest quintile, %, mean of the six full-set
\providecommand{\rhedQFiveShare}{}\renewcommand{\rhedQFiveShare}{43}  % % of the summed mid return earned in the highest quintile, six full-set forecasts
\providecommand{\rhedQLowRMin}{}\renewcommand{\rhedQLowRMin}{\ensuremath{-}13}  % most negative mean straddle return R in quintiles 1-4, %, mean of the six
\providecommand{\rhedQLowRMax}{}\renewcommand{\rhedQLowRMax}{\ensuremath{-}6}  % least negative mean straddle return R in quintiles 1-4, %, mean of the six
\providecommand{\rhedHarQOneRVIV}{}\renewcommand{\rhedHarQOneRVIV}{0.74}  % mean RV/IV^2 in the lowest quintile, mean of the six Har
\providecommand{\rhedHarQFiveRVIV}{}\renewcommand{\rhedHarQFiveRVIV}{1.02}  % mean RV/IV^2 in the highest quintile, mean of the six Har
\providecommand{\rhedHarQFiveR}{}\renewcommand{\rhedHarQFiveR}{7}  % mean straddle return R in the highest quintile, %, mean of the six Har
\providecommand{\rhedHarQFiveFc}{}\renewcommand{\rhedHarQFiveFc}{1.23}  % geometric-mean forecast / IV^2 in the highest quintile, mean of the six Har
\providecommand{\rhedHarQFourFc}{}\renewcommand{\rhedHarQFourFc}{0.98}  % ... in the fourth quintile
\providecommand{\rhedHarQFiveLong}{}\renewcommand{\rhedHarQFiveLong}{100}  % % long days in the highest quintile, mean of the six Har
\providecommand{\rhedHarQFiveHit}{}\renewcommand{\rhedHarQFiveHit}{47}  % hit rate in the highest quintile, %, mean of the six Har
\providecommand{\rhedHarQOneHit}{}\renewcommand{\rhedHarQOneHit}{89}  % hit rate in the lowest quintile, %, mean of the six Har
\providecommand{\rhedHarQFiveShare}{}\renewcommand{\rhedHarQFiveShare}{23}  % % of the summed mid return earned in the highest quintile, six Har forecasts
\providecommand{\rhedHarQLowRMin}{}\renewcommand{\rhedHarQLowRMin}{\ensuremath{-}13}  % most negative mean straddle return R in quintiles 1-4, %, mean of the six
\providecommand{\rhedHarQLowRMax}{}\renewcommand{\rhedHarQLowRMax}{5}  % least negative mean straddle return R in quintiles 1-4, %, mean of the six
\providecommand{\rhedQFiveShareRngMin}{}\renewcommand{\rhedQFiveShareRngMin}{40}  % lowest % of the summed mid return in quintile 5 across the six full-set pooled forecasts, 614 days
\providecommand{\rhedQFiveShareRngMax}{}\renewcommand{\rhedQFiveShareRngMax}{48}  % highest % of the summed mid return in quintile 5 across the six full-set pooled forecasts, 614 days
\providecommand{\rhedBestUnion}{}\renewcommand{\rhedBestUnion}{11}  % days in the union of the six full-set forecasts' ten best days
\providecommand{\rhedWorstUnion}{}\renewcommand{\rhedWorstUnion}{11}  % days in the union of the six full-set forecasts' ten worst days
\providecommand{\rhedBestAllSix}{}\renewcommand{\rhedBestAllSix}{9}  % days among the ten best of all six forecasts
\providecommand{\rhedWorstAllSix}{}\renewcommand{\rhedWorstAllSix}{9}  % days among the ten worst of all six forecasts
\providecommand{\rhedTopDate}{}\renewcommand{\rhedTopDate}{22~November~2021}  % the day of the largest straddle return among the best days
\providecommand{\rhedTopR}{}\renewcommand{\rhedTopR}{1{,}032}  % ... the straddle's mid return that day, % of premium
\providecommand{\rhedTopRVIV}{}\renewcommand{\rhedTopRVIV}{2.79}  % ... RV / IV^2 that day
\providecommand{\rhedTopMove}{}\renewcommand{\rhedTopMove}{\ensuremath{-}0.83}  % ... S&P 500 15:30-16:00 log return that day, %
\providecommand{\rhedTopShareMin}{}\renewcommand{\rhedTopShareMin}{12}  % lowest % of a full-set forecast's summed mid return earned on that one day
\providecommand{\rhedTopShareMax}{}\renewcommand{\rhedTopShareMax}{15}  % highest % of a full-set forecast's summed mid return earned on that one day
\providecommand{\rhedBestFomcN}{}\renewcommand{\rhedBestFomcN}{1}  % best-day union: FOMC statement days
\providecommand{\rhedWorstFomcN}{}\renewcommand{\rhedWorstFomcN}{1}  % worst-day union: FOMC statement days
\providecommand{\rhedWorstMeN}{}\renewcommand{\rhedWorstMeN}{3}  % worst-day union: month ends
\providecommand{\rhedBestMeN}{}\renewcommand{\rhedBestMeN}{0}  % best-day union: month ends
\providecommand{\rhedBestOpexN}{}\renewcommand{\rhedBestOpexN}{0}  % best-day union: monthly expirations
\providecommand{\rhedWorstOpexN}{}\renewcommand{\rhedWorstOpexN}{0}  % worst-day union: monthly expirations
\providecommand{\rhedBestNoTagN}{}\renewcommand{\rhedBestNoTagN}{8}  % best-day union: days without any event tag (calendar as covered)
\providecommand{\rhedBestAfterRelN}{}\renewcommand{\rhedBestAfterRelN}{5}  % best-day union: days after the release calendar ends (dagger)
\providecommand{\rhedWorstAfterRelN}{}\renewcommand{\rhedWorstAfterRelN}{4}  % worst-day union: days after the release calendar ends
\providecommand{\rhedDShortAllMin}{}\renewcommand{\rhedDShortAllMin}{1.21}  % lowest Delta SR mid vs the short straddle, all days across the six full-set pooled forecasts, 614 days
\providecommand{\rhedDShortAllMax}{}\renewcommand{\rhedDShortAllMax}{1.56}  % highest Delta SR mid vs the short straddle, all days across the six full-set pooled forecasts, 614 days
\providecommand{\rhedDShortExBestMin}{}\renewcommand{\rhedDShortExBestMin}{\ensuremath{-}0.54}  % lowest Delta SR mid vs the short straddle without the forecast's ten best days across the six full-set pooled forecasts, 614 days
\providecommand{\rhedDShortExBestMax}{}\renewcommand{\rhedDShortExBestMax}{\ensuremath{-}0.12}  % highest Delta SR mid vs the short straddle without the forecast's ten best days across the six full-set pooled forecasts, 614 days
\providecommand{\rhedDShortExBothMin}{}\renewcommand{\rhedDShortExBothMin}{\ensuremath{-}0.64}  % lowest ... without its ten best and ten worst days across the six full-set pooled forecasts, 614 days
\providecommand{\rhedDShortExBothMax}{}\renewcommand{\rhedDShortExBothMax}{\ensuremath{-}0.14}  % highest ... without its ten best and ten worst days across the six full-set pooled forecasts, 614 days
\providecommand{\rhedBestShareMin}{}\renewcommand{\rhedBestShareMin}{107}  % lowest % of the summed gain over the short straddle earned on the ten best days across the six full-set pooled forecasts, 614 days
\providecommand{\rhedBestShareMax}{}\renewcommand{\rhedBestShareMax}{138}  % highest % of the summed gain over the short straddle earned on the ten best days across the six full-set pooled forecasts, 614 days
\providecommand{\rhedDHarAllMin}{}\renewcommand{\rhedDHarAllMin}{0.56}  % lowest Delta SR mid vs the HAR twin, all days across the six full-set pooled forecasts, 614 days
\providecommand{\rhedDHarAllMax}{}\renewcommand{\rhedDHarAllMax}{1.56}  % highest Delta SR mid vs the HAR twin, all days across the six full-set pooled forecasts, 614 days
\providecommand{\rhedDHarExBestMin}{}\renewcommand{\rhedDHarExBestMin}{0.30}  % lowest Delta SR mid vs the HAR twin without the forecast's ten best days across the six full-set pooled forecasts, 614 days
\providecommand{\rhedDHarExBestMax}{}\renewcommand{\rhedDHarExBestMax}{0.83}  % highest Delta SR mid vs the HAR twin without the forecast's ten best days across the six full-set pooled forecasts, 614 days
\providecommand{\rhedDHarExBothMin}{}\renewcommand{\rhedDHarExBothMin}{0.36}  % lowest ... without its ten best and ten worst days across the six full-set pooled forecasts, 614 days
\providecommand{\rhedDHarExBothMax}{}\renewcommand{\rhedDHarExBothMax}{0.97}  % highest ... without its ten best and ten worst days across the six full-set pooled forecasts, 614 days
\providecommand{\rhedHarBestShareMin}{}\renewcommand{\rhedHarBestShareMin}{47}  % lowest % of the summed gain over the HAR twin earned on the ten best days across the six full-set pooled forecasts, 614 days
\providecommand{\rhedHarBestShareMax}{}\renewcommand{\rhedHarBestShareMax}{61}  % highest % of the summed gain over the HAR twin earned on the ten best days across the six full-set pooled forecasts, 614 days
\providecommand{\rhedDShortExBestNeg}{}\renewcommand{\rhedDShortExBestNeg}{6}  % full-set forecasts whose Delta SR vs short is negative without the ten best days
\providecommand{\rhedDHarExBestNeg}{}\renewcommand{\rhedDHarExBestNeg}{0}  % full-set forecasts whose Delta SR vs the HAR twin is negative without the ten best days
\providecommand{\rhedFomcN}{}\renewcommand{\rhedFomcN}{21}  % FOMC days among the 614 post-warmup days
\providecommand{\rhedFomcR}{}\renewcommand{\rhedFomcR}{25}  % mean mid straddle return R on FOMC days, %
\providecommand{\rhedFomcRT}{}\renewcommand{\rhedFomcRT}{1.02}  % ... its t (i.i.d.; the days are weeks apart)
\providecommand{\rhedFomcLong}{}\renewcommand{\rhedFomcLong}{25}  % % of FOMC days long, mean of the six full-set forecasts
\providecommand{\rhedFomcStrat}{}\renewcommand{\rhedFomcStrat}{\ensuremath{-}6}  % mean mid strategy return on FOMC days, %, mean of the six full-set
\providecommand{\rhedFomcOtherR}{}\renewcommand{\rhedFomcOtherR}{\ensuremath{-}3}  % mean straddle R on the other (covered) days, %
\providecommand{\rhedFomcOtherStrat}{}\renewcommand{\rhedFomcOtherStrat}{13}  % mean strategy return on the other (covered) days, %
\providecommand{\rhedFomcRVIV}{}\renewcommand{\rhedFomcRVIV}{0.76}  % mean RV_t / IV_t^2 on FOMC days
\providecommand{\rhedFomcMoveSqIV}{}\renewcommand{\rhedFomcMoveSqIV}{1.43}  % mean squared S&P 500 half-hour log return / IV_t^2 on FOMC days
\providecommand{\rhedFomcShare}{}\renewcommand{\rhedFomcShare}{\ensuremath{-}2}  % % of the summed mid return earned on FOMC days, mean of the six full-set forecasts
\providecommand{\rhedCpiN}{}\renewcommand{\rhedCpiN}{25}  % CPI days among the 614 post-warmup days
\providecommand{\rhedCpiR}{}\renewcommand{\rhedCpiR}{\ensuremath{-}12}  % mean mid straddle return R on CPI days, %
\providecommand{\rhedCpiRT}{}\renewcommand{\rhedCpiRT}{\ensuremath{-}0.55}  % ... its t (i.i.d.; the days are weeks apart)
\providecommand{\rhedCpiLong}{}\renewcommand{\rhedCpiLong}{16}  % % of CPI days long, mean of the six full-set forecasts
\providecommand{\rhedCpiStrat}{}\renewcommand{\rhedCpiStrat}{22}  % mean mid strategy return on CPI days, %, mean of the six full-set
\providecommand{\rhedCpiOtherR}{}\renewcommand{\rhedCpiOtherR}{\ensuremath{-}6}  % mean straddle R on the other (covered) days, %
\providecommand{\rhedCpiOtherStrat}{}\renewcommand{\rhedCpiOtherStrat}{11}  % mean strategy return on the other (covered) days, %
\providecommand{\rhedCpiRVIV}{}\renewcommand{\rhedCpiRVIV}{0.80}  % mean RV_t / IV_t^2 on CPI days
\providecommand{\rhedCpiMoveSqIV}{}\renewcommand{\rhedCpiMoveSqIV}{0.86}  % mean squared S&P 500 half-hour log return / IV_t^2 on CPI days
\providecommand{\rhedCpiShare}{}\renewcommand{\rhedCpiShare}{7}  % % of the summed mid return earned on CPI days, mean of the six full-set forecasts
\providecommand{\rhedNfpN}{}\renewcommand{\rhedNfpN}{27}  % NFP days among the 614 post-warmup days
\providecommand{\rhedNfpR}{}\renewcommand{\rhedNfpR}{1}  % mean mid straddle return R on NFP days, %
\providecommand{\rhedNfpRT}{}\renewcommand{\rhedNfpRT}{0.06}  % ... its t (i.i.d.; the days are weeks apart)
\providecommand{\rhedNfpLong}{}\renewcommand{\rhedNfpLong}{12}  % % of NFP days long, mean of the six full-set forecasts
\providecommand{\rhedNfpStrat}{}\renewcommand{\rhedNfpStrat}{20}  % mean mid strategy return on NFP days, %, mean of the six full-set
\providecommand{\rhedNfpOtherR}{}\renewcommand{\rhedNfpOtherR}{\ensuremath{-}6}  % mean straddle R on the other (covered) days, %
\providecommand{\rhedNfpOtherStrat}{}\renewcommand{\rhedNfpOtherStrat}{11}  % mean strategy return on the other (covered) days, %
\providecommand{\rhedNfpRVIV}{}\renewcommand{\rhedNfpRVIV}{0.73}  % mean RV_t / IV_t^2 on NFP days
\providecommand{\rhedNfpMoveSqIV}{}\renewcommand{\rhedNfpMoveSqIV}{0.93}  % mean squared S&P 500 half-hour log return / IV_t^2 on NFP days
\providecommand{\rhedNfpShare}{}\renewcommand{\rhedNfpShare}{7}  % % of the summed mid return earned on NFP days, mean of the six full-set forecasts
\providecommand{\rhedOpexN}{}\renewcommand{\rhedOpexN}{33}  % OpEx days among the 614 post-warmup days
\providecommand{\rhedOpexR}{}\renewcommand{\rhedOpexR}{\ensuremath{-}45}  % mean mid straddle return R on OpEx days, %
\providecommand{\rhedOpexRT}{}\renewcommand{\rhedOpexRT}{\ensuremath{-}4.01}  % ... its t (i.i.d.; the days are weeks apart)
\providecommand{\rhedOpexLong}{}\renewcommand{\rhedOpexLong}{9}  % % of OpEx days long, mean of the six full-set forecasts
\providecommand{\rhedOpexStrat}{}\renewcommand{\rhedOpexStrat}{36}  % mean mid strategy return on OpEx days, %, mean of the six full-set
\providecommand{\rhedOpexOtherR}{}\renewcommand{\rhedOpexOtherR}{0}  % mean straddle R on the other (covered) days, %
\providecommand{\rhedOpexOtherStrat}{}\renewcommand{\rhedOpexOtherStrat}{11}  % mean strategy return on the other (covered) days, %
\providecommand{\rhedOpexRVIV}{}\renewcommand{\rhedOpexRVIV}{0.83}  % mean RV_t / IV_t^2 on OpEx days
\providecommand{\rhedOpexMoveSqIV}{}\renewcommand{\rhedOpexMoveSqIV}{0.40}  % mean squared S&P 500 half-hour log return / IV_t^2 on OpEx days
\providecommand{\rhedOpexShare}{}\renewcommand{\rhedOpexShare}{15}  % % of the summed mid return earned on OpEx days, mean of the six full-set forecasts
\providecommand{\rhedTwN}{}\renewcommand{\rhedTwN}{11}  % TW days among the 614 post-warmup days
\providecommand{\rhedTwR}{}\renewcommand{\rhedTwR}{\ensuremath{-}52}  % mean mid straddle return R on TW days, %
\providecommand{\rhedTwRT}{}\renewcommand{\rhedTwRT}{\ensuremath{-}3.98}  % ... its t (i.i.d.; the days are weeks apart)
\providecommand{\rhedTwLong}{}\renewcommand{\rhedTwLong}{8}  % % of TW days long, mean of the six full-set forecasts
\providecommand{\rhedTwStrat}{}\renewcommand{\rhedTwStrat}{37}  % mean mid strategy return on TW days, %, mean of the six full-set
\providecommand{\rhedTwOtherR}{}\renewcommand{\rhedTwOtherR}{\ensuremath{-}2}  % mean straddle R on the other (covered) days, %
\providecommand{\rhedTwOtherStrat}{}\renewcommand{\rhedTwOtherStrat}{12}  % mean strategy return on the other (covered) days, %
\providecommand{\rhedTwRVIV}{}\renewcommand{\rhedTwRVIV}{0.89}  % mean RV_t / IV_t^2 on TW days
\providecommand{\rhedTwMoveSqIV}{}\renewcommand{\rhedTwMoveSqIV}{0.31}  % mean squared S&P 500 half-hour log return / IV_t^2 on TW days
\providecommand{\rhedTwShare}{}\renewcommand{\rhedTwShare}{5}  % % of the summed mid return earned on TW days, mean of the six full-set forecasts
\providecommand{\rhedMeN}{}\renewcommand{\rhedMeN}{33}  % ME days among the 614 post-warmup days
\providecommand{\rhedMeR}{}\renewcommand{\rhedMeR}{70}  % mean mid straddle return R on ME days, %
\providecommand{\rhedMeRT}{}\renewcommand{\rhedMeRT}{2.79}  % ... its t (i.i.d.; the days are weeks apart)
\providecommand{\rhedMeLong}{}\renewcommand{\rhedMeLong}{6}  % % of ME days long, mean of the six full-set forecasts
\providecommand{\rhedMeStrat}{}\renewcommand{\rhedMeStrat}{\ensuremath{-}58}  % mean mid strategy return on ME days, %, mean of the six full-set
\providecommand{\rhedMeOtherR}{}\renewcommand{\rhedMeOtherR}{\ensuremath{-}7}  % mean straddle R on the other (covered) days, %
\providecommand{\rhedMeOtherStrat}{}\renewcommand{\rhedMeOtherStrat}{17}  % mean strategy return on the other (covered) days, %
\providecommand{\rhedMeRVIV}{}\renewcommand{\rhedMeRVIV}{1.04}  % mean RV_t / IV_t^2 on ME days
\providecommand{\rhedMeMoveSqIV}{}\renewcommand{\rhedMeMoveSqIV}{1.93}  % mean squared S&P 500 half-hour log return / IV_t^2 on ME days
\providecommand{\rhedMeShare}{}\renewcommand{\rhedMeShare}{\ensuremath{-}24}  % % of the summed mid return earned on ME days, mean of the six full-set forecasts
\providecommand{\rhedRelEndText}{}\renewcommand{\rhedRelEndText}{26~October~2023}  % earliest last flag of the six macro-release series in data/releases.parquet
\providecommand{\rhedRelFomcEndText}{}\renewcommand{\rhedRelFomcEndText}{1~November~2023}  % last FOMC flag in data/releases.parquet
\providecommand{\rhedFomcSuppN}{}\renewcommand{\rhedFomcSuppN}{3}  % FOMC statement days of the sample taken from live/close_signal/state/fomc_statement_dates.csv
\providecommand{\rhedRelAfterN}{}\renewcommand{\rhedRelAfterN}{126}  % post-warmup days after the release calendar ends
\providecommand{\rhedWorthN}{}\renewcommand{\rhedWorthN}{137}  % post-warmup days with R_t = -1 (= rheWorthlessDays)
\providecommand{\rhedWorthPct}{}\renewcommand{\rhedWorthPct}{22.3}  % ... % of the 614 days
\providecommand{\rhedWorthExpIV}{}\renewcommand{\rhedWorthExpIV}{128}  % expected R = -1 days under a normal half-hour log move with volatility IV_t
\providecommand{\rhedWorthExpIVPct}{}\renewcommand{\rhedWorthExpIVPct}{20.9}  % ... as % of the 614 days
\providecommand{\rhedWorthZ}{}\renewcommand{\rhedWorthZ}{0.90}  % (observed - expected) / sd of the count under that normal (independent days)
\providecommand{\rhedWorthExpEmp}{}\renewcommand{\rhedWorthExpEmp}{132}  % expected R = -1 days under the empirical distribution of the earlier days' moves standardized by IV
\providecommand{\rhedWorthExpRaw}{}\renewcommand{\rhedWorthExpRaw}{118}  % expected R = -1 days under the earlier days' raw half-hour moves
\providecommand{\rhedGapFiveN}{}\renewcommand{\rhedGapFiveN}{611}  % post-warmup days with K_C - K_P = 5 points
\providecommand{\rhedGapOtherN}{}\renewcommand{\rhedGapOtherN}{3}  % post-warmup days with another strike gap
\providecommand{\rhedMoveSdPts}{}\renewcommand{\rhedMoveSdPts}{9.6}  % median one-sd implied half-hour move, index points
\providecommand{\rhedStrikeDist}{}\renewcommand{\rhedStrikeDist}{1.23}  % mean distance of the close to the nearest 5-point strike, points
\providecommand{\rhedStrikeDistUnif}{}\renewcommand{\rhedStrikeDistUnif}{1.25}  % ... its mean under a uniform close
\providecommand{\rhedKsD}{}\renewcommand{\rhedKsD}{0.026}  % Kolmogorov-Smirnov D of the close's position within its 5-point strike interval against a uniform
\providecommand{\rhedKsCV}{}\renewcommand{\rhedKsCV}{0.055}  % ... asymptotic 5% critical value
\providecommand{\rhedZSd}{}\renewcommand{\rhedZSd}{1.00}  % sd of the half-hour log move standardized by IV_t, 614 days
\providecommand{\rhedZKurt}{}\renewcommand{\rhedZKurt}{2.0}  % excess kurtosis of that standardized move
\providecommand{\rhedWorthShortMin}{}\renewcommand{\rhedWorthShortMin}{110}  % lowest R = -1 days on which the forecast is short across the six full-set pooled forecasts, 614 days
\providecommand{\rhedWorthShortMax}{}\renewcommand{\rhedWorthShortMax}{115}  % highest R = -1 days on which the forecast is short across the six full-set pooled forecasts, 614 days
\providecommand{\rhedWorthOpexN}{}\renewcommand{\rhedWorthOpexN}{10}  % R = -1 days among the monthly-expiration days
\providecommand{\rhedWorthOpexExp}{}\renewcommand{\rhedWorthOpexExp}{6.1}  % ... expected under the normal with IV_t
\providecommand{\rhedWorthOtherPct}{}\renewcommand{\rhedWorthOtherPct}{21.9}  % % of the other days with R = -1
\providecommand{\rhedWorthOtherExpPct}{}\renewcommand{\rhedWorthOtherExpPct}{21.0}  % ... expected under the normal with IV_t, %
\providecommand{\rhedWorthOpexZ}{}\renewcommand{\rhedWorthOpexZ}{1.77}  % (observed - expected) / sd of the count on the monthly-expiration days
\providecommand{\rhedRegAlphaAMin}{}\renewcommand{\rhedRegAlphaAMin}{11.0}  % lowest alpha, (a) on the short straddle, mid, % per day across the six full-set pooled forecasts, 614 days
\providecommand{\rhedRegAlphaAMax}{}\renewcommand{\rhedRegAlphaAMax}{13.4}  % highest alpha, (a) on the short straddle, mid, % per day across the six full-set pooled forecasts, 614 days
\providecommand{\rhedRegAlphaAXMin}{}\renewcommand{\rhedRegAlphaAXMin}{8.3}  % lowest alpha, (a) on the short straddle, bid/ask, % per day across the six full-set pooled forecasts, 614 days
\providecommand{\rhedRegAlphaAXMax}{}\renewcommand{\rhedRegAlphaAXMax}{10.7}  % highest alpha, (a) on the short straddle, bid/ask, % per day across the six full-set pooled forecasts, 614 days
\providecommand{\rhedRegBetaShortMin}{}\renewcommand{\rhedRegBetaShortMin}{0.23}  % lowest beta on the short-straddle return, (a), mid across the six full-set pooled forecasts, 614 days
\providecommand{\rhedRegBetaShortMax}{}\renewcommand{\rhedRegBetaShortMax}{0.28}  % highest beta on the short-straddle return, (a), mid across the six full-set pooled forecasts, 614 days
\providecommand{\rhedRegAlphaAPMin}{}\renewcommand{\rhedRegAlphaAPMin}{0.005}  % lowest fixed-b p of alpha, (a), mid across the six full-set pooled forecasts, 614 days
\providecommand{\rhedRegAlphaAPMax}{}\renewcommand{\rhedRegAlphaAPMax}{0.025}  % highest fixed-b p of alpha, (a), mid across the six full-set pooled forecasts, 614 days
\providecommand{\rhedRegAlphaAXPMin}{}\renewcommand{\rhedRegAlphaAXPMin}{0.021}  % lowest fixed-b p of alpha, (a), bid/ask across the six full-set pooled forecasts, 614 days
\providecommand{\rhedRegAlphaAXPMax}{}\renewcommand{\rhedRegAlphaAXPMax}{0.077}  % highest fixed-b p of alpha, (a), bid/ask across the six full-set pooled forecasts, 614 days
\providecommand{\rhedRegBetaShortTMin}{}\renewcommand{\rhedRegBetaShortTMin}{1.22}  % lowest t of the beta on the short straddle, (a), mid across the six full-set pooled forecasts, 614 days
\providecommand{\rhedRegBetaShortTMax}{}\renewcommand{\rhedRegBetaShortTMax}{1.42}  % highest t of the beta on the short straddle, (a), mid across the six full-set pooled forecasts, 614 days
\providecommand{\rhedRegRsqAMin}{}\renewcommand{\rhedRegRsqAMin}{0.06}  % lowest R^2 of (a), mid across the six full-set pooled forecasts, 614 days
\providecommand{\rhedRegRsqAMax}{}\renewcommand{\rhedRegRsqAMax}{0.08}  % highest R^2 of (a), mid across the six full-set pooled forecasts, 614 days
\providecommand{\rhedRegBetaAbsMin}{}\renewcommand{\rhedRegBetaAbsMin}{\ensuremath{-}1.87}  % lowest beta on |r30| (bp), (b), mid, % of premium per bp across the six full-set pooled forecasts, 614 days
\providecommand{\rhedRegBetaAbsMax}{}\renewcommand{\rhedRegBetaAbsMax}{\ensuremath{-}1.40}  % highest beta on |r30| (bp), (b), mid, % of premium per bp across the six full-set pooled forecasts, 614 days
\providecommand{\rhedRegBetaAbsTMin}{}\renewcommand{\rhedRegBetaAbsTMin}{\ensuremath{-}3.49}  % lowest t of the beta on |r30|, (b), mid across the six full-set pooled forecasts, 614 days
\providecommand{\rhedRegBetaAbsTMax}{}\renewcommand{\rhedRegBetaAbsTMax}{\ensuremath{-}2.31}  % highest t of the beta on |r30|, (b), mid across the six full-set pooled forecasts, 614 days
\providecommand{\rhedRegBetaMoveMin}{}\renewcommand{\rhedRegBetaMoveMin}{\ensuremath{-}0.01}  % lowest beta on r30 (bp), (b), mid across the six full-set pooled forecasts, 614 days
\providecommand{\rhedRegBetaMoveMax}{}\renewcommand{\rhedRegBetaMoveMax}{0.14}  % highest beta on r30 (bp), (b), mid across the six full-set pooled forecasts, 614 days
\providecommand{\rhedRegMeanMin}{}\renewcommand{\rhedRegMeanMin}{11.5}  % lowest mean return, mid, % per day across the six full-set pooled forecasts, 614 days
\providecommand{\rhedRegMeanMax}{}\renewcommand{\rhedRegMeanMax}{14.1}  % highest mean return, mid, % per day across the six full-set pooled forecasts, 614 days
\providecommand{\rhedRegMeanXMin}{}\renewcommand{\rhedRegMeanXMin}{8.9}  % lowest mean return, bid/ask, % per day across the six full-set pooled forecasts, 614 days
\providecommand{\rhedRegMeanXMax}{}\renewcommand{\rhedRegMeanXMax}{11.4}  % highest mean return, bid/ask, % per day across the six full-set pooled forecasts, 614 days
\providecommand{\rhedRegRsqCMin}{}\renewcommand{\rhedRegRsqCMin}{0.00}  % lowest R^2 of (c), mid across the six full-set pooled forecasts, 614 days
\providecommand{\rhedRegRsqCMax}{}\renewcommand{\rhedRegRsqCMax}{0.00}  % highest R^2 of (c), mid across the six full-set pooled forecasts, 614 days
\providecommand{\rhedRegCPMaxMin}{}\renewcommand{\rhedRegCPMaxMin}{0.74}  % lowest larger fixed-b p of the two betas of (c), mid across the six full-set pooled forecasts, 614 days
\providecommand{\rhedRegCPMaxMax}{}\renewcommand{\rhedRegCPMaxMax}{0.97}  % highest larger fixed-b p of the two betas of (c), mid across the six full-set pooled forecasts, 614 days
\providecommand{\rhedRegCPMinMin}{}\renewcommand{\rhedRegCPMinMin}{0.56}  % lowest smaller fixed-b p of the two betas of (c), mid across the six full-set pooled forecasts, 614 days
\providecommand{\rhedRegCPMinMax}{}\renewcommand{\rhedRegCPMinMax}{0.82}  % highest smaller fixed-b p of the two betas of (c), mid across the six full-set pooled forecasts, 614 days
\providecommand{\rhedRegAlphaDMin}{}\renewcommand{\rhedRegAlphaDMin}{11.4}  % lowest alpha, (d) joint, mid, % per day across the six full-set pooled forecasts, 614 days
\providecommand{\rhedRegAlphaDMax}{}\renewcommand{\rhedRegAlphaDMax}{13.6}  % highest alpha, (d) joint, mid, % per day across the six full-set pooled forecasts, 614 days
\providecommand{\rhedRegAlphaDXMin}{}\renewcommand{\rhedRegAlphaDXMin}{8.7}  % lowest alpha, (d) joint, bid/ask, % per day across the six full-set pooled forecasts, 614 days
\providecommand{\rhedRegAlphaDXMax}{}\renewcommand{\rhedRegAlphaDXMax}{10.9}  % highest alpha, (d) joint, bid/ask, % per day across the six full-set pooled forecasts, 614 days
\providecommand{\rhedRegAlphaDPMin}{}\renewcommand{\rhedRegAlphaDPMin}{0.004}  % lowest fixed-b p of alpha, (d), mid across the six full-set pooled forecasts, 614 days
\providecommand{\rhedRegAlphaDPMax}{}\renewcommand{\rhedRegAlphaDPMax}{0.026}  % highest fixed-b p of alpha, (d), mid across the six full-set pooled forecasts, 614 days
\providecommand{\rhedRegAlphaDXPMin}{}\renewcommand{\rhedRegAlphaDXPMin}{0.017}  % lowest fixed-b p of alpha, (d), bid/ask across the six full-set pooled forecasts, 614 days
\providecommand{\rhedRegAlphaDXPMax}{}\renewcommand{\rhedRegAlphaDXPMax}{0.076}  % highest fixed-b p of alpha, (d), bid/ask across the six full-set pooled forecasts, 614 days
\providecommand{\rhedRegAvgAlphaM}{}\renewcommand{\rhedRegAvgAlphaM}{12.7}  % alpha of the mean of the six full-set returns, (m), mid, % per day
\providecommand{\rhedRegAvgAlphaMT}{}\renewcommand{\rhedRegAvgAlphaMT}{3.48}  % ... its t
\providecommand{\rhedRegAvgAlphaMP}{}\renewcommand{\rhedRegAvgAlphaMP}{0.002}  % ... its fixed-b p
\providecommand{\rhedRegAvgAlphaA}{}\renewcommand{\rhedRegAvgAlphaA}{12.1}  % alpha of the mean of the six full-set returns, (a), mid, % per day
\providecommand{\rhedRegAvgAlphaAT}{}\renewcommand{\rhedRegAvgAlphaAT}{2.83}  % ... its t
\providecommand{\rhedRegAvgAlphaAP}{}\renewcommand{\rhedRegAvgAlphaAP}{0.0100}  % ... its fixed-b p
\providecommand{\rhedRegAvgAlphaD}{}\renewcommand{\rhedRegAvgAlphaD}{12.5}  % alpha of the mean of the six full-set returns, (d), mid, % per day
\providecommand{\rhedRegAvgAlphaDT}{}\renewcommand{\rhedRegAvgAlphaDT}{2.86}  % ... its t
\providecommand{\rhedRegAvgAlphaDP}{}\renewcommand{\rhedRegAvgAlphaDP}{0.009}  % ... its fixed-b p
\providecommand{\rhedRegAvgAlphaMX}{}\renewcommand{\rhedRegAvgAlphaMX}{10.1}  % alpha of the mean of the six full-set returns, (m), ba, % per day
\providecommand{\rhedRegAvgAlphaMXT}{}\renewcommand{\rhedRegAvgAlphaMXT}{2.78}  % ... its t
\providecommand{\rhedRegAvgAlphaMXP}{}\renewcommand{\rhedRegAvgAlphaMXP}{0.011}  % ... its fixed-b p
\providecommand{\rhedRegAvgAlphaAX}{}\renewcommand{\rhedRegAvgAlphaAX}{9.4}  % alpha of the mean of the six full-set returns, (a), ba, % per day
\providecommand{\rhedRegAvgAlphaAXT}{}\renewcommand{\rhedRegAvgAlphaAXT}{2.21}  % ... its t
\providecommand{\rhedRegAvgAlphaAXP}{}\renewcommand{\rhedRegAvgAlphaAXP}{0.040}  % ... its fixed-b p
\providecommand{\rhedRegAvgAlphaDX}{}\renewcommand{\rhedRegAvgAlphaDX}{9.8}  % alpha of the mean of the six full-set returns, (d), ba, % per day
\providecommand{\rhedRegAvgAlphaDXT}{}\renewcommand{\rhedRegAvgAlphaDXT}{2.26}  % ... its t
\providecommand{\rhedRegAvgAlphaDXP}{}\renewcommand{\rhedRegAvgAlphaDXP}{0.036}  % ... its fixed-b p
\providecommand{\rhedRegAvgBetaShort}{}\renewcommand{\rhedRegAvgBetaShort}{0.26}  % beta of the mean full-set return on the short straddle, (a), mid
\providecommand{\rhedRegAvgBetaShortT}{}\renewcommand{\rhedRegAvgBetaShortT}{1.34}  % ... its t
\providecommand{\rhedRegAvgBetaShortP}{}\renewcommand{\rhedRegAvgBetaShortP}{0.20}  % ... its fixed-b p
\providecommand{\rhedRegAvgRsqA}{}\renewcommand{\rhedRegAvgRsqA}{0.07}  % R^2 of (a) for the mean full-set return
\providecommand{\rhedRegAvgBetaAbs}{}\renewcommand{\rhedRegAvgBetaAbs}{\ensuremath{-}1.68}  % beta of the mean full-set return on |r30| (bp), (b), mid
\providecommand{\rhedRegAvgBetaAbsT}{}\renewcommand{\rhedRegAvgBetaAbsT}{\ensuremath{-}3.26}  % ... its t
\providecommand{\rhedRegShortBetaAbs}{}\renewcommand{\rhedRegShortBetaAbs}{\ensuremath{-}3.96}  % beta of the short straddle on |r30| (bp), (b), mid
\providecommand{\rhedRegShortRsqB}{}\renewcommand{\rhedRegShortRsqB}{0.48}  % R^2 of (b) for the short straddle
\providecommand{\rhedRegAbsRatio}{}\renewcommand{\rhedRegAbsRatio}{43}  % beta on |r30| of the mean full-set return as % of the short straddle's
\providecommand{\rhedRegAvgBetaAbsJ}{}\renewcommand{\rhedRegAvgBetaAbsJ}{\ensuremath{-}1.30}  % beta of the mean full-set return on |r30|, (d) joint, mid
\providecommand{\rhedRegAvgBetaAbsJT}{}\renewcommand{\rhedRegAvgBetaAbsJT}{\ensuremath{-}1.90}  % ... its t
\providecommand{\rhedHarRegAlphaAMin}{}\renewcommand{\rhedHarRegAlphaAMin}{1.1}  % lowest alpha, (a), mid, % per day across the six HAR-benchmark pooled forecasts, 614 days
\providecommand{\rhedHarRegAlphaAMax}{}\renewcommand{\rhedHarRegAlphaAMax}{7.0}  % highest alpha, (a), mid, % per day across the six HAR-benchmark pooled forecasts, 614 days
\providecommand{\rhedHarRegAlphaAPMin}{}\renewcommand{\rhedHarRegAlphaAPMin}{0.082}  % lowest fixed-b p of alpha, (a), mid across the six HAR-benchmark pooled forecasts, 614 days
\providecommand{\rhedHarRegAlphaAPMax}{}\renewcommand{\rhedHarRegAlphaAPMax}{0.745}  % highest fixed-b p of alpha, (a), mid across the six HAR-benchmark pooled forecasts, 614 days
\providecommand{\rhedOlsRegAlphaA}{}\renewcommand{\rhedOlsRegAlphaA}{8.1}  % alpha, (a), mid, closing-interval HAR forecast, % per day
\providecommand{\rhedOlsRegAlphaAP}{}\renewcommand{\rhedOlsRegAlphaAP}{0.051}  % ... its fixed-b p
\providecommand{\rhedRegAlphaAExBestMin}{}\renewcommand{\rhedRegAlphaAExBestMin}{\ensuremath{-}0.5}  % lowest alpha, (a), mid, without the forecast's ten best days, % per day across the six full-set pooled forecasts, 614 days
\providecommand{\rhedRegAlphaAExBestMax}{}\renewcommand{\rhedRegAlphaAExBestMax}{2.1}  % highest alpha, (a), mid, without the forecast's ten best days, % per day across the six full-set pooled forecasts, 614 days
\providecommand{\rhedRegAlphaAExBestPMin}{}\renewcommand{\rhedRegAlphaAExBestPMin}{0.43}  % lowest ... its fixed-b p across the six full-set pooled forecasts, 614 days
\providecommand{\rhedRegAlphaAExBestPMax}{}\renewcommand{\rhedRegAlphaAExBestPMax}{0.97}  % highest ... its fixed-b p across the six full-set pooled forecasts, 614 days
\fi
\begin{table}[tbp]
\centering
\small
\setlength{\tabcolsep}{2.5pt}
\begin{threeparttable}
\caption{Long and short days of the sign strategy}
\label{tab:rh_econ_longshort}
\begin{tabular}{@{}lrrrrrrrr@{}}
\toprule
 & \multicolumn{6}{c}{Full predictor set} & \multicolumn{2}{c}{Mean} \\
\cmidrule(lr){2-7}\cmidrule(l){8-9}
 & Ridge & Lasso & \makecell[r]{Elastic\\net} & LightGBM & XGBoost & \makecell[r]{Random\\forest} & Full & HAR \\
\midrule
\multicolumn{9}{@{}l}{\textit{Panel A: Long days ($q_t = +1$)}} \\
\addlinespace[2pt]
Days & $123$ & $121$ & $117$ & $146$ & $143$ & $132$ & $130$ & $167$ \\
Mean return, mid (\%) & $24.7$ & $26.7$ & $30.1$ & $19.2$ & $19.7$ & $27.1$ & $24.6$ & $6.8$ \\
Mean return, bid/ask (\%) & $21.2$ & $23.1$ & $26.4$ & $15.8$ & $16.3$ & $23.6$ & $21.1$ & $3.9$ \\
$\RV_t > \mathrm{IV}_t^2$ (\%) & $48.8$ & $47.1$ & $53.0$ & $51.4$ & $49.7$ & $50.8$ & $50.1$ & $47.0$ \\
Positive return (\%) & $43.1$ & $43.8$ & $45.3$ & $42.5$ & $42.7$ & $45.5$ & $43.8$ & $39.9$ \\
Share of return, mid (\%) & $40.1$ & $40.6$ & $41.3$ & $39.5$ & $39.5$ & $41.4$ & $40.5$ & $29.9$ \\
\addlinespace[4pt]
\multicolumn{9}{@{}l}{\textit{Panel B: Short days ($q_t = -1$)}} \\
\addlinespace[2pt]
Days & $491$ & $493$ & $497$ & $468$ & $471$ & $482$ & $484$ & $447$ \\
Mean return, mid (\%) & $9.2$ & $9.6$ & $10.1$ & $9.2$ & $9.1$ & $10.5$ & $9.6$ & $5.8$ \\
Mean return, bid/ask (\%) & $6.8$ & $7.2$ & $7.7$ & $6.7$ & $6.7$ & $8.1$ & $7.2$ & $3.3$ \\
$\RV_t < \mathrm{IV}_t^2$ (\%) & $77.0$ & $76.5$ & $77.7$ & $79.1$ & $78.3$ & $78.0$ & $77.8$ & $78.8$ \\
Positive return (\%) & $63.5$ & $63.7$ & $64.0$ & $63.7$ & $63.7$ & $64.3$ & $63.8$ & $63.0$ \\
Share of return, mid (\%) & $59.9$ & $59.4$ & $58.7$ & $60.5$ & $60.5$ & $58.6$ & $59.5$ & $70.1$ \\
\addlinespace[4pt]
\multicolumn{9}{@{}l}{\textit{Panel C: All \rheSzDays{} post-warmup days}} \\
\addlinespace[2pt]
Mean return, mid (\%) & $12.3$ & $13.0$ & $13.9$ & $11.5$ & $11.6$ & $14.1$ & $12.7$ & $6.0$ \\
\quad Short straddle, $-\bar R$ (\%) & $2.4$ & $2.4$ & $2.4$ & $2.4$ & $2.4$ & $2.4$ & $2.4$ & $2.4$ \\
\quad Gain, $\bar\Delta$ (\%) & $9.9^{\phantom{***}}$ & $10.5^{*\phantom{**}}$ & $11.5^{*\phantom{**}}$ & $9.1^{\phantom{***}}$ & $9.2^{\phantom{***}}$ & $11.6^{**\phantom{*}}$ & $10.3^{\phantom{***}}$ & $3.6^{\phantom{***}}$ \\
 & $(1.70)^{\phantom{***}}$ & $(1.93)^{\phantom{***}}$ & $(1.95)^{\phantom{***}}$ & $(1.68)^{\phantom{***}}$ & $(1.61)^{\phantom{***}}$ & $(2.14)^{\phantom{***}}$ &  &  \\[2pt]
Hit rate (\%) & $71.3$ & $70.7$ & $73.0$ & $72.5$ & $71.7$ & $72.1$ & $71.9$ & $70.1$ \\
Positive return (\%) & $59.4$ & $59.8$ & $60.4$ & $58.6$ & $58.8$ & $60.3$ & $59.6$ & $56.7$ \\
\addlinespace[4pt]
\multicolumn{9}{@{}l}{\textit{Panel D: Calibration; a bias-corrected forecast (diagnostic)}} \\
\addlinespace[2pt]
Mean $\hat\sigma^2_t$ / mean $\RV_t$ & $0.97$ & $0.95$ & $0.96$ & $1.01$ & $1.01$ & $1.01$ & $0.98$ & $1.00$ \\
Mean of $\RV_t/\hat\sigma^2_t$ & $1.19$ & $1.20$ & $1.19$ & $1.09$ & $1.10$ & $1.11$ & $1.15$ & $1.07$ \\
Long days, corrected & $228$ & $241$ & $225$ & $201$ & $204$ & $222$ & $220$ & $212$ \\
Positions changed & $105$ & $120$ & $108$ & $55$ & $61$ & $90$ & $90$ & $45$ \\
$\bar R$ on changed days (\%) & $-10.4$ & $-7.9$ & $-18.3$ & $-9.0$ & $-14.8$ & $-24.4$ & $-14.1$ & $6.1$ \\
$\SR$, mid: uncorrected & $1.75$ & $1.75$ & $1.95$ & $1.67$ & $1.71$ & $1.99$ & $1.80$ & $0.75$ \\
$\SR$, mid: corrected & $1.08$ & $1.19$ & $0.88$ & $1.42$ & $1.24$ & $0.99$ & $1.13$ & $1.05$ \\
$\SR$, bid/ask: uncorrected & $1.40$ & $1.40$ & $1.59$ & $1.31$ & $1.36$ & $1.63$ & $1.45$ & $0.39$ \\
$\SR$, bid/ask: corrected & $0.73$ & $0.85$ & $0.53$ & $1.08$ & $0.89$ & $0.65$ & $0.79$ & $0.69$ \\
\bottomrule
\end{tabular}
\begin{tablenotes}[flushleft]
\footnotesize
\item \textit{Notes:} Sign strategy $q_t = \sign(\hat\sigma^2_t - \mathrm{IV}_t^2)$ of each pooled intraday forecast
on the full predictor set, on the \rheSzDays{} post-warmup option days (Section~\ref{sec:rh_econ_strategy}); returns
in \% of the premium. Mean: the mean over the six full-set forecasts and over their six HAR-benchmark counterparts (the
three penalized linear models on the HAR benchmark take the same position on every day); in the share rows, the summed
returns of the six. $\RV_t > \mathrm{IV}_t^2$ and $\RV_t < \mathrm{IV}_t^2$: the share of the long (short) days on
which realized variance exceeds (falls short of) implied variance, which happens on \rhedUncRVAboveIV\% of all days;
Positive return: the share of days with $r_t > 0$ at mid-quotes; Share of return: the part of the summed mid-quote
return earned on those days; Hit rate: the share of days on which $q_t$ has the sign of $\RV_t - \mathrm{IV}_t^2$
(\rhedUncShortHit\% for the unconditional short straddle). Panel C: because
$r_t = -R_t + 2R_t\mathbf{1}\{q_t = +1\}$, the mean return is that of the unconditional short straddle plus the gain
$\bar\Delta$ over it, earned on the long days only; in parentheses, the $t$-statistic of $\bar\Delta$ (Newey--West,
bandwidth $S = \rheHACS$, fixed-$b$ critical values; Section~\ref{sec:rh_methods_tests}). Panel D: the corrected
forecast is $c_t\hat\sigma^2_t$, with $c_t$ the mean of $\RV_{t'}/\hat\sigma^2_{t'}$ over the earlier option days
$t' < t$; a diagnostic, not a strategy of the paper. $\SR$ is the paper's Sharpe ratio
(Section~\ref{sec:rh_econ_measures}).
$^{*}$, $^{**}$ and $^{***}$ denote significance at the 10\%, 5\% and 1\% levels.
\end{tablenotes}
\end{threeparttable}
\end{table}
    % tab:rh_econ_longshort
%   % AUTO-GENERATED by writeup/make_rehaul_econ_decomp.py -- do not edit; re-run the script.
% Table tab:rh_econ_quintiles: the sign strategy by quintile of the forecast variance gap
% (writeup/make_rehaul_econ_decomp.py, statistic 2).
\ifdefined\rhedNDays\else% AUTO-GENERATED by writeup/make_rehaul_econ_decomp.py -- do not edit; re-run the script.
% \rhed... macros: where the straddle strategy's profits come from (Section 7 decomposition).
\providecommand{\rhedNDays}{}\renewcommand{\rhedNDays}{614}  % post-warmup option days (= rheSzDays)
\providecommand{\rhedNFull}{}\renewcommand{\rhedNFull}{6}  % pooled intraday forecasts on the full predictor set
\providecommand{\rhedNTop}{}\renewcommand{\rhedNTop}{10}  % best / worst days per forecast
\providecommand{\rhedUncRVAboveIV}{}\renewcommand{\rhedUncRVAboveIV}{28.2}  % % of the 614 days with RV_t > IV_t^2
\providecommand{\rhedUncRVIVMean}{}\renewcommand{\rhedUncRVIVMean}{0.87}  % mean RV_t / IV_t^2, 614 days
\providecommand{\rhedUncRVIVMed}{}\renewcommand{\rhedUncRVIVMed}{0.77}  % median RV_t / IV_t^2, 614 days
\providecommand{\rhedUncShortMean}{}\renewcommand{\rhedUncShortMean}{2.4}  % mean mid return of the unconditional short straddle, % of premium
\providecommand{\rhedUncShortHit}{}\renewcommand{\rhedUncShortHit}{71.8}  % % of days on which q = -1 matches sign(RV - IV^2) (always short)
\providecommand{\rhedLongNMin}{}\renewcommand{\rhedLongNMin}{117}  % lowest number of long days across the six full-set pooled forecasts, 614 days
\providecommand{\rhedLongNMax}{}\renewcommand{\rhedLongNMax}{146}  % highest number of long days across the six full-set pooled forecasts, 614 days
\providecommand{\rhedLongPctMin}{}\renewcommand{\rhedLongPctMin}{19}  % lowest % long days across the six full-set pooled forecasts, 614 days
\providecommand{\rhedLongPctMax}{}\renewcommand{\rhedLongPctMax}{24}  % highest % long days across the six full-set pooled forecasts, 614 days
\providecommand{\rhedHarLongNMin}{}\renewcommand{\rhedHarLongNMin}{150}  % lowest number of long days across the six HAR-benchmark pooled forecasts, 614 days
\providecommand{\rhedHarLongNMax}{}\renewcommand{\rhedHarLongNMax}{193}  % highest number of long days across the six HAR-benchmark pooled forecasts, 614 days
\providecommand{\rhedLongMidMin}{}\renewcommand{\rhedLongMidMin}{19}  % lowest mean mid return on long days, % across the six full-set pooled forecasts, 614 days
\providecommand{\rhedLongMidMax}{}\renewcommand{\rhedLongMidMax}{30}  % highest mean mid return on long days, % across the six full-set pooled forecasts, 614 days
\providecommand{\rhedLongBaMin}{}\renewcommand{\rhedLongBaMin}{16}  % lowest mean bid/ask return on long days, % across the six full-set pooled forecasts, 614 days
\providecommand{\rhedLongBaMax}{}\renewcommand{\rhedLongBaMax}{26}  % highest mean bid/ask return on long days, % across the six full-set pooled forecasts, 614 days
\providecommand{\rhedShortMidMin}{}\renewcommand{\rhedShortMidMin}{9}  % lowest mean mid return on short days, % across the six full-set pooled forecasts, 614 days
\providecommand{\rhedShortMidMax}{}\renewcommand{\rhedShortMidMax}{11}  % highest mean mid return on short days, % across the six full-set pooled forecasts, 614 days
\providecommand{\rhedShortBaMin}{}\renewcommand{\rhedShortBaMin}{7}  % lowest mean bid/ask return on short days, % across the six full-set pooled forecasts, 614 days
\providecommand{\rhedShortBaMax}{}\renewcommand{\rhedShortBaMax}{8}  % highest mean bid/ask return on short days, % across the six full-set pooled forecasts, 614 days
\providecommand{\rhedHarLongMidMin}{}\renewcommand{\rhedHarLongMidMin}{0}  % lowest mean mid return on long days, % across the six HAR-benchmark pooled forecasts, 614 days
\providecommand{\rhedHarLongMidMax}{}\renewcommand{\rhedHarLongMidMax}{10}  % highest mean mid return on long days, % across the six HAR-benchmark pooled forecasts, 614 days
\providecommand{\rhedHarShortMidMin}{}\renewcommand{\rhedHarShortMidMin}{3}  % lowest mean mid return on short days, % across the six HAR-benchmark pooled forecasts, 614 days
\providecommand{\rhedHarShortMidMax}{}\renewcommand{\rhedHarShortMidMax}{7}  % highest mean mid return on short days, % across the six HAR-benchmark pooled forecasts, 614 days
\providecommand{\rhedLongHitMin}{}\renewcommand{\rhedLongHitMin}{47}  % lowest % of long days with RV > IV^2 across the six full-set pooled forecasts, 614 days
\providecommand{\rhedLongHitMax}{}\renewcommand{\rhedLongHitMax}{53}  % highest % of long days with RV > IV^2 across the six full-set pooled forecasts, 614 days
\providecommand{\rhedHarLongHitMin}{}\renewcommand{\rhedHarLongHitMin}{46}  % lowest % of long days with RV > IV^2 across the six HAR-benchmark pooled forecasts, 614 days
\providecommand{\rhedHarLongHitMax}{}\renewcommand{\rhedHarLongHitMax}{48}  % highest % of long days with RV > IV^2 across the six HAR-benchmark pooled forecasts, 614 days
\providecommand{\rhedLongPosMin}{}\renewcommand{\rhedLongPosMin}{42}  % lowest % of long days with a positive return across the six full-set pooled forecasts, 614 days
\providecommand{\rhedLongPosMax}{}\renewcommand{\rhedLongPosMax}{45}  % highest % of long days with a positive return across the six full-set pooled forecasts, 614 days
\providecommand{\rhedHitMin}{}\renewcommand{\rhedHitMin}{70.7}  % lowest % of days on which q matches sign(RV - IV^2) across the six full-set pooled forecasts, 614 days
\providecommand{\rhedHitMax}{}\renewcommand{\rhedHitMax}{73.0}  % highest % of days on which q matches sign(RV - IV^2) across the six full-set pooled forecasts, 614 days
\providecommand{\rhedPosMin}{}\renewcommand{\rhedPosMin}{58.6}  % lowest % of days with a positive mid return across the six full-set pooled forecasts, 614 days
\providecommand{\rhedPosMax}{}\renewcommand{\rhedPosMax}{60.4}  % highest % of days with a positive mid return across the six full-set pooled forecasts, 614 days
\providecommand{\rhedLongShareMin}{}\renewcommand{\rhedLongShareMin}{39}  % lowest % of the summed mid return earned on long days across the six full-set pooled forecasts, 614 days
\providecommand{\rhedLongShareMax}{}\renewcommand{\rhedLongShareMax}{41}  % highest % of the summed mid return earned on long days across the six full-set pooled forecasts, 614 days
\providecommand{\rhedLongShareXMin}{}\renewcommand{\rhedLongShareXMin}{42}  % lowest % of the summed bid/ask return earned on long days across the six full-set pooled forecasts, 614 days
\providecommand{\rhedLongShareXMax}{}\renewcommand{\rhedLongShareXMax}{45}  % highest % of the summed bid/ask return earned on long days across the six full-set pooled forecasts, 614 days
\providecommand{\rhedMeanMin}{}\renewcommand{\rhedMeanMin}{11.5}  % lowest mean mid return, % across the six full-set pooled forecasts, 614 days
\providecommand{\rhedMeanMax}{}\renewcommand{\rhedMeanMax}{14.1}  % highest mean mid return, % across the six full-set pooled forecasts, 614 days
\providecommand{\rhedGainMin}{}\renewcommand{\rhedGainMin}{9.1}  % lowest mean gain over the short straddle (2 R 1(long)), % across the six full-set pooled forecasts, 614 days
\providecommand{\rhedGainMax}{}\renewcommand{\rhedGainMax}{11.6}  % highest mean gain over the short straddle (2 R 1(long)), % across the six full-set pooled forecasts, 614 days
\providecommand{\rhedGainShareMin}{}\renewcommand{\rhedGainShareMin}{79}  % lowest gain over short / mean return, % across the six full-set pooled forecasts, 614 days
\providecommand{\rhedGainShareMax}{}\renewcommand{\rhedGainShareMax}{83}  % highest gain over short / mean return, % across the six full-set pooled forecasts, 614 days
\providecommand{\rhedHarGainMin}{}\renewcommand{\rhedHarGainMin}{0.0}  % lowest mean gain over the short straddle, % across the six HAR-benchmark pooled forecasts, 614 days
\providecommand{\rhedHarGainMax}{}\renewcommand{\rhedHarGainMax}{5.0}  % highest mean gain over the short straddle, % across the six HAR-benchmark pooled forecasts, 614 days
\providecommand{\rhedOlsLongN}{}\renewcommand{\rhedOlsLongN}{198}  % long days of the closing-interval HAR forecast (sub_ols_baseline)
\providecommand{\rhedOlsLongMid}{}\renewcommand{\rhedOlsLongMid}{9}  % ... its mean mid return on long days, %
\providecommand{\rhedOlsShortMid}{}\renewcommand{\rhedOlsShortMid}{8}  % ... its mean mid return on short days, %
\providecommand{\rhedOlsLongHit}{}\renewcommand{\rhedOlsLongHit}{46}  % ... % of its long days with RV > IV^2
\providecommand{\rhedOlsSh}{}\renewcommand{\rhedOlsSh}{1.11}  % ... its Sharpe ratio, mid (= rhrOlsSharpe)
\providecommand{\rhedCalLevelMin}{}\renewcommand{\rhedCalLevelMin}{0.95}  % lowest mean forecast / mean RV across the six full-set pooled forecasts, 614 days
\providecommand{\rhedCalLevelMax}{}\renewcommand{\rhedCalLevelMax}{1.01}  % highest mean forecast / mean RV across the six full-set pooled forecasts, 614 days
\providecommand{\rhedCalScaleMin}{}\renewcommand{\rhedCalScaleMin}{1.09}  % lowest mean of RV / forecast (QLIKE-optimal scale) across the six full-set pooled forecasts, 614 days
\providecommand{\rhedCalScaleMax}{}\renewcommand{\rhedCalScaleMax}{1.20}  % highest mean of RV / forecast (QLIKE-optimal scale) across the six full-set pooled forecasts, 614 days
\providecommand{\rhedHarCalScaleMin}{}\renewcommand{\rhedHarCalScaleMin}{1.04}  % lowest mean of RV / forecast across the six HAR-benchmark pooled forecasts, 614 days
\providecommand{\rhedHarCalScaleMax}{}\renewcommand{\rhedHarCalScaleMax}{1.10}  % highest mean of RV / forecast across the six HAR-benchmark pooled forecasts, 614 days
\providecommand{\rhedMedFcIVMin}{}\renewcommand{\rhedMedFcIVMin}{0.75}  % lowest median forecast / IV^2 across the six full-set pooled forecasts, 614 days
\providecommand{\rhedMedFcIVMax}{}\renewcommand{\rhedMedFcIVMax}{0.82}  % highest median forecast / IV^2 across the six full-set pooled forecasts, 614 days
\providecommand{\rhedBcScaleLoMin}{}\renewcommand{\rhedBcScaleLoMin}{1.06}  % lowest lowest causal scale c_t over the 614 days across the six full-set pooled forecasts, 614 days
\providecommand{\rhedBcScaleLoMax}{}\renewcommand{\rhedBcScaleLoMax}{1.15}  % highest lowest causal scale c_t over the 614 days across the six full-set pooled forecasts, 614 days
\providecommand{\rhedBcScaleHiMin}{}\renewcommand{\rhedBcScaleHiMin}{1.10}  % lowest highest causal scale c_t over the 614 days across the six full-set pooled forecasts, 614 days
\providecommand{\rhedBcScaleHiMax}{}\renewcommand{\rhedBcScaleHiMax}{1.22}  % highest highest causal scale c_t over the 614 days across the six full-set pooled forecasts, 614 days
\providecommand{\rhedBcFlipMin}{}\renewcommand{\rhedBcFlipMin}{55}  % lowest positions changed by the bias correction across the six full-set pooled forecasts, 614 days
\providecommand{\rhedBcFlipMax}{}\renewcommand{\rhedBcFlipMax}{120}  % highest positions changed by the bias correction across the six full-set pooled forecasts, 614 days
\providecommand{\rhedBcLongMin}{}\renewcommand{\rhedBcLongMin}{201}  % lowest long days under the bias correction across the six full-set pooled forecasts, 614 days
\providecommand{\rhedBcLongMax}{}\renewcommand{\rhedBcLongMax}{241}  % highest long days under the bias correction across the six full-set pooled forecasts, 614 days
\providecommand{\rhedBcFlipRMin}{}\renewcommand{\rhedBcFlipRMin}{\ensuremath{-}24}  % lowest mean mid straddle return R on the changed days, % across the six full-set pooled forecasts, 614 days
\providecommand{\rhedBcFlipRMax}{}\renewcommand{\rhedBcFlipRMax}{\ensuremath{-}8}  % highest mean mid straddle return R on the changed days, % across the six full-set pooled forecasts, 614 days
\providecommand{\rhedBcShMin}{}\renewcommand{\rhedBcShMin}{0.92}  % lowest Sharpe ratio mid under the bias correction (diagnostic) across the six full-set pooled forecasts, 614 days
\providecommand{\rhedBcShMax}{}\renewcommand{\rhedBcShMax}{1.32}  % highest Sharpe ratio mid under the bias correction (diagnostic) across the six full-set pooled forecasts, 614 days
\providecommand{\rhedBcShXMin}{}\renewcommand{\rhedBcShXMin}{0.57}  % lowest Sharpe ratio bid/ask under the bias correction across the six full-set pooled forecasts, 614 days
\providecommand{\rhedBcShXMax}{}\renewcommand{\rhedBcShXMax}{0.97}  % highest Sharpe ratio bid/ask under the bias correction across the six full-set pooled forecasts, 614 days
\providecommand{\rhedLevFlipMin}{}\renewcommand{\rhedLevFlipMin}{4}  % lowest positions changed by the level correction (sum RV / sum fc) across the six full-set pooled forecasts, 614 days
\providecommand{\rhedLevFlipMax}{}\renewcommand{\rhedLevFlipMax}{19}  % highest positions changed by the level correction (sum RV / sum fc) across the six full-set pooled forecasts, 614 days
\providecommand{\rhedLevShMin}{}\renewcommand{\rhedLevShMin}{1.61}  % lowest Sharpe ratio mid under the level correction across the six full-set pooled forecasts, 614 days
\providecommand{\rhedLevShMax}{}\renewcommand{\rhedLevShMax}{2.01}  % highest Sharpe ratio mid under the level correction across the six full-set pooled forecasts, 614 days
\providecommand{\rhedHarBcShMin}{}\renewcommand{\rhedHarBcShMin}{0.76}  % lowest Sharpe ratio mid under the bias correction across the six HAR-benchmark pooled forecasts, 614 days
\providecommand{\rhedHarBcShMax}{}\renewcommand{\rhedHarBcShMax}{1.05}  % highest Sharpe ratio mid under the bias correction across the six HAR-benchmark pooled forecasts, 614 days
\providecommand{\rhedBcShUpN}{}\renewcommand{\rhedBcShUpN}{0}  % full-set forecasts whose mid SR rises under the bias correction
\providecommand{\rhedBcShDownN}{}\renewcommand{\rhedBcShDownN}{6}  % full-set forecasts whose mid SR falls under the bias correction
\providecommand{\rhedHarBcShUpN}{}\renewcommand{\rhedHarBcShUpN}{2}  % HAR-benchmark forecasts whose mid SR rises under the bias correction
\providecommand{\rhedMoveSqIV}{}\renewcommand{\rhedMoveSqIV}{1.01}  % mean of the squared S&P 500 half-hour log return over IV_t^2, 614 days
\providecommand{\rhedMoveSqRV}{}\renewcommand{\rhedMoveSqRV}{1.15}  % mean squared S&P 500 half-hour log return over mean RV_t, 614 days
\providecommand{\rhedOlsCalScale}{}\renewcommand{\rhedOlsCalScale}{1.01}  % mean of RV / forecast, closing-interval HAR forecast
\providecommand{\rhedOlsBcSh}{}\renewcommand{\rhedOlsBcSh}{1.12}  % Sharpe ratio mid under the bias correction, closing-interval HAR
\providecommand{\rhedBcFlipRidge}{}\renewcommand{\rhedBcFlipRidge}{105}  % positions changed by the bias correction, poolpt13_ridge_all_features
\providecommand{\rhedBcShRidge}{}\renewcommand{\rhedBcShRidge}{1.17}  % Sharpe ratio mid under the bias correction, poolpt13_ridge_all_features
\providecommand{\rhedBcFlipLasso}{}\renewcommand{\rhedBcFlipLasso}{120}  % positions changed by the bias correction, poolpt13_lasso_all_features
\providecommand{\rhedBcShLasso}{}\renewcommand{\rhedBcShLasso}{1.31}  % Sharpe ratio mid under the bias correction, poolpt13_lasso_all_features
\providecommand{\rhedBcFlipEnet}{}\renewcommand{\rhedBcFlipEnet}{108}  % positions changed by the bias correction, poolpt13_enet_all_features
\providecommand{\rhedBcShEnet}{}\renewcommand{\rhedBcShEnet}{0.99}  % Sharpe ratio mid under the bias correction, poolpt13_enet_all_features
\providecommand{\rhedBcFlipLgbm}{}\renewcommand{\rhedBcFlipLgbm}{55}  % positions changed by the bias correction, pooltree_lgbm_all_features
\providecommand{\rhedBcShLgbm}{}\renewcommand{\rhedBcShLgbm}{1.32}  % Sharpe ratio mid under the bias correction, pooltree_lgbm_all_features
\providecommand{\rhedBcFlipXgb}{}\renewcommand{\rhedBcFlipXgb}{61}  % positions changed by the bias correction, pooltree_xgb_all_features
\providecommand{\rhedBcShXgb}{}\renewcommand{\rhedBcShXgb}{1.15}  % Sharpe ratio mid under the bias correction, pooltree_xgb_all_features
\providecommand{\rhedBcFlipRf}{}\renewcommand{\rhedBcFlipRf}{90}  % positions changed by the bias correction, pooltree_rf_all_features
\providecommand{\rhedBcShRf}{}\renewcommand{\rhedBcShRf}{0.92}  % Sharpe ratio mid under the bias correction, pooltree_rf_all_features
\providecommand{\rhedQMonoN}{}\renewcommand{\rhedQMonoN}{6}  % full-set forecasts whose mean RV/IV^2 rises across all five quintiles
\providecommand{\rhedHarQMonoN}{}\renewcommand{\rhedHarQMonoN}{4}  % HAR-benchmark forecasts whose mean RV/IV^2 rises across all five quintiles
\providecommand{\rhedQRPatternN}{}\renewcommand{\rhedQRPatternN}{6}  % full-set forecasts with mean R < 0 in quintiles 1-4 and > 0 in 5
\providecommand{\rhedHarQRPatternN}{}\renewcommand{\rhedHarQRPatternN}{0}  % HAR-benchmark forecasts with that pattern
\providecommand{\rhedQOneRVIV}{}\renewcommand{\rhedQOneRVIV}{0.73}  % mean RV/IV^2 in the lowest quintile, mean of the six full-set
\providecommand{\rhedQFiveRVIV}{}\renewcommand{\rhedQFiveRVIV}{1.05}  % mean RV/IV^2 in the highest quintile, mean of the six full-set
\providecommand{\rhedQFiveR}{}\renewcommand{\rhedQFiveR}{26}  % mean straddle return R in the highest quintile, %, mean of the six full-set
\providecommand{\rhedQFiveFc}{}\renewcommand{\rhedQFiveFc}{1.18}  % geometric-mean forecast / IV^2 in the highest quintile, mean of the six full-set
\providecommand{\rhedQFourFc}{}\renewcommand{\rhedQFourFc}{0.93}  % ... in the fourth quintile
\providecommand{\rhedQFiveLong}{}\renewcommand{\rhedQFiveLong}{99}  % % long days in the highest quintile, mean of the six full-set
\providecommand{\rhedQFiveHit}{}\renewcommand{\rhedQFiveHit}{52}  % hit rate in the highest quintile, %, mean of the six full-set
\providecommand{\rhedQOneHit}{}\renewcommand{\rhedQOneHit}{88}  % hit rate in the lowest quintile, %, mean of the six full-set
\providecommand{\rhedQFiveShare}{}\renewcommand{\rhedQFiveShare}{43}  % % of the summed mid return earned in the highest quintile, six full-set forecasts
\providecommand{\rhedQLowRMin}{}\renewcommand{\rhedQLowRMin}{\ensuremath{-}13}  % most negative mean straddle return R in quintiles 1-4, %, mean of the six
\providecommand{\rhedQLowRMax}{}\renewcommand{\rhedQLowRMax}{\ensuremath{-}6}  % least negative mean straddle return R in quintiles 1-4, %, mean of the six
\providecommand{\rhedHarQOneRVIV}{}\renewcommand{\rhedHarQOneRVIV}{0.74}  % mean RV/IV^2 in the lowest quintile, mean of the six Har
\providecommand{\rhedHarQFiveRVIV}{}\renewcommand{\rhedHarQFiveRVIV}{1.02}  % mean RV/IV^2 in the highest quintile, mean of the six Har
\providecommand{\rhedHarQFiveR}{}\renewcommand{\rhedHarQFiveR}{7}  % mean straddle return R in the highest quintile, %, mean of the six Har
\providecommand{\rhedHarQFiveFc}{}\renewcommand{\rhedHarQFiveFc}{1.23}  % geometric-mean forecast / IV^2 in the highest quintile, mean of the six Har
\providecommand{\rhedHarQFourFc}{}\renewcommand{\rhedHarQFourFc}{0.98}  % ... in the fourth quintile
\providecommand{\rhedHarQFiveLong}{}\renewcommand{\rhedHarQFiveLong}{100}  % % long days in the highest quintile, mean of the six Har
\providecommand{\rhedHarQFiveHit}{}\renewcommand{\rhedHarQFiveHit}{47}  % hit rate in the highest quintile, %, mean of the six Har
\providecommand{\rhedHarQOneHit}{}\renewcommand{\rhedHarQOneHit}{89}  % hit rate in the lowest quintile, %, mean of the six Har
\providecommand{\rhedHarQFiveShare}{}\renewcommand{\rhedHarQFiveShare}{23}  % % of the summed mid return earned in the highest quintile, six Har forecasts
\providecommand{\rhedHarQLowRMin}{}\renewcommand{\rhedHarQLowRMin}{\ensuremath{-}13}  % most negative mean straddle return R in quintiles 1-4, %, mean of the six
\providecommand{\rhedHarQLowRMax}{}\renewcommand{\rhedHarQLowRMax}{5}  % least negative mean straddle return R in quintiles 1-4, %, mean of the six
\providecommand{\rhedQFiveShareRngMin}{}\renewcommand{\rhedQFiveShareRngMin}{40}  % lowest % of the summed mid return in quintile 5 across the six full-set pooled forecasts, 614 days
\providecommand{\rhedQFiveShareRngMax}{}\renewcommand{\rhedQFiveShareRngMax}{48}  % highest % of the summed mid return in quintile 5 across the six full-set pooled forecasts, 614 days
\providecommand{\rhedBestUnion}{}\renewcommand{\rhedBestUnion}{11}  % days in the union of the six full-set forecasts' ten best days
\providecommand{\rhedWorstUnion}{}\renewcommand{\rhedWorstUnion}{11}  % days in the union of the six full-set forecasts' ten worst days
\providecommand{\rhedBestAllSix}{}\renewcommand{\rhedBestAllSix}{9}  % days among the ten best of all six forecasts
\providecommand{\rhedWorstAllSix}{}\renewcommand{\rhedWorstAllSix}{9}  % days among the ten worst of all six forecasts
\providecommand{\rhedTopDate}{}\renewcommand{\rhedTopDate}{22~November~2021}  % the day of the largest straddle return among the best days
\providecommand{\rhedTopR}{}\renewcommand{\rhedTopR}{1{,}032}  % ... the straddle's mid return that day, % of premium
\providecommand{\rhedTopRVIV}{}\renewcommand{\rhedTopRVIV}{2.79}  % ... RV / IV^2 that day
\providecommand{\rhedTopMove}{}\renewcommand{\rhedTopMove}{\ensuremath{-}0.83}  % ... S&P 500 15:30-16:00 log return that day, %
\providecommand{\rhedTopShareMin}{}\renewcommand{\rhedTopShareMin}{12}  % lowest % of a full-set forecast's summed mid return earned on that one day
\providecommand{\rhedTopShareMax}{}\renewcommand{\rhedTopShareMax}{15}  % highest % of a full-set forecast's summed mid return earned on that one day
\providecommand{\rhedBestFomcN}{}\renewcommand{\rhedBestFomcN}{1}  % best-day union: FOMC statement days
\providecommand{\rhedWorstFomcN}{}\renewcommand{\rhedWorstFomcN}{1}  % worst-day union: FOMC statement days
\providecommand{\rhedWorstMeN}{}\renewcommand{\rhedWorstMeN}{3}  % worst-day union: month ends
\providecommand{\rhedBestMeN}{}\renewcommand{\rhedBestMeN}{0}  % best-day union: month ends
\providecommand{\rhedBestOpexN}{}\renewcommand{\rhedBestOpexN}{0}  % best-day union: monthly expirations
\providecommand{\rhedWorstOpexN}{}\renewcommand{\rhedWorstOpexN}{0}  % worst-day union: monthly expirations
\providecommand{\rhedBestNoTagN}{}\renewcommand{\rhedBestNoTagN}{8}  % best-day union: days without any event tag (calendar as covered)
\providecommand{\rhedBestAfterRelN}{}\renewcommand{\rhedBestAfterRelN}{5}  % best-day union: days after the release calendar ends (dagger)
\providecommand{\rhedWorstAfterRelN}{}\renewcommand{\rhedWorstAfterRelN}{4}  % worst-day union: days after the release calendar ends
\providecommand{\rhedDShortAllMin}{}\renewcommand{\rhedDShortAllMin}{1.21}  % lowest Delta SR mid vs the short straddle, all days across the six full-set pooled forecasts, 614 days
\providecommand{\rhedDShortAllMax}{}\renewcommand{\rhedDShortAllMax}{1.56}  % highest Delta SR mid vs the short straddle, all days across the six full-set pooled forecasts, 614 days
\providecommand{\rhedDShortExBestMin}{}\renewcommand{\rhedDShortExBestMin}{\ensuremath{-}0.54}  % lowest Delta SR mid vs the short straddle without the forecast's ten best days across the six full-set pooled forecasts, 614 days
\providecommand{\rhedDShortExBestMax}{}\renewcommand{\rhedDShortExBestMax}{\ensuremath{-}0.12}  % highest Delta SR mid vs the short straddle without the forecast's ten best days across the six full-set pooled forecasts, 614 days
\providecommand{\rhedDShortExBothMin}{}\renewcommand{\rhedDShortExBothMin}{\ensuremath{-}0.64}  % lowest ... without its ten best and ten worst days across the six full-set pooled forecasts, 614 days
\providecommand{\rhedDShortExBothMax}{}\renewcommand{\rhedDShortExBothMax}{\ensuremath{-}0.14}  % highest ... without its ten best and ten worst days across the six full-set pooled forecasts, 614 days
\providecommand{\rhedBestShareMin}{}\renewcommand{\rhedBestShareMin}{107}  % lowest % of the summed gain over the short straddle earned on the ten best days across the six full-set pooled forecasts, 614 days
\providecommand{\rhedBestShareMax}{}\renewcommand{\rhedBestShareMax}{138}  % highest % of the summed gain over the short straddle earned on the ten best days across the six full-set pooled forecasts, 614 days
\providecommand{\rhedDHarAllMin}{}\renewcommand{\rhedDHarAllMin}{0.56}  % lowest Delta SR mid vs the HAR twin, all days across the six full-set pooled forecasts, 614 days
\providecommand{\rhedDHarAllMax}{}\renewcommand{\rhedDHarAllMax}{1.56}  % highest Delta SR mid vs the HAR twin, all days across the six full-set pooled forecasts, 614 days
\providecommand{\rhedDHarExBestMin}{}\renewcommand{\rhedDHarExBestMin}{0.30}  % lowest Delta SR mid vs the HAR twin without the forecast's ten best days across the six full-set pooled forecasts, 614 days
\providecommand{\rhedDHarExBestMax}{}\renewcommand{\rhedDHarExBestMax}{0.83}  % highest Delta SR mid vs the HAR twin without the forecast's ten best days across the six full-set pooled forecasts, 614 days
\providecommand{\rhedDHarExBothMin}{}\renewcommand{\rhedDHarExBothMin}{0.36}  % lowest ... without its ten best and ten worst days across the six full-set pooled forecasts, 614 days
\providecommand{\rhedDHarExBothMax}{}\renewcommand{\rhedDHarExBothMax}{0.97}  % highest ... without its ten best and ten worst days across the six full-set pooled forecasts, 614 days
\providecommand{\rhedHarBestShareMin}{}\renewcommand{\rhedHarBestShareMin}{47}  % lowest % of the summed gain over the HAR twin earned on the ten best days across the six full-set pooled forecasts, 614 days
\providecommand{\rhedHarBestShareMax}{}\renewcommand{\rhedHarBestShareMax}{61}  % highest % of the summed gain over the HAR twin earned on the ten best days across the six full-set pooled forecasts, 614 days
\providecommand{\rhedDShortExBestNeg}{}\renewcommand{\rhedDShortExBestNeg}{6}  % full-set forecasts whose Delta SR vs short is negative without the ten best days
\providecommand{\rhedDHarExBestNeg}{}\renewcommand{\rhedDHarExBestNeg}{0}  % full-set forecasts whose Delta SR vs the HAR twin is negative without the ten best days
\providecommand{\rhedFomcN}{}\renewcommand{\rhedFomcN}{21}  % FOMC days among the 614 post-warmup days
\providecommand{\rhedFomcR}{}\renewcommand{\rhedFomcR}{25}  % mean mid straddle return R on FOMC days, %
\providecommand{\rhedFomcRT}{}\renewcommand{\rhedFomcRT}{1.02}  % ... its t (i.i.d.; the days are weeks apart)
\providecommand{\rhedFomcLong}{}\renewcommand{\rhedFomcLong}{25}  % % of FOMC days long, mean of the six full-set forecasts
\providecommand{\rhedFomcStrat}{}\renewcommand{\rhedFomcStrat}{\ensuremath{-}6}  % mean mid strategy return on FOMC days, %, mean of the six full-set
\providecommand{\rhedFomcOtherR}{}\renewcommand{\rhedFomcOtherR}{\ensuremath{-}3}  % mean straddle R on the other (covered) days, %
\providecommand{\rhedFomcOtherStrat}{}\renewcommand{\rhedFomcOtherStrat}{13}  % mean strategy return on the other (covered) days, %
\providecommand{\rhedFomcRVIV}{}\renewcommand{\rhedFomcRVIV}{0.76}  % mean RV_t / IV_t^2 on FOMC days
\providecommand{\rhedFomcMoveSqIV}{}\renewcommand{\rhedFomcMoveSqIV}{1.43}  % mean squared S&P 500 half-hour log return / IV_t^2 on FOMC days
\providecommand{\rhedFomcShare}{}\renewcommand{\rhedFomcShare}{\ensuremath{-}2}  % % of the summed mid return earned on FOMC days, mean of the six full-set forecasts
\providecommand{\rhedCpiN}{}\renewcommand{\rhedCpiN}{25}  % CPI days among the 614 post-warmup days
\providecommand{\rhedCpiR}{}\renewcommand{\rhedCpiR}{\ensuremath{-}12}  % mean mid straddle return R on CPI days, %
\providecommand{\rhedCpiRT}{}\renewcommand{\rhedCpiRT}{\ensuremath{-}0.55}  % ... its t (i.i.d.; the days are weeks apart)
\providecommand{\rhedCpiLong}{}\renewcommand{\rhedCpiLong}{16}  % % of CPI days long, mean of the six full-set forecasts
\providecommand{\rhedCpiStrat}{}\renewcommand{\rhedCpiStrat}{22}  % mean mid strategy return on CPI days, %, mean of the six full-set
\providecommand{\rhedCpiOtherR}{}\renewcommand{\rhedCpiOtherR}{\ensuremath{-}6}  % mean straddle R on the other (covered) days, %
\providecommand{\rhedCpiOtherStrat}{}\renewcommand{\rhedCpiOtherStrat}{11}  % mean strategy return on the other (covered) days, %
\providecommand{\rhedCpiRVIV}{}\renewcommand{\rhedCpiRVIV}{0.80}  % mean RV_t / IV_t^2 on CPI days
\providecommand{\rhedCpiMoveSqIV}{}\renewcommand{\rhedCpiMoveSqIV}{0.86}  % mean squared S&P 500 half-hour log return / IV_t^2 on CPI days
\providecommand{\rhedCpiShare}{}\renewcommand{\rhedCpiShare}{7}  % % of the summed mid return earned on CPI days, mean of the six full-set forecasts
\providecommand{\rhedNfpN}{}\renewcommand{\rhedNfpN}{27}  % NFP days among the 614 post-warmup days
\providecommand{\rhedNfpR}{}\renewcommand{\rhedNfpR}{1}  % mean mid straddle return R on NFP days, %
\providecommand{\rhedNfpRT}{}\renewcommand{\rhedNfpRT}{0.06}  % ... its t (i.i.d.; the days are weeks apart)
\providecommand{\rhedNfpLong}{}\renewcommand{\rhedNfpLong}{12}  % % of NFP days long, mean of the six full-set forecasts
\providecommand{\rhedNfpStrat}{}\renewcommand{\rhedNfpStrat}{20}  % mean mid strategy return on NFP days, %, mean of the six full-set
\providecommand{\rhedNfpOtherR}{}\renewcommand{\rhedNfpOtherR}{\ensuremath{-}6}  % mean straddle R on the other (covered) days, %
\providecommand{\rhedNfpOtherStrat}{}\renewcommand{\rhedNfpOtherStrat}{11}  % mean strategy return on the other (covered) days, %
\providecommand{\rhedNfpRVIV}{}\renewcommand{\rhedNfpRVIV}{0.73}  % mean RV_t / IV_t^2 on NFP days
\providecommand{\rhedNfpMoveSqIV}{}\renewcommand{\rhedNfpMoveSqIV}{0.93}  % mean squared S&P 500 half-hour log return / IV_t^2 on NFP days
\providecommand{\rhedNfpShare}{}\renewcommand{\rhedNfpShare}{7}  % % of the summed mid return earned on NFP days, mean of the six full-set forecasts
\providecommand{\rhedOpexN}{}\renewcommand{\rhedOpexN}{33}  % OpEx days among the 614 post-warmup days
\providecommand{\rhedOpexR}{}\renewcommand{\rhedOpexR}{\ensuremath{-}45}  % mean mid straddle return R on OpEx days, %
\providecommand{\rhedOpexRT}{}\renewcommand{\rhedOpexRT}{\ensuremath{-}4.01}  % ... its t (i.i.d.; the days are weeks apart)
\providecommand{\rhedOpexLong}{}\renewcommand{\rhedOpexLong}{9}  % % of OpEx days long, mean of the six full-set forecasts
\providecommand{\rhedOpexStrat}{}\renewcommand{\rhedOpexStrat}{36}  % mean mid strategy return on OpEx days, %, mean of the six full-set
\providecommand{\rhedOpexOtherR}{}\renewcommand{\rhedOpexOtherR}{0}  % mean straddle R on the other (covered) days, %
\providecommand{\rhedOpexOtherStrat}{}\renewcommand{\rhedOpexOtherStrat}{11}  % mean strategy return on the other (covered) days, %
\providecommand{\rhedOpexRVIV}{}\renewcommand{\rhedOpexRVIV}{0.83}  % mean RV_t / IV_t^2 on OpEx days
\providecommand{\rhedOpexMoveSqIV}{}\renewcommand{\rhedOpexMoveSqIV}{0.40}  % mean squared S&P 500 half-hour log return / IV_t^2 on OpEx days
\providecommand{\rhedOpexShare}{}\renewcommand{\rhedOpexShare}{15}  % % of the summed mid return earned on OpEx days, mean of the six full-set forecasts
\providecommand{\rhedTwN}{}\renewcommand{\rhedTwN}{11}  % TW days among the 614 post-warmup days
\providecommand{\rhedTwR}{}\renewcommand{\rhedTwR}{\ensuremath{-}52}  % mean mid straddle return R on TW days, %
\providecommand{\rhedTwRT}{}\renewcommand{\rhedTwRT}{\ensuremath{-}3.98}  % ... its t (i.i.d.; the days are weeks apart)
\providecommand{\rhedTwLong}{}\renewcommand{\rhedTwLong}{8}  % % of TW days long, mean of the six full-set forecasts
\providecommand{\rhedTwStrat}{}\renewcommand{\rhedTwStrat}{37}  % mean mid strategy return on TW days, %, mean of the six full-set
\providecommand{\rhedTwOtherR}{}\renewcommand{\rhedTwOtherR}{\ensuremath{-}2}  % mean straddle R on the other (covered) days, %
\providecommand{\rhedTwOtherStrat}{}\renewcommand{\rhedTwOtherStrat}{12}  % mean strategy return on the other (covered) days, %
\providecommand{\rhedTwRVIV}{}\renewcommand{\rhedTwRVIV}{0.89}  % mean RV_t / IV_t^2 on TW days
\providecommand{\rhedTwMoveSqIV}{}\renewcommand{\rhedTwMoveSqIV}{0.31}  % mean squared S&P 500 half-hour log return / IV_t^2 on TW days
\providecommand{\rhedTwShare}{}\renewcommand{\rhedTwShare}{5}  % % of the summed mid return earned on TW days, mean of the six full-set forecasts
\providecommand{\rhedMeN}{}\renewcommand{\rhedMeN}{33}  % ME days among the 614 post-warmup days
\providecommand{\rhedMeR}{}\renewcommand{\rhedMeR}{70}  % mean mid straddle return R on ME days, %
\providecommand{\rhedMeRT}{}\renewcommand{\rhedMeRT}{2.79}  % ... its t (i.i.d.; the days are weeks apart)
\providecommand{\rhedMeLong}{}\renewcommand{\rhedMeLong}{6}  % % of ME days long, mean of the six full-set forecasts
\providecommand{\rhedMeStrat}{}\renewcommand{\rhedMeStrat}{\ensuremath{-}58}  % mean mid strategy return on ME days, %, mean of the six full-set
\providecommand{\rhedMeOtherR}{}\renewcommand{\rhedMeOtherR}{\ensuremath{-}7}  % mean straddle R on the other (covered) days, %
\providecommand{\rhedMeOtherStrat}{}\renewcommand{\rhedMeOtherStrat}{17}  % mean strategy return on the other (covered) days, %
\providecommand{\rhedMeRVIV}{}\renewcommand{\rhedMeRVIV}{1.04}  % mean RV_t / IV_t^2 on ME days
\providecommand{\rhedMeMoveSqIV}{}\renewcommand{\rhedMeMoveSqIV}{1.93}  % mean squared S&P 500 half-hour log return / IV_t^2 on ME days
\providecommand{\rhedMeShare}{}\renewcommand{\rhedMeShare}{\ensuremath{-}24}  % % of the summed mid return earned on ME days, mean of the six full-set forecasts
\providecommand{\rhedRelEndText}{}\renewcommand{\rhedRelEndText}{26~October~2023}  % earliest last flag of the six macro-release series in data/releases.parquet
\providecommand{\rhedRelFomcEndText}{}\renewcommand{\rhedRelFomcEndText}{1~November~2023}  % last FOMC flag in data/releases.parquet
\providecommand{\rhedFomcSuppN}{}\renewcommand{\rhedFomcSuppN}{3}  % FOMC statement days of the sample taken from live/close_signal/state/fomc_statement_dates.csv
\providecommand{\rhedRelAfterN}{}\renewcommand{\rhedRelAfterN}{126}  % post-warmup days after the release calendar ends
\providecommand{\rhedWorthN}{}\renewcommand{\rhedWorthN}{137}  % post-warmup days with R_t = -1 (= rheWorthlessDays)
\providecommand{\rhedWorthPct}{}\renewcommand{\rhedWorthPct}{22.3}  % ... % of the 614 days
\providecommand{\rhedWorthExpIV}{}\renewcommand{\rhedWorthExpIV}{128}  % expected R = -1 days under a normal half-hour log move with volatility IV_t
\providecommand{\rhedWorthExpIVPct}{}\renewcommand{\rhedWorthExpIVPct}{20.9}  % ... as % of the 614 days
\providecommand{\rhedWorthZ}{}\renewcommand{\rhedWorthZ}{0.90}  % (observed - expected) / sd of the count under that normal (independent days)
\providecommand{\rhedWorthExpEmp}{}\renewcommand{\rhedWorthExpEmp}{132}  % expected R = -1 days under the empirical distribution of the earlier days' moves standardized by IV
\providecommand{\rhedWorthExpRaw}{}\renewcommand{\rhedWorthExpRaw}{118}  % expected R = -1 days under the earlier days' raw half-hour moves
\providecommand{\rhedGapFiveN}{}\renewcommand{\rhedGapFiveN}{611}  % post-warmup days with K_C - K_P = 5 points
\providecommand{\rhedGapOtherN}{}\renewcommand{\rhedGapOtherN}{3}  % post-warmup days with another strike gap
\providecommand{\rhedMoveSdPts}{}\renewcommand{\rhedMoveSdPts}{9.6}  % median one-sd implied half-hour move, index points
\providecommand{\rhedStrikeDist}{}\renewcommand{\rhedStrikeDist}{1.23}  % mean distance of the close to the nearest 5-point strike, points
\providecommand{\rhedStrikeDistUnif}{}\renewcommand{\rhedStrikeDistUnif}{1.25}  % ... its mean under a uniform close
\providecommand{\rhedKsD}{}\renewcommand{\rhedKsD}{0.026}  % Kolmogorov-Smirnov D of the close's position within its 5-point strike interval against a uniform
\providecommand{\rhedKsCV}{}\renewcommand{\rhedKsCV}{0.055}  % ... asymptotic 5% critical value
\providecommand{\rhedZSd}{}\renewcommand{\rhedZSd}{1.00}  % sd of the half-hour log move standardized by IV_t, 614 days
\providecommand{\rhedZKurt}{}\renewcommand{\rhedZKurt}{2.0}  % excess kurtosis of that standardized move
\providecommand{\rhedWorthShortMin}{}\renewcommand{\rhedWorthShortMin}{110}  % lowest R = -1 days on which the forecast is short across the six full-set pooled forecasts, 614 days
\providecommand{\rhedWorthShortMax}{}\renewcommand{\rhedWorthShortMax}{115}  % highest R = -1 days on which the forecast is short across the six full-set pooled forecasts, 614 days
\providecommand{\rhedWorthOpexN}{}\renewcommand{\rhedWorthOpexN}{10}  % R = -1 days among the monthly-expiration days
\providecommand{\rhedWorthOpexExp}{}\renewcommand{\rhedWorthOpexExp}{6.1}  % ... expected under the normal with IV_t
\providecommand{\rhedWorthOtherPct}{}\renewcommand{\rhedWorthOtherPct}{21.9}  % % of the other days with R = -1
\providecommand{\rhedWorthOtherExpPct}{}\renewcommand{\rhedWorthOtherExpPct}{21.0}  % ... expected under the normal with IV_t, %
\providecommand{\rhedWorthOpexZ}{}\renewcommand{\rhedWorthOpexZ}{1.77}  % (observed - expected) / sd of the count on the monthly-expiration days
\providecommand{\rhedRegAlphaAMin}{}\renewcommand{\rhedRegAlphaAMin}{11.0}  % lowest alpha, (a) on the short straddle, mid, % per day across the six full-set pooled forecasts, 614 days
\providecommand{\rhedRegAlphaAMax}{}\renewcommand{\rhedRegAlphaAMax}{13.4}  % highest alpha, (a) on the short straddle, mid, % per day across the six full-set pooled forecasts, 614 days
\providecommand{\rhedRegAlphaAXMin}{}\renewcommand{\rhedRegAlphaAXMin}{8.3}  % lowest alpha, (a) on the short straddle, bid/ask, % per day across the six full-set pooled forecasts, 614 days
\providecommand{\rhedRegAlphaAXMax}{}\renewcommand{\rhedRegAlphaAXMax}{10.7}  % highest alpha, (a) on the short straddle, bid/ask, % per day across the six full-set pooled forecasts, 614 days
\providecommand{\rhedRegBetaShortMin}{}\renewcommand{\rhedRegBetaShortMin}{0.23}  % lowest beta on the short-straddle return, (a), mid across the six full-set pooled forecasts, 614 days
\providecommand{\rhedRegBetaShortMax}{}\renewcommand{\rhedRegBetaShortMax}{0.28}  % highest beta on the short-straddle return, (a), mid across the six full-set pooled forecasts, 614 days
\providecommand{\rhedRegAlphaAPMin}{}\renewcommand{\rhedRegAlphaAPMin}{0.005}  % lowest fixed-b p of alpha, (a), mid across the six full-set pooled forecasts, 614 days
\providecommand{\rhedRegAlphaAPMax}{}\renewcommand{\rhedRegAlphaAPMax}{0.025}  % highest fixed-b p of alpha, (a), mid across the six full-set pooled forecasts, 614 days
\providecommand{\rhedRegAlphaAXPMin}{}\renewcommand{\rhedRegAlphaAXPMin}{0.021}  % lowest fixed-b p of alpha, (a), bid/ask across the six full-set pooled forecasts, 614 days
\providecommand{\rhedRegAlphaAXPMax}{}\renewcommand{\rhedRegAlphaAXPMax}{0.077}  % highest fixed-b p of alpha, (a), bid/ask across the six full-set pooled forecasts, 614 days
\providecommand{\rhedRegBetaShortTMin}{}\renewcommand{\rhedRegBetaShortTMin}{1.22}  % lowest t of the beta on the short straddle, (a), mid across the six full-set pooled forecasts, 614 days
\providecommand{\rhedRegBetaShortTMax}{}\renewcommand{\rhedRegBetaShortTMax}{1.42}  % highest t of the beta on the short straddle, (a), mid across the six full-set pooled forecasts, 614 days
\providecommand{\rhedRegRsqAMin}{}\renewcommand{\rhedRegRsqAMin}{0.06}  % lowest R^2 of (a), mid across the six full-set pooled forecasts, 614 days
\providecommand{\rhedRegRsqAMax}{}\renewcommand{\rhedRegRsqAMax}{0.08}  % highest R^2 of (a), mid across the six full-set pooled forecasts, 614 days
\providecommand{\rhedRegBetaAbsMin}{}\renewcommand{\rhedRegBetaAbsMin}{\ensuremath{-}1.87}  % lowest beta on |r30| (bp), (b), mid, % of premium per bp across the six full-set pooled forecasts, 614 days
\providecommand{\rhedRegBetaAbsMax}{}\renewcommand{\rhedRegBetaAbsMax}{\ensuremath{-}1.40}  % highest beta on |r30| (bp), (b), mid, % of premium per bp across the six full-set pooled forecasts, 614 days
\providecommand{\rhedRegBetaAbsTMin}{}\renewcommand{\rhedRegBetaAbsTMin}{\ensuremath{-}3.49}  % lowest t of the beta on |r30|, (b), mid across the six full-set pooled forecasts, 614 days
\providecommand{\rhedRegBetaAbsTMax}{}\renewcommand{\rhedRegBetaAbsTMax}{\ensuremath{-}2.31}  % highest t of the beta on |r30|, (b), mid across the six full-set pooled forecasts, 614 days
\providecommand{\rhedRegBetaMoveMin}{}\renewcommand{\rhedRegBetaMoveMin}{\ensuremath{-}0.01}  % lowest beta on r30 (bp), (b), mid across the six full-set pooled forecasts, 614 days
\providecommand{\rhedRegBetaMoveMax}{}\renewcommand{\rhedRegBetaMoveMax}{0.14}  % highest beta on r30 (bp), (b), mid across the six full-set pooled forecasts, 614 days
\providecommand{\rhedRegMeanMin}{}\renewcommand{\rhedRegMeanMin}{11.5}  % lowest mean return, mid, % per day across the six full-set pooled forecasts, 614 days
\providecommand{\rhedRegMeanMax}{}\renewcommand{\rhedRegMeanMax}{14.1}  % highest mean return, mid, % per day across the six full-set pooled forecasts, 614 days
\providecommand{\rhedRegMeanXMin}{}\renewcommand{\rhedRegMeanXMin}{8.9}  % lowest mean return, bid/ask, % per day across the six full-set pooled forecasts, 614 days
\providecommand{\rhedRegMeanXMax}{}\renewcommand{\rhedRegMeanXMax}{11.4}  % highest mean return, bid/ask, % per day across the six full-set pooled forecasts, 614 days
\providecommand{\rhedRegRsqCMin}{}\renewcommand{\rhedRegRsqCMin}{0.00}  % lowest R^2 of (c), mid across the six full-set pooled forecasts, 614 days
\providecommand{\rhedRegRsqCMax}{}\renewcommand{\rhedRegRsqCMax}{0.00}  % highest R^2 of (c), mid across the six full-set pooled forecasts, 614 days
\providecommand{\rhedRegCPMaxMin}{}\renewcommand{\rhedRegCPMaxMin}{0.74}  % lowest larger fixed-b p of the two betas of (c), mid across the six full-set pooled forecasts, 614 days
\providecommand{\rhedRegCPMaxMax}{}\renewcommand{\rhedRegCPMaxMax}{0.97}  % highest larger fixed-b p of the two betas of (c), mid across the six full-set pooled forecasts, 614 days
\providecommand{\rhedRegCPMinMin}{}\renewcommand{\rhedRegCPMinMin}{0.56}  % lowest smaller fixed-b p of the two betas of (c), mid across the six full-set pooled forecasts, 614 days
\providecommand{\rhedRegCPMinMax}{}\renewcommand{\rhedRegCPMinMax}{0.82}  % highest smaller fixed-b p of the two betas of (c), mid across the six full-set pooled forecasts, 614 days
\providecommand{\rhedRegAlphaDMin}{}\renewcommand{\rhedRegAlphaDMin}{11.4}  % lowest alpha, (d) joint, mid, % per day across the six full-set pooled forecasts, 614 days
\providecommand{\rhedRegAlphaDMax}{}\renewcommand{\rhedRegAlphaDMax}{13.6}  % highest alpha, (d) joint, mid, % per day across the six full-set pooled forecasts, 614 days
\providecommand{\rhedRegAlphaDXMin}{}\renewcommand{\rhedRegAlphaDXMin}{8.7}  % lowest alpha, (d) joint, bid/ask, % per day across the six full-set pooled forecasts, 614 days
\providecommand{\rhedRegAlphaDXMax}{}\renewcommand{\rhedRegAlphaDXMax}{10.9}  % highest alpha, (d) joint, bid/ask, % per day across the six full-set pooled forecasts, 614 days
\providecommand{\rhedRegAlphaDPMin}{}\renewcommand{\rhedRegAlphaDPMin}{0.004}  % lowest fixed-b p of alpha, (d), mid across the six full-set pooled forecasts, 614 days
\providecommand{\rhedRegAlphaDPMax}{}\renewcommand{\rhedRegAlphaDPMax}{0.026}  % highest fixed-b p of alpha, (d), mid across the six full-set pooled forecasts, 614 days
\providecommand{\rhedRegAlphaDXPMin}{}\renewcommand{\rhedRegAlphaDXPMin}{0.017}  % lowest fixed-b p of alpha, (d), bid/ask across the six full-set pooled forecasts, 614 days
\providecommand{\rhedRegAlphaDXPMax}{}\renewcommand{\rhedRegAlphaDXPMax}{0.076}  % highest fixed-b p of alpha, (d), bid/ask across the six full-set pooled forecasts, 614 days
\providecommand{\rhedRegAvgAlphaM}{}\renewcommand{\rhedRegAvgAlphaM}{12.7}  % alpha of the mean of the six full-set returns, (m), mid, % per day
\providecommand{\rhedRegAvgAlphaMT}{}\renewcommand{\rhedRegAvgAlphaMT}{3.48}  % ... its t
\providecommand{\rhedRegAvgAlphaMP}{}\renewcommand{\rhedRegAvgAlphaMP}{0.002}  % ... its fixed-b p
\providecommand{\rhedRegAvgAlphaA}{}\renewcommand{\rhedRegAvgAlphaA}{12.1}  % alpha of the mean of the six full-set returns, (a), mid, % per day
\providecommand{\rhedRegAvgAlphaAT}{}\renewcommand{\rhedRegAvgAlphaAT}{2.83}  % ... its t
\providecommand{\rhedRegAvgAlphaAP}{}\renewcommand{\rhedRegAvgAlphaAP}{0.0100}  % ... its fixed-b p
\providecommand{\rhedRegAvgAlphaD}{}\renewcommand{\rhedRegAvgAlphaD}{12.5}  % alpha of the mean of the six full-set returns, (d), mid, % per day
\providecommand{\rhedRegAvgAlphaDT}{}\renewcommand{\rhedRegAvgAlphaDT}{2.86}  % ... its t
\providecommand{\rhedRegAvgAlphaDP}{}\renewcommand{\rhedRegAvgAlphaDP}{0.009}  % ... its fixed-b p
\providecommand{\rhedRegAvgAlphaMX}{}\renewcommand{\rhedRegAvgAlphaMX}{10.1}  % alpha of the mean of the six full-set returns, (m), ba, % per day
\providecommand{\rhedRegAvgAlphaMXT}{}\renewcommand{\rhedRegAvgAlphaMXT}{2.78}  % ... its t
\providecommand{\rhedRegAvgAlphaMXP}{}\renewcommand{\rhedRegAvgAlphaMXP}{0.011}  % ... its fixed-b p
\providecommand{\rhedRegAvgAlphaAX}{}\renewcommand{\rhedRegAvgAlphaAX}{9.4}  % alpha of the mean of the six full-set returns, (a), ba, % per day
\providecommand{\rhedRegAvgAlphaAXT}{}\renewcommand{\rhedRegAvgAlphaAXT}{2.21}  % ... its t
\providecommand{\rhedRegAvgAlphaAXP}{}\renewcommand{\rhedRegAvgAlphaAXP}{0.040}  % ... its fixed-b p
\providecommand{\rhedRegAvgAlphaDX}{}\renewcommand{\rhedRegAvgAlphaDX}{9.8}  % alpha of the mean of the six full-set returns, (d), ba, % per day
\providecommand{\rhedRegAvgAlphaDXT}{}\renewcommand{\rhedRegAvgAlphaDXT}{2.26}  % ... its t
\providecommand{\rhedRegAvgAlphaDXP}{}\renewcommand{\rhedRegAvgAlphaDXP}{0.036}  % ... its fixed-b p
\providecommand{\rhedRegAvgBetaShort}{}\renewcommand{\rhedRegAvgBetaShort}{0.26}  % beta of the mean full-set return on the short straddle, (a), mid
\providecommand{\rhedRegAvgBetaShortT}{}\renewcommand{\rhedRegAvgBetaShortT}{1.34}  % ... its t
\providecommand{\rhedRegAvgBetaShortP}{}\renewcommand{\rhedRegAvgBetaShortP}{0.20}  % ... its fixed-b p
\providecommand{\rhedRegAvgRsqA}{}\renewcommand{\rhedRegAvgRsqA}{0.07}  % R^2 of (a) for the mean full-set return
\providecommand{\rhedRegAvgBetaAbs}{}\renewcommand{\rhedRegAvgBetaAbs}{\ensuremath{-}1.68}  % beta of the mean full-set return on |r30| (bp), (b), mid
\providecommand{\rhedRegAvgBetaAbsT}{}\renewcommand{\rhedRegAvgBetaAbsT}{\ensuremath{-}3.26}  % ... its t
\providecommand{\rhedRegShortBetaAbs}{}\renewcommand{\rhedRegShortBetaAbs}{\ensuremath{-}3.96}  % beta of the short straddle on |r30| (bp), (b), mid
\providecommand{\rhedRegShortRsqB}{}\renewcommand{\rhedRegShortRsqB}{0.48}  % R^2 of (b) for the short straddle
\providecommand{\rhedRegAbsRatio}{}\renewcommand{\rhedRegAbsRatio}{43}  % beta on |r30| of the mean full-set return as % of the short straddle's
\providecommand{\rhedRegAvgBetaAbsJ}{}\renewcommand{\rhedRegAvgBetaAbsJ}{\ensuremath{-}1.30}  % beta of the mean full-set return on |r30|, (d) joint, mid
\providecommand{\rhedRegAvgBetaAbsJT}{}\renewcommand{\rhedRegAvgBetaAbsJT}{\ensuremath{-}1.90}  % ... its t
\providecommand{\rhedHarRegAlphaAMin}{}\renewcommand{\rhedHarRegAlphaAMin}{1.1}  % lowest alpha, (a), mid, % per day across the six HAR-benchmark pooled forecasts, 614 days
\providecommand{\rhedHarRegAlphaAMax}{}\renewcommand{\rhedHarRegAlphaAMax}{7.0}  % highest alpha, (a), mid, % per day across the six HAR-benchmark pooled forecasts, 614 days
\providecommand{\rhedHarRegAlphaAPMin}{}\renewcommand{\rhedHarRegAlphaAPMin}{0.082}  % lowest fixed-b p of alpha, (a), mid across the six HAR-benchmark pooled forecasts, 614 days
\providecommand{\rhedHarRegAlphaAPMax}{}\renewcommand{\rhedHarRegAlphaAPMax}{0.745}  % highest fixed-b p of alpha, (a), mid across the six HAR-benchmark pooled forecasts, 614 days
\providecommand{\rhedOlsRegAlphaA}{}\renewcommand{\rhedOlsRegAlphaA}{8.1}  % alpha, (a), mid, closing-interval HAR forecast, % per day
\providecommand{\rhedOlsRegAlphaAP}{}\renewcommand{\rhedOlsRegAlphaAP}{0.051}  % ... its fixed-b p
\providecommand{\rhedRegAlphaAExBestMin}{}\renewcommand{\rhedRegAlphaAExBestMin}{\ensuremath{-}0.5}  % lowest alpha, (a), mid, without the forecast's ten best days, % per day across the six full-set pooled forecasts, 614 days
\providecommand{\rhedRegAlphaAExBestMax}{}\renewcommand{\rhedRegAlphaAExBestMax}{2.1}  % highest alpha, (a), mid, without the forecast's ten best days, % per day across the six full-set pooled forecasts, 614 days
\providecommand{\rhedRegAlphaAExBestPMin}{}\renewcommand{\rhedRegAlphaAExBestPMin}{0.43}  % lowest ... its fixed-b p across the six full-set pooled forecasts, 614 days
\providecommand{\rhedRegAlphaAExBestPMax}{}\renewcommand{\rhedRegAlphaAExBestPMax}{0.97}  % highest ... its fixed-b p across the six full-set pooled forecasts, 614 days
\fi
\begin{table}[tbp]
\centering
\small
\setlength{\tabcolsep}{4pt}
\begin{threeparttable}
\caption{Straddle returns by quintile of the forecast variance gap}
\label{tab:rh_econ_quintiles}
\begin{tabular}{@{}lrrrrrrrr@{}}
\toprule
 & & & & Straddle & \multicolumn{2}{c}{Strategy} & & Share of \\
\cmidrule(lr){6-7}
Quintile & $\hat\sigma^2_t/\mathrm{IV}_t^2$ & Long (\%) & $\RV_t/\mathrm{IV}_t^2$ & $R_t$ (\%) & Mid (\%) & Bid/ask (\%) & Hit (\%) & return (\%) \\
\midrule
\multicolumn{9}{@{}l}{\textit{Panel A: Full predictor set (mean of the six forecasts)}} \\
\addlinespace[2pt]
1 (lowest) & $0.47$ & $0.0$ & $0.73$ & $-9.3$ & $9.3$ & $7.0$ & $88.1$ & $14.7$ \\
2 & $0.65$ & $0.0$ & $0.78$ & $-13.3$ & $13.3$ & $11.1$ & $80.8$ & $21.0$ \\
3 & $0.78$ & $0.0$ & $0.85$ & $-6.0$ & $6.0$ & $3.3$ & $74.7$ & $9.3$ \\
4 & $0.93$ & $7.0$ & $0.93$ & $-10.0$ & $7.8$ & $5.3$ & $64.0$ & $12.2$ \\
5 (highest) & $1.18$ & $98.9$ & $1.05$ & $26.5$ & $27.2$ & $23.7$ & $51.9$ & $42.8$ \\
\addlinespace[2pt]
All & $0.77$ & $21.2$ & $0.87$ & $-2.4$ & $12.7$ & $10.1$ & $71.9$ & $100.0$ \\
\addlinespace[4pt]
\multicolumn{9}{@{}l}{\textit{Panel B: HAR benchmark (mean of the six forecasts)}} \\
\addlinespace[2pt]
1 (lowest) & $0.53$ & $0.0$ & $0.74$ & $-8.1$ & $8.1$ & $5.9$ & $88.9$ & $26.7$ \\
2 & $0.71$ & $0.0$ & $0.81$ & $-3.8$ & $3.8$ & $1.3$ & $80.2$ & $12.7$ \\
3 & $0.83$ & $0.0$ & $0.81$ & $-12.8$ & $12.8$ & $10.4$ & $76.1$ & $42.0$ \\
4 & $0.98$ & $35.9$ & $0.96$ & $5.4$ & $-1.5$ & $-4.6$ & $58.1$ & $-4.9$ \\
5 (highest) & $1.23$ & $100.0$ & $1.02$ & $7.1$ & $7.1$ & $4.2$ & $47.3$ & $23.5$ \\
\addlinespace[2pt]
All & $0.82$ & $27.2$ & $0.87$ & $-2.4$ & $6.0$ & $3.4$ & $70.1$ & $100.0$ \\
\bottomrule
\end{tabular}
\begin{tablenotes}[flushleft]
\footnotesize
\item \textit{Notes:} For each pooled intraday forecast, the \rheSzDays{} post-warmup option days are sorted into
quintiles of $\log(\hat\sigma^2_t/\mathrm{IV}_t^2)$, the forecast variance gap that signs the straddle position
(Section~\ref{sec:rh_econ_strategy}); breakpoints from the whole sample, about 123 days a quintile. Each entry is the mean
over the six forecasts of a panel of the quintile's statistic: $\hat\sigma^2_t/\mathrm{IV}_t^2$, the geometric mean of
the forecast ratio; Long, the share of days with $q_t = +1$; $\RV_t/\mathrm{IV}_t^2$, the mean ratio of realized to
implied variance; $R_t$, the mean return of a long straddle at mid-quotes; Strategy, the mean return of the sign strategy
at mid-quotes and at quoted bid/ask prices; Hit, the share of days on which $q_t$ has the sign of
$\RV_t - \mathrm{IV}_t^2$; Share of return, the quintile's part of the summed mid-quote returns of the six
forecasts. Returns in \% of the premium.
\end{tablenotes}
\end{threeparttable}
\end{table}
    % tab:rh_econ_quintiles
%   % AUTO-GENERATED by writeup/make_rehaul_econ_decomp.py -- do not edit; re-run the script.
% Table tab:rh_econ_bestdays: the ten best and ten worst days of each full-set pooled forecast, their union
% over the six, and what they contribute to Delta SR (writeup/make_rehaul_econ_decomp.py, statistic 4).
\ifdefined\rhedNDays\else% AUTO-GENERATED by writeup/make_rehaul_econ_decomp.py -- do not edit; re-run the script.
% \rhed... macros: where the straddle strategy's profits come from (Section 7 decomposition).
\providecommand{\rhedNDays}{}\renewcommand{\rhedNDays}{614}  % post-warmup option days (= rheSzDays)
\providecommand{\rhedNFull}{}\renewcommand{\rhedNFull}{6}  % pooled intraday forecasts on the full predictor set
\providecommand{\rhedNTop}{}\renewcommand{\rhedNTop}{10}  % best / worst days per forecast
\providecommand{\rhedUncRVAboveIV}{}\renewcommand{\rhedUncRVAboveIV}{28.2}  % % of the 614 days with RV_t > IV_t^2
\providecommand{\rhedUncRVIVMean}{}\renewcommand{\rhedUncRVIVMean}{0.87}  % mean RV_t / IV_t^2, 614 days
\providecommand{\rhedUncRVIVMed}{}\renewcommand{\rhedUncRVIVMed}{0.77}  % median RV_t / IV_t^2, 614 days
\providecommand{\rhedUncShortMean}{}\renewcommand{\rhedUncShortMean}{2.4}  % mean mid return of the unconditional short straddle, % of premium
\providecommand{\rhedUncShortHit}{}\renewcommand{\rhedUncShortHit}{71.8}  % % of days on which q = -1 matches sign(RV - IV^2) (always short)
\providecommand{\rhedLongNMin}{}\renewcommand{\rhedLongNMin}{117}  % lowest number of long days across the six full-set pooled forecasts, 614 days
\providecommand{\rhedLongNMax}{}\renewcommand{\rhedLongNMax}{146}  % highest number of long days across the six full-set pooled forecasts, 614 days
\providecommand{\rhedLongPctMin}{}\renewcommand{\rhedLongPctMin}{19}  % lowest % long days across the six full-set pooled forecasts, 614 days
\providecommand{\rhedLongPctMax}{}\renewcommand{\rhedLongPctMax}{24}  % highest % long days across the six full-set pooled forecasts, 614 days
\providecommand{\rhedHarLongNMin}{}\renewcommand{\rhedHarLongNMin}{150}  % lowest number of long days across the six HAR-benchmark pooled forecasts, 614 days
\providecommand{\rhedHarLongNMax}{}\renewcommand{\rhedHarLongNMax}{193}  % highest number of long days across the six HAR-benchmark pooled forecasts, 614 days
\providecommand{\rhedLongMidMin}{}\renewcommand{\rhedLongMidMin}{19}  % lowest mean mid return on long days, % across the six full-set pooled forecasts, 614 days
\providecommand{\rhedLongMidMax}{}\renewcommand{\rhedLongMidMax}{30}  % highest mean mid return on long days, % across the six full-set pooled forecasts, 614 days
\providecommand{\rhedLongBaMin}{}\renewcommand{\rhedLongBaMin}{16}  % lowest mean bid/ask return on long days, % across the six full-set pooled forecasts, 614 days
\providecommand{\rhedLongBaMax}{}\renewcommand{\rhedLongBaMax}{26}  % highest mean bid/ask return on long days, % across the six full-set pooled forecasts, 614 days
\providecommand{\rhedShortMidMin}{}\renewcommand{\rhedShortMidMin}{9}  % lowest mean mid return on short days, % across the six full-set pooled forecasts, 614 days
\providecommand{\rhedShortMidMax}{}\renewcommand{\rhedShortMidMax}{11}  % highest mean mid return on short days, % across the six full-set pooled forecasts, 614 days
\providecommand{\rhedShortBaMin}{}\renewcommand{\rhedShortBaMin}{7}  % lowest mean bid/ask return on short days, % across the six full-set pooled forecasts, 614 days
\providecommand{\rhedShortBaMax}{}\renewcommand{\rhedShortBaMax}{8}  % highest mean bid/ask return on short days, % across the six full-set pooled forecasts, 614 days
\providecommand{\rhedHarLongMidMin}{}\renewcommand{\rhedHarLongMidMin}{0}  % lowest mean mid return on long days, % across the six HAR-benchmark pooled forecasts, 614 days
\providecommand{\rhedHarLongMidMax}{}\renewcommand{\rhedHarLongMidMax}{10}  % highest mean mid return on long days, % across the six HAR-benchmark pooled forecasts, 614 days
\providecommand{\rhedHarShortMidMin}{}\renewcommand{\rhedHarShortMidMin}{3}  % lowest mean mid return on short days, % across the six HAR-benchmark pooled forecasts, 614 days
\providecommand{\rhedHarShortMidMax}{}\renewcommand{\rhedHarShortMidMax}{7}  % highest mean mid return on short days, % across the six HAR-benchmark pooled forecasts, 614 days
\providecommand{\rhedLongHitMin}{}\renewcommand{\rhedLongHitMin}{47}  % lowest % of long days with RV > IV^2 across the six full-set pooled forecasts, 614 days
\providecommand{\rhedLongHitMax}{}\renewcommand{\rhedLongHitMax}{53}  % highest % of long days with RV > IV^2 across the six full-set pooled forecasts, 614 days
\providecommand{\rhedHarLongHitMin}{}\renewcommand{\rhedHarLongHitMin}{46}  % lowest % of long days with RV > IV^2 across the six HAR-benchmark pooled forecasts, 614 days
\providecommand{\rhedHarLongHitMax}{}\renewcommand{\rhedHarLongHitMax}{48}  % highest % of long days with RV > IV^2 across the six HAR-benchmark pooled forecasts, 614 days
\providecommand{\rhedLongPosMin}{}\renewcommand{\rhedLongPosMin}{42}  % lowest % of long days with a positive return across the six full-set pooled forecasts, 614 days
\providecommand{\rhedLongPosMax}{}\renewcommand{\rhedLongPosMax}{45}  % highest % of long days with a positive return across the six full-set pooled forecasts, 614 days
\providecommand{\rhedHitMin}{}\renewcommand{\rhedHitMin}{70.7}  % lowest % of days on which q matches sign(RV - IV^2) across the six full-set pooled forecasts, 614 days
\providecommand{\rhedHitMax}{}\renewcommand{\rhedHitMax}{73.0}  % highest % of days on which q matches sign(RV - IV^2) across the six full-set pooled forecasts, 614 days
\providecommand{\rhedPosMin}{}\renewcommand{\rhedPosMin}{58.6}  % lowest % of days with a positive mid return across the six full-set pooled forecasts, 614 days
\providecommand{\rhedPosMax}{}\renewcommand{\rhedPosMax}{60.4}  % highest % of days with a positive mid return across the six full-set pooled forecasts, 614 days
\providecommand{\rhedLongShareMin}{}\renewcommand{\rhedLongShareMin}{39}  % lowest % of the summed mid return earned on long days across the six full-set pooled forecasts, 614 days
\providecommand{\rhedLongShareMax}{}\renewcommand{\rhedLongShareMax}{41}  % highest % of the summed mid return earned on long days across the six full-set pooled forecasts, 614 days
\providecommand{\rhedLongShareXMin}{}\renewcommand{\rhedLongShareXMin}{42}  % lowest % of the summed bid/ask return earned on long days across the six full-set pooled forecasts, 614 days
\providecommand{\rhedLongShareXMax}{}\renewcommand{\rhedLongShareXMax}{45}  % highest % of the summed bid/ask return earned on long days across the six full-set pooled forecasts, 614 days
\providecommand{\rhedMeanMin}{}\renewcommand{\rhedMeanMin}{11.5}  % lowest mean mid return, % across the six full-set pooled forecasts, 614 days
\providecommand{\rhedMeanMax}{}\renewcommand{\rhedMeanMax}{14.1}  % highest mean mid return, % across the six full-set pooled forecasts, 614 days
\providecommand{\rhedGainMin}{}\renewcommand{\rhedGainMin}{9.1}  % lowest mean gain over the short straddle (2 R 1(long)), % across the six full-set pooled forecasts, 614 days
\providecommand{\rhedGainMax}{}\renewcommand{\rhedGainMax}{11.6}  % highest mean gain over the short straddle (2 R 1(long)), % across the six full-set pooled forecasts, 614 days
\providecommand{\rhedGainShareMin}{}\renewcommand{\rhedGainShareMin}{79}  % lowest gain over short / mean return, % across the six full-set pooled forecasts, 614 days
\providecommand{\rhedGainShareMax}{}\renewcommand{\rhedGainShareMax}{83}  % highest gain over short / mean return, % across the six full-set pooled forecasts, 614 days
\providecommand{\rhedHarGainMin}{}\renewcommand{\rhedHarGainMin}{0.0}  % lowest mean gain over the short straddle, % across the six HAR-benchmark pooled forecasts, 614 days
\providecommand{\rhedHarGainMax}{}\renewcommand{\rhedHarGainMax}{5.0}  % highest mean gain over the short straddle, % across the six HAR-benchmark pooled forecasts, 614 days
\providecommand{\rhedOlsLongN}{}\renewcommand{\rhedOlsLongN}{198}  % long days of the closing-interval HAR forecast (sub_ols_baseline)
\providecommand{\rhedOlsLongMid}{}\renewcommand{\rhedOlsLongMid}{9}  % ... its mean mid return on long days, %
\providecommand{\rhedOlsShortMid}{}\renewcommand{\rhedOlsShortMid}{8}  % ... its mean mid return on short days, %
\providecommand{\rhedOlsLongHit}{}\renewcommand{\rhedOlsLongHit}{46}  % ... % of its long days with RV > IV^2
\providecommand{\rhedOlsSh}{}\renewcommand{\rhedOlsSh}{1.11}  % ... its Sharpe ratio, mid (= rhrOlsSharpe)
\providecommand{\rhedCalLevelMin}{}\renewcommand{\rhedCalLevelMin}{0.95}  % lowest mean forecast / mean RV across the six full-set pooled forecasts, 614 days
\providecommand{\rhedCalLevelMax}{}\renewcommand{\rhedCalLevelMax}{1.01}  % highest mean forecast / mean RV across the six full-set pooled forecasts, 614 days
\providecommand{\rhedCalScaleMin}{}\renewcommand{\rhedCalScaleMin}{1.09}  % lowest mean of RV / forecast (QLIKE-optimal scale) across the six full-set pooled forecasts, 614 days
\providecommand{\rhedCalScaleMax}{}\renewcommand{\rhedCalScaleMax}{1.20}  % highest mean of RV / forecast (QLIKE-optimal scale) across the six full-set pooled forecasts, 614 days
\providecommand{\rhedHarCalScaleMin}{}\renewcommand{\rhedHarCalScaleMin}{1.04}  % lowest mean of RV / forecast across the six HAR-benchmark pooled forecasts, 614 days
\providecommand{\rhedHarCalScaleMax}{}\renewcommand{\rhedHarCalScaleMax}{1.10}  % highest mean of RV / forecast across the six HAR-benchmark pooled forecasts, 614 days
\providecommand{\rhedMedFcIVMin}{}\renewcommand{\rhedMedFcIVMin}{0.75}  % lowest median forecast / IV^2 across the six full-set pooled forecasts, 614 days
\providecommand{\rhedMedFcIVMax}{}\renewcommand{\rhedMedFcIVMax}{0.82}  % highest median forecast / IV^2 across the six full-set pooled forecasts, 614 days
\providecommand{\rhedBcScaleLoMin}{}\renewcommand{\rhedBcScaleLoMin}{1.06}  % lowest lowest causal scale c_t over the 614 days across the six full-set pooled forecasts, 614 days
\providecommand{\rhedBcScaleLoMax}{}\renewcommand{\rhedBcScaleLoMax}{1.15}  % highest lowest causal scale c_t over the 614 days across the six full-set pooled forecasts, 614 days
\providecommand{\rhedBcScaleHiMin}{}\renewcommand{\rhedBcScaleHiMin}{1.10}  % lowest highest causal scale c_t over the 614 days across the six full-set pooled forecasts, 614 days
\providecommand{\rhedBcScaleHiMax}{}\renewcommand{\rhedBcScaleHiMax}{1.22}  % highest highest causal scale c_t over the 614 days across the six full-set pooled forecasts, 614 days
\providecommand{\rhedBcFlipMin}{}\renewcommand{\rhedBcFlipMin}{55}  % lowest positions changed by the bias correction across the six full-set pooled forecasts, 614 days
\providecommand{\rhedBcFlipMax}{}\renewcommand{\rhedBcFlipMax}{120}  % highest positions changed by the bias correction across the six full-set pooled forecasts, 614 days
\providecommand{\rhedBcLongMin}{}\renewcommand{\rhedBcLongMin}{201}  % lowest long days under the bias correction across the six full-set pooled forecasts, 614 days
\providecommand{\rhedBcLongMax}{}\renewcommand{\rhedBcLongMax}{241}  % highest long days under the bias correction across the six full-set pooled forecasts, 614 days
\providecommand{\rhedBcFlipRMin}{}\renewcommand{\rhedBcFlipRMin}{\ensuremath{-}24}  % lowest mean mid straddle return R on the changed days, % across the six full-set pooled forecasts, 614 days
\providecommand{\rhedBcFlipRMax}{}\renewcommand{\rhedBcFlipRMax}{\ensuremath{-}8}  % highest mean mid straddle return R on the changed days, % across the six full-set pooled forecasts, 614 days
\providecommand{\rhedBcShMin}{}\renewcommand{\rhedBcShMin}{0.92}  % lowest Sharpe ratio mid under the bias correction (diagnostic) across the six full-set pooled forecasts, 614 days
\providecommand{\rhedBcShMax}{}\renewcommand{\rhedBcShMax}{1.32}  % highest Sharpe ratio mid under the bias correction (diagnostic) across the six full-set pooled forecasts, 614 days
\providecommand{\rhedBcShXMin}{}\renewcommand{\rhedBcShXMin}{0.57}  % lowest Sharpe ratio bid/ask under the bias correction across the six full-set pooled forecasts, 614 days
\providecommand{\rhedBcShXMax}{}\renewcommand{\rhedBcShXMax}{0.97}  % highest Sharpe ratio bid/ask under the bias correction across the six full-set pooled forecasts, 614 days
\providecommand{\rhedLevFlipMin}{}\renewcommand{\rhedLevFlipMin}{4}  % lowest positions changed by the level correction (sum RV / sum fc) across the six full-set pooled forecasts, 614 days
\providecommand{\rhedLevFlipMax}{}\renewcommand{\rhedLevFlipMax}{19}  % highest positions changed by the level correction (sum RV / sum fc) across the six full-set pooled forecasts, 614 days
\providecommand{\rhedLevShMin}{}\renewcommand{\rhedLevShMin}{1.61}  % lowest Sharpe ratio mid under the level correction across the six full-set pooled forecasts, 614 days
\providecommand{\rhedLevShMax}{}\renewcommand{\rhedLevShMax}{2.01}  % highest Sharpe ratio mid under the level correction across the six full-set pooled forecasts, 614 days
\providecommand{\rhedHarBcShMin}{}\renewcommand{\rhedHarBcShMin}{0.76}  % lowest Sharpe ratio mid under the bias correction across the six HAR-benchmark pooled forecasts, 614 days
\providecommand{\rhedHarBcShMax}{}\renewcommand{\rhedHarBcShMax}{1.05}  % highest Sharpe ratio mid under the bias correction across the six HAR-benchmark pooled forecasts, 614 days
\providecommand{\rhedBcShUpN}{}\renewcommand{\rhedBcShUpN}{0}  % full-set forecasts whose mid SR rises under the bias correction
\providecommand{\rhedBcShDownN}{}\renewcommand{\rhedBcShDownN}{6}  % full-set forecasts whose mid SR falls under the bias correction
\providecommand{\rhedHarBcShUpN}{}\renewcommand{\rhedHarBcShUpN}{2}  % HAR-benchmark forecasts whose mid SR rises under the bias correction
\providecommand{\rhedMoveSqIV}{}\renewcommand{\rhedMoveSqIV}{1.01}  % mean of the squared S&P 500 half-hour log return over IV_t^2, 614 days
\providecommand{\rhedMoveSqRV}{}\renewcommand{\rhedMoveSqRV}{1.15}  % mean squared S&P 500 half-hour log return over mean RV_t, 614 days
\providecommand{\rhedOlsCalScale}{}\renewcommand{\rhedOlsCalScale}{1.01}  % mean of RV / forecast, closing-interval HAR forecast
\providecommand{\rhedOlsBcSh}{}\renewcommand{\rhedOlsBcSh}{1.12}  % Sharpe ratio mid under the bias correction, closing-interval HAR
\providecommand{\rhedBcFlipRidge}{}\renewcommand{\rhedBcFlipRidge}{105}  % positions changed by the bias correction, poolpt13_ridge_all_features
\providecommand{\rhedBcShRidge}{}\renewcommand{\rhedBcShRidge}{1.17}  % Sharpe ratio mid under the bias correction, poolpt13_ridge_all_features
\providecommand{\rhedBcFlipLasso}{}\renewcommand{\rhedBcFlipLasso}{120}  % positions changed by the bias correction, poolpt13_lasso_all_features
\providecommand{\rhedBcShLasso}{}\renewcommand{\rhedBcShLasso}{1.31}  % Sharpe ratio mid under the bias correction, poolpt13_lasso_all_features
\providecommand{\rhedBcFlipEnet}{}\renewcommand{\rhedBcFlipEnet}{108}  % positions changed by the bias correction, poolpt13_enet_all_features
\providecommand{\rhedBcShEnet}{}\renewcommand{\rhedBcShEnet}{0.99}  % Sharpe ratio mid under the bias correction, poolpt13_enet_all_features
\providecommand{\rhedBcFlipLgbm}{}\renewcommand{\rhedBcFlipLgbm}{55}  % positions changed by the bias correction, pooltree_lgbm_all_features
\providecommand{\rhedBcShLgbm}{}\renewcommand{\rhedBcShLgbm}{1.32}  % Sharpe ratio mid under the bias correction, pooltree_lgbm_all_features
\providecommand{\rhedBcFlipXgb}{}\renewcommand{\rhedBcFlipXgb}{61}  % positions changed by the bias correction, pooltree_xgb_all_features
\providecommand{\rhedBcShXgb}{}\renewcommand{\rhedBcShXgb}{1.15}  % Sharpe ratio mid under the bias correction, pooltree_xgb_all_features
\providecommand{\rhedBcFlipRf}{}\renewcommand{\rhedBcFlipRf}{90}  % positions changed by the bias correction, pooltree_rf_all_features
\providecommand{\rhedBcShRf}{}\renewcommand{\rhedBcShRf}{0.92}  % Sharpe ratio mid under the bias correction, pooltree_rf_all_features
\providecommand{\rhedQMonoN}{}\renewcommand{\rhedQMonoN}{6}  % full-set forecasts whose mean RV/IV^2 rises across all five quintiles
\providecommand{\rhedHarQMonoN}{}\renewcommand{\rhedHarQMonoN}{4}  % HAR-benchmark forecasts whose mean RV/IV^2 rises across all five quintiles
\providecommand{\rhedQRPatternN}{}\renewcommand{\rhedQRPatternN}{6}  % full-set forecasts with mean R < 0 in quintiles 1-4 and > 0 in 5
\providecommand{\rhedHarQRPatternN}{}\renewcommand{\rhedHarQRPatternN}{0}  % HAR-benchmark forecasts with that pattern
\providecommand{\rhedQOneRVIV}{}\renewcommand{\rhedQOneRVIV}{0.73}  % mean RV/IV^2 in the lowest quintile, mean of the six full-set
\providecommand{\rhedQFiveRVIV}{}\renewcommand{\rhedQFiveRVIV}{1.05}  % mean RV/IV^2 in the highest quintile, mean of the six full-set
\providecommand{\rhedQFiveR}{}\renewcommand{\rhedQFiveR}{26}  % mean straddle return R in the highest quintile, %, mean of the six full-set
\providecommand{\rhedQFiveFc}{}\renewcommand{\rhedQFiveFc}{1.18}  % geometric-mean forecast / IV^2 in the highest quintile, mean of the six full-set
\providecommand{\rhedQFourFc}{}\renewcommand{\rhedQFourFc}{0.93}  % ... in the fourth quintile
\providecommand{\rhedQFiveLong}{}\renewcommand{\rhedQFiveLong}{99}  % % long days in the highest quintile, mean of the six full-set
\providecommand{\rhedQFiveHit}{}\renewcommand{\rhedQFiveHit}{52}  % hit rate in the highest quintile, %, mean of the six full-set
\providecommand{\rhedQOneHit}{}\renewcommand{\rhedQOneHit}{88}  % hit rate in the lowest quintile, %, mean of the six full-set
\providecommand{\rhedQFiveShare}{}\renewcommand{\rhedQFiveShare}{43}  % % of the summed mid return earned in the highest quintile, six full-set forecasts
\providecommand{\rhedQLowRMin}{}\renewcommand{\rhedQLowRMin}{\ensuremath{-}13}  % most negative mean straddle return R in quintiles 1-4, %, mean of the six
\providecommand{\rhedQLowRMax}{}\renewcommand{\rhedQLowRMax}{\ensuremath{-}6}  % least negative mean straddle return R in quintiles 1-4, %, mean of the six
\providecommand{\rhedHarQOneRVIV}{}\renewcommand{\rhedHarQOneRVIV}{0.74}  % mean RV/IV^2 in the lowest quintile, mean of the six Har
\providecommand{\rhedHarQFiveRVIV}{}\renewcommand{\rhedHarQFiveRVIV}{1.02}  % mean RV/IV^2 in the highest quintile, mean of the six Har
\providecommand{\rhedHarQFiveR}{}\renewcommand{\rhedHarQFiveR}{7}  % mean straddle return R in the highest quintile, %, mean of the six Har
\providecommand{\rhedHarQFiveFc}{}\renewcommand{\rhedHarQFiveFc}{1.23}  % geometric-mean forecast / IV^2 in the highest quintile, mean of the six Har
\providecommand{\rhedHarQFourFc}{}\renewcommand{\rhedHarQFourFc}{0.98}  % ... in the fourth quintile
\providecommand{\rhedHarQFiveLong}{}\renewcommand{\rhedHarQFiveLong}{100}  % % long days in the highest quintile, mean of the six Har
\providecommand{\rhedHarQFiveHit}{}\renewcommand{\rhedHarQFiveHit}{47}  % hit rate in the highest quintile, %, mean of the six Har
\providecommand{\rhedHarQOneHit}{}\renewcommand{\rhedHarQOneHit}{89}  % hit rate in the lowest quintile, %, mean of the six Har
\providecommand{\rhedHarQFiveShare}{}\renewcommand{\rhedHarQFiveShare}{23}  % % of the summed mid return earned in the highest quintile, six Har forecasts
\providecommand{\rhedHarQLowRMin}{}\renewcommand{\rhedHarQLowRMin}{\ensuremath{-}13}  % most negative mean straddle return R in quintiles 1-4, %, mean of the six
\providecommand{\rhedHarQLowRMax}{}\renewcommand{\rhedHarQLowRMax}{5}  % least negative mean straddle return R in quintiles 1-4, %, mean of the six
\providecommand{\rhedQFiveShareRngMin}{}\renewcommand{\rhedQFiveShareRngMin}{40}  % lowest % of the summed mid return in quintile 5 across the six full-set pooled forecasts, 614 days
\providecommand{\rhedQFiveShareRngMax}{}\renewcommand{\rhedQFiveShareRngMax}{48}  % highest % of the summed mid return in quintile 5 across the six full-set pooled forecasts, 614 days
\providecommand{\rhedBestUnion}{}\renewcommand{\rhedBestUnion}{11}  % days in the union of the six full-set forecasts' ten best days
\providecommand{\rhedWorstUnion}{}\renewcommand{\rhedWorstUnion}{11}  % days in the union of the six full-set forecasts' ten worst days
\providecommand{\rhedBestAllSix}{}\renewcommand{\rhedBestAllSix}{9}  % days among the ten best of all six forecasts
\providecommand{\rhedWorstAllSix}{}\renewcommand{\rhedWorstAllSix}{9}  % days among the ten worst of all six forecasts
\providecommand{\rhedTopDate}{}\renewcommand{\rhedTopDate}{22~November~2021}  % the day of the largest straddle return among the best days
\providecommand{\rhedTopR}{}\renewcommand{\rhedTopR}{1{,}032}  % ... the straddle's mid return that day, % of premium
\providecommand{\rhedTopRVIV}{}\renewcommand{\rhedTopRVIV}{2.79}  % ... RV / IV^2 that day
\providecommand{\rhedTopMove}{}\renewcommand{\rhedTopMove}{\ensuremath{-}0.83}  % ... S&P 500 15:30-16:00 log return that day, %
\providecommand{\rhedTopShareMin}{}\renewcommand{\rhedTopShareMin}{12}  % lowest % of a full-set forecast's summed mid return earned on that one day
\providecommand{\rhedTopShareMax}{}\renewcommand{\rhedTopShareMax}{15}  % highest % of a full-set forecast's summed mid return earned on that one day
\providecommand{\rhedBestFomcN}{}\renewcommand{\rhedBestFomcN}{1}  % best-day union: FOMC statement days
\providecommand{\rhedWorstFomcN}{}\renewcommand{\rhedWorstFomcN}{1}  % worst-day union: FOMC statement days
\providecommand{\rhedWorstMeN}{}\renewcommand{\rhedWorstMeN}{3}  % worst-day union: month ends
\providecommand{\rhedBestMeN}{}\renewcommand{\rhedBestMeN}{0}  % best-day union: month ends
\providecommand{\rhedBestOpexN}{}\renewcommand{\rhedBestOpexN}{0}  % best-day union: monthly expirations
\providecommand{\rhedWorstOpexN}{}\renewcommand{\rhedWorstOpexN}{0}  % worst-day union: monthly expirations
\providecommand{\rhedBestNoTagN}{}\renewcommand{\rhedBestNoTagN}{8}  % best-day union: days without any event tag (calendar as covered)
\providecommand{\rhedBestAfterRelN}{}\renewcommand{\rhedBestAfterRelN}{5}  % best-day union: days after the release calendar ends (dagger)
\providecommand{\rhedWorstAfterRelN}{}\renewcommand{\rhedWorstAfterRelN}{4}  % worst-day union: days after the release calendar ends
\providecommand{\rhedDShortAllMin}{}\renewcommand{\rhedDShortAllMin}{1.21}  % lowest Delta SR mid vs the short straddle, all days across the six full-set pooled forecasts, 614 days
\providecommand{\rhedDShortAllMax}{}\renewcommand{\rhedDShortAllMax}{1.56}  % highest Delta SR mid vs the short straddle, all days across the six full-set pooled forecasts, 614 days
\providecommand{\rhedDShortExBestMin}{}\renewcommand{\rhedDShortExBestMin}{\ensuremath{-}0.54}  % lowest Delta SR mid vs the short straddle without the forecast's ten best days across the six full-set pooled forecasts, 614 days
\providecommand{\rhedDShortExBestMax}{}\renewcommand{\rhedDShortExBestMax}{\ensuremath{-}0.12}  % highest Delta SR mid vs the short straddle without the forecast's ten best days across the six full-set pooled forecasts, 614 days
\providecommand{\rhedDShortExBothMin}{}\renewcommand{\rhedDShortExBothMin}{\ensuremath{-}0.64}  % lowest ... without its ten best and ten worst days across the six full-set pooled forecasts, 614 days
\providecommand{\rhedDShortExBothMax}{}\renewcommand{\rhedDShortExBothMax}{\ensuremath{-}0.14}  % highest ... without its ten best and ten worst days across the six full-set pooled forecasts, 614 days
\providecommand{\rhedBestShareMin}{}\renewcommand{\rhedBestShareMin}{107}  % lowest % of the summed gain over the short straddle earned on the ten best days across the six full-set pooled forecasts, 614 days
\providecommand{\rhedBestShareMax}{}\renewcommand{\rhedBestShareMax}{138}  % highest % of the summed gain over the short straddle earned on the ten best days across the six full-set pooled forecasts, 614 days
\providecommand{\rhedDHarAllMin}{}\renewcommand{\rhedDHarAllMin}{0.56}  % lowest Delta SR mid vs the HAR twin, all days across the six full-set pooled forecasts, 614 days
\providecommand{\rhedDHarAllMax}{}\renewcommand{\rhedDHarAllMax}{1.56}  % highest Delta SR mid vs the HAR twin, all days across the six full-set pooled forecasts, 614 days
\providecommand{\rhedDHarExBestMin}{}\renewcommand{\rhedDHarExBestMin}{0.30}  % lowest Delta SR mid vs the HAR twin without the forecast's ten best days across the six full-set pooled forecasts, 614 days
\providecommand{\rhedDHarExBestMax}{}\renewcommand{\rhedDHarExBestMax}{0.83}  % highest Delta SR mid vs the HAR twin without the forecast's ten best days across the six full-set pooled forecasts, 614 days
\providecommand{\rhedDHarExBothMin}{}\renewcommand{\rhedDHarExBothMin}{0.36}  % lowest ... without its ten best and ten worst days across the six full-set pooled forecasts, 614 days
\providecommand{\rhedDHarExBothMax}{}\renewcommand{\rhedDHarExBothMax}{0.97}  % highest ... without its ten best and ten worst days across the six full-set pooled forecasts, 614 days
\providecommand{\rhedHarBestShareMin}{}\renewcommand{\rhedHarBestShareMin}{47}  % lowest % of the summed gain over the HAR twin earned on the ten best days across the six full-set pooled forecasts, 614 days
\providecommand{\rhedHarBestShareMax}{}\renewcommand{\rhedHarBestShareMax}{61}  % highest % of the summed gain over the HAR twin earned on the ten best days across the six full-set pooled forecasts, 614 days
\providecommand{\rhedDShortExBestNeg}{}\renewcommand{\rhedDShortExBestNeg}{6}  % full-set forecasts whose Delta SR vs short is negative without the ten best days
\providecommand{\rhedDHarExBestNeg}{}\renewcommand{\rhedDHarExBestNeg}{0}  % full-set forecasts whose Delta SR vs the HAR twin is negative without the ten best days
\providecommand{\rhedFomcN}{}\renewcommand{\rhedFomcN}{21}  % FOMC days among the 614 post-warmup days
\providecommand{\rhedFomcR}{}\renewcommand{\rhedFomcR}{25}  % mean mid straddle return R on FOMC days, %
\providecommand{\rhedFomcRT}{}\renewcommand{\rhedFomcRT}{1.02}  % ... its t (i.i.d.; the days are weeks apart)
\providecommand{\rhedFomcLong}{}\renewcommand{\rhedFomcLong}{25}  % % of FOMC days long, mean of the six full-set forecasts
\providecommand{\rhedFomcStrat}{}\renewcommand{\rhedFomcStrat}{\ensuremath{-}6}  % mean mid strategy return on FOMC days, %, mean of the six full-set
\providecommand{\rhedFomcOtherR}{}\renewcommand{\rhedFomcOtherR}{\ensuremath{-}3}  % mean straddle R on the other (covered) days, %
\providecommand{\rhedFomcOtherStrat}{}\renewcommand{\rhedFomcOtherStrat}{13}  % mean strategy return on the other (covered) days, %
\providecommand{\rhedFomcRVIV}{}\renewcommand{\rhedFomcRVIV}{0.76}  % mean RV_t / IV_t^2 on FOMC days
\providecommand{\rhedFomcMoveSqIV}{}\renewcommand{\rhedFomcMoveSqIV}{1.43}  % mean squared S&P 500 half-hour log return / IV_t^2 on FOMC days
\providecommand{\rhedFomcShare}{}\renewcommand{\rhedFomcShare}{\ensuremath{-}2}  % % of the summed mid return earned on FOMC days, mean of the six full-set forecasts
\providecommand{\rhedCpiN}{}\renewcommand{\rhedCpiN}{25}  % CPI days among the 614 post-warmup days
\providecommand{\rhedCpiR}{}\renewcommand{\rhedCpiR}{\ensuremath{-}12}  % mean mid straddle return R on CPI days, %
\providecommand{\rhedCpiRT}{}\renewcommand{\rhedCpiRT}{\ensuremath{-}0.55}  % ... its t (i.i.d.; the days are weeks apart)
\providecommand{\rhedCpiLong}{}\renewcommand{\rhedCpiLong}{16}  % % of CPI days long, mean of the six full-set forecasts
\providecommand{\rhedCpiStrat}{}\renewcommand{\rhedCpiStrat}{22}  % mean mid strategy return on CPI days, %, mean of the six full-set
\providecommand{\rhedCpiOtherR}{}\renewcommand{\rhedCpiOtherR}{\ensuremath{-}6}  % mean straddle R on the other (covered) days, %
\providecommand{\rhedCpiOtherStrat}{}\renewcommand{\rhedCpiOtherStrat}{11}  % mean strategy return on the other (covered) days, %
\providecommand{\rhedCpiRVIV}{}\renewcommand{\rhedCpiRVIV}{0.80}  % mean RV_t / IV_t^2 on CPI days
\providecommand{\rhedCpiMoveSqIV}{}\renewcommand{\rhedCpiMoveSqIV}{0.86}  % mean squared S&P 500 half-hour log return / IV_t^2 on CPI days
\providecommand{\rhedCpiShare}{}\renewcommand{\rhedCpiShare}{7}  % % of the summed mid return earned on CPI days, mean of the six full-set forecasts
\providecommand{\rhedNfpN}{}\renewcommand{\rhedNfpN}{27}  % NFP days among the 614 post-warmup days
\providecommand{\rhedNfpR}{}\renewcommand{\rhedNfpR}{1}  % mean mid straddle return R on NFP days, %
\providecommand{\rhedNfpRT}{}\renewcommand{\rhedNfpRT}{0.06}  % ... its t (i.i.d.; the days are weeks apart)
\providecommand{\rhedNfpLong}{}\renewcommand{\rhedNfpLong}{12}  % % of NFP days long, mean of the six full-set forecasts
\providecommand{\rhedNfpStrat}{}\renewcommand{\rhedNfpStrat}{20}  % mean mid strategy return on NFP days, %, mean of the six full-set
\providecommand{\rhedNfpOtherR}{}\renewcommand{\rhedNfpOtherR}{\ensuremath{-}6}  % mean straddle R on the other (covered) days, %
\providecommand{\rhedNfpOtherStrat}{}\renewcommand{\rhedNfpOtherStrat}{11}  % mean strategy return on the other (covered) days, %
\providecommand{\rhedNfpRVIV}{}\renewcommand{\rhedNfpRVIV}{0.73}  % mean RV_t / IV_t^2 on NFP days
\providecommand{\rhedNfpMoveSqIV}{}\renewcommand{\rhedNfpMoveSqIV}{0.93}  % mean squared S&P 500 half-hour log return / IV_t^2 on NFP days
\providecommand{\rhedNfpShare}{}\renewcommand{\rhedNfpShare}{7}  % % of the summed mid return earned on NFP days, mean of the six full-set forecasts
\providecommand{\rhedOpexN}{}\renewcommand{\rhedOpexN}{33}  % OpEx days among the 614 post-warmup days
\providecommand{\rhedOpexR}{}\renewcommand{\rhedOpexR}{\ensuremath{-}45}  % mean mid straddle return R on OpEx days, %
\providecommand{\rhedOpexRT}{}\renewcommand{\rhedOpexRT}{\ensuremath{-}4.01}  % ... its t (i.i.d.; the days are weeks apart)
\providecommand{\rhedOpexLong}{}\renewcommand{\rhedOpexLong}{9}  % % of OpEx days long, mean of the six full-set forecasts
\providecommand{\rhedOpexStrat}{}\renewcommand{\rhedOpexStrat}{36}  % mean mid strategy return on OpEx days, %, mean of the six full-set
\providecommand{\rhedOpexOtherR}{}\renewcommand{\rhedOpexOtherR}{0}  % mean straddle R on the other (covered) days, %
\providecommand{\rhedOpexOtherStrat}{}\renewcommand{\rhedOpexOtherStrat}{11}  % mean strategy return on the other (covered) days, %
\providecommand{\rhedOpexRVIV}{}\renewcommand{\rhedOpexRVIV}{0.83}  % mean RV_t / IV_t^2 on OpEx days
\providecommand{\rhedOpexMoveSqIV}{}\renewcommand{\rhedOpexMoveSqIV}{0.40}  % mean squared S&P 500 half-hour log return / IV_t^2 on OpEx days
\providecommand{\rhedOpexShare}{}\renewcommand{\rhedOpexShare}{15}  % % of the summed mid return earned on OpEx days, mean of the six full-set forecasts
\providecommand{\rhedTwN}{}\renewcommand{\rhedTwN}{11}  % TW days among the 614 post-warmup days
\providecommand{\rhedTwR}{}\renewcommand{\rhedTwR}{\ensuremath{-}52}  % mean mid straddle return R on TW days, %
\providecommand{\rhedTwRT}{}\renewcommand{\rhedTwRT}{\ensuremath{-}3.98}  % ... its t (i.i.d.; the days are weeks apart)
\providecommand{\rhedTwLong}{}\renewcommand{\rhedTwLong}{8}  % % of TW days long, mean of the six full-set forecasts
\providecommand{\rhedTwStrat}{}\renewcommand{\rhedTwStrat}{37}  % mean mid strategy return on TW days, %, mean of the six full-set
\providecommand{\rhedTwOtherR}{}\renewcommand{\rhedTwOtherR}{\ensuremath{-}2}  % mean straddle R on the other (covered) days, %
\providecommand{\rhedTwOtherStrat}{}\renewcommand{\rhedTwOtherStrat}{12}  % mean strategy return on the other (covered) days, %
\providecommand{\rhedTwRVIV}{}\renewcommand{\rhedTwRVIV}{0.89}  % mean RV_t / IV_t^2 on TW days
\providecommand{\rhedTwMoveSqIV}{}\renewcommand{\rhedTwMoveSqIV}{0.31}  % mean squared S&P 500 half-hour log return / IV_t^2 on TW days
\providecommand{\rhedTwShare}{}\renewcommand{\rhedTwShare}{5}  % % of the summed mid return earned on TW days, mean of the six full-set forecasts
\providecommand{\rhedMeN}{}\renewcommand{\rhedMeN}{33}  % ME days among the 614 post-warmup days
\providecommand{\rhedMeR}{}\renewcommand{\rhedMeR}{70}  % mean mid straddle return R on ME days, %
\providecommand{\rhedMeRT}{}\renewcommand{\rhedMeRT}{2.79}  % ... its t (i.i.d.; the days are weeks apart)
\providecommand{\rhedMeLong}{}\renewcommand{\rhedMeLong}{6}  % % of ME days long, mean of the six full-set forecasts
\providecommand{\rhedMeStrat}{}\renewcommand{\rhedMeStrat}{\ensuremath{-}58}  % mean mid strategy return on ME days, %, mean of the six full-set
\providecommand{\rhedMeOtherR}{}\renewcommand{\rhedMeOtherR}{\ensuremath{-}7}  % mean straddle R on the other (covered) days, %
\providecommand{\rhedMeOtherStrat}{}\renewcommand{\rhedMeOtherStrat}{17}  % mean strategy return on the other (covered) days, %
\providecommand{\rhedMeRVIV}{}\renewcommand{\rhedMeRVIV}{1.04}  % mean RV_t / IV_t^2 on ME days
\providecommand{\rhedMeMoveSqIV}{}\renewcommand{\rhedMeMoveSqIV}{1.93}  % mean squared S&P 500 half-hour log return / IV_t^2 on ME days
\providecommand{\rhedMeShare}{}\renewcommand{\rhedMeShare}{\ensuremath{-}24}  % % of the summed mid return earned on ME days, mean of the six full-set forecasts
\providecommand{\rhedRelEndText}{}\renewcommand{\rhedRelEndText}{26~October~2023}  % earliest last flag of the six macro-release series in data/releases.parquet
\providecommand{\rhedRelFomcEndText}{}\renewcommand{\rhedRelFomcEndText}{1~November~2023}  % last FOMC flag in data/releases.parquet
\providecommand{\rhedFomcSuppN}{}\renewcommand{\rhedFomcSuppN}{3}  % FOMC statement days of the sample taken from live/close_signal/state/fomc_statement_dates.csv
\providecommand{\rhedRelAfterN}{}\renewcommand{\rhedRelAfterN}{126}  % post-warmup days after the release calendar ends
\providecommand{\rhedWorthN}{}\renewcommand{\rhedWorthN}{137}  % post-warmup days with R_t = -1 (= rheWorthlessDays)
\providecommand{\rhedWorthPct}{}\renewcommand{\rhedWorthPct}{22.3}  % ... % of the 614 days
\providecommand{\rhedWorthExpIV}{}\renewcommand{\rhedWorthExpIV}{128}  % expected R = -1 days under a normal half-hour log move with volatility IV_t
\providecommand{\rhedWorthExpIVPct}{}\renewcommand{\rhedWorthExpIVPct}{20.9}  % ... as % of the 614 days
\providecommand{\rhedWorthZ}{}\renewcommand{\rhedWorthZ}{0.90}  % (observed - expected) / sd of the count under that normal (independent days)
\providecommand{\rhedWorthExpEmp}{}\renewcommand{\rhedWorthExpEmp}{132}  % expected R = -1 days under the empirical distribution of the earlier days' moves standardized by IV
\providecommand{\rhedWorthExpRaw}{}\renewcommand{\rhedWorthExpRaw}{118}  % expected R = -1 days under the earlier days' raw half-hour moves
\providecommand{\rhedGapFiveN}{}\renewcommand{\rhedGapFiveN}{611}  % post-warmup days with K_C - K_P = 5 points
\providecommand{\rhedGapOtherN}{}\renewcommand{\rhedGapOtherN}{3}  % post-warmup days with another strike gap
\providecommand{\rhedMoveSdPts}{}\renewcommand{\rhedMoveSdPts}{9.6}  % median one-sd implied half-hour move, index points
\providecommand{\rhedStrikeDist}{}\renewcommand{\rhedStrikeDist}{1.23}  % mean distance of the close to the nearest 5-point strike, points
\providecommand{\rhedStrikeDistUnif}{}\renewcommand{\rhedStrikeDistUnif}{1.25}  % ... its mean under a uniform close
\providecommand{\rhedKsD}{}\renewcommand{\rhedKsD}{0.026}  % Kolmogorov-Smirnov D of the close's position within its 5-point strike interval against a uniform
\providecommand{\rhedKsCV}{}\renewcommand{\rhedKsCV}{0.055}  % ... asymptotic 5% critical value
\providecommand{\rhedZSd}{}\renewcommand{\rhedZSd}{1.00}  % sd of the half-hour log move standardized by IV_t, 614 days
\providecommand{\rhedZKurt}{}\renewcommand{\rhedZKurt}{2.0}  % excess kurtosis of that standardized move
\providecommand{\rhedWorthShortMin}{}\renewcommand{\rhedWorthShortMin}{110}  % lowest R = -1 days on which the forecast is short across the six full-set pooled forecasts, 614 days
\providecommand{\rhedWorthShortMax}{}\renewcommand{\rhedWorthShortMax}{115}  % highest R = -1 days on which the forecast is short across the six full-set pooled forecasts, 614 days
\providecommand{\rhedWorthOpexN}{}\renewcommand{\rhedWorthOpexN}{10}  % R = -1 days among the monthly-expiration days
\providecommand{\rhedWorthOpexExp}{}\renewcommand{\rhedWorthOpexExp}{6.1}  % ... expected under the normal with IV_t
\providecommand{\rhedWorthOtherPct}{}\renewcommand{\rhedWorthOtherPct}{21.9}  % % of the other days with R = -1
\providecommand{\rhedWorthOtherExpPct}{}\renewcommand{\rhedWorthOtherExpPct}{21.0}  % ... expected under the normal with IV_t, %
\providecommand{\rhedWorthOpexZ}{}\renewcommand{\rhedWorthOpexZ}{1.77}  % (observed - expected) / sd of the count on the monthly-expiration days
\providecommand{\rhedRegAlphaAMin}{}\renewcommand{\rhedRegAlphaAMin}{11.0}  % lowest alpha, (a) on the short straddle, mid, % per day across the six full-set pooled forecasts, 614 days
\providecommand{\rhedRegAlphaAMax}{}\renewcommand{\rhedRegAlphaAMax}{13.4}  % highest alpha, (a) on the short straddle, mid, % per day across the six full-set pooled forecasts, 614 days
\providecommand{\rhedRegAlphaAXMin}{}\renewcommand{\rhedRegAlphaAXMin}{8.3}  % lowest alpha, (a) on the short straddle, bid/ask, % per day across the six full-set pooled forecasts, 614 days
\providecommand{\rhedRegAlphaAXMax}{}\renewcommand{\rhedRegAlphaAXMax}{10.7}  % highest alpha, (a) on the short straddle, bid/ask, % per day across the six full-set pooled forecasts, 614 days
\providecommand{\rhedRegBetaShortMin}{}\renewcommand{\rhedRegBetaShortMin}{0.23}  % lowest beta on the short-straddle return, (a), mid across the six full-set pooled forecasts, 614 days
\providecommand{\rhedRegBetaShortMax}{}\renewcommand{\rhedRegBetaShortMax}{0.28}  % highest beta on the short-straddle return, (a), mid across the six full-set pooled forecasts, 614 days
\providecommand{\rhedRegAlphaAPMin}{}\renewcommand{\rhedRegAlphaAPMin}{0.005}  % lowest fixed-b p of alpha, (a), mid across the six full-set pooled forecasts, 614 days
\providecommand{\rhedRegAlphaAPMax}{}\renewcommand{\rhedRegAlphaAPMax}{0.025}  % highest fixed-b p of alpha, (a), mid across the six full-set pooled forecasts, 614 days
\providecommand{\rhedRegAlphaAXPMin}{}\renewcommand{\rhedRegAlphaAXPMin}{0.021}  % lowest fixed-b p of alpha, (a), bid/ask across the six full-set pooled forecasts, 614 days
\providecommand{\rhedRegAlphaAXPMax}{}\renewcommand{\rhedRegAlphaAXPMax}{0.077}  % highest fixed-b p of alpha, (a), bid/ask across the six full-set pooled forecasts, 614 days
\providecommand{\rhedRegBetaShortTMin}{}\renewcommand{\rhedRegBetaShortTMin}{1.22}  % lowest t of the beta on the short straddle, (a), mid across the six full-set pooled forecasts, 614 days
\providecommand{\rhedRegBetaShortTMax}{}\renewcommand{\rhedRegBetaShortTMax}{1.42}  % highest t of the beta on the short straddle, (a), mid across the six full-set pooled forecasts, 614 days
\providecommand{\rhedRegRsqAMin}{}\renewcommand{\rhedRegRsqAMin}{0.06}  % lowest R^2 of (a), mid across the six full-set pooled forecasts, 614 days
\providecommand{\rhedRegRsqAMax}{}\renewcommand{\rhedRegRsqAMax}{0.08}  % highest R^2 of (a), mid across the six full-set pooled forecasts, 614 days
\providecommand{\rhedRegBetaAbsMin}{}\renewcommand{\rhedRegBetaAbsMin}{\ensuremath{-}1.87}  % lowest beta on |r30| (bp), (b), mid, % of premium per bp across the six full-set pooled forecasts, 614 days
\providecommand{\rhedRegBetaAbsMax}{}\renewcommand{\rhedRegBetaAbsMax}{\ensuremath{-}1.40}  % highest beta on |r30| (bp), (b), mid, % of premium per bp across the six full-set pooled forecasts, 614 days
\providecommand{\rhedRegBetaAbsTMin}{}\renewcommand{\rhedRegBetaAbsTMin}{\ensuremath{-}3.49}  % lowest t of the beta on |r30|, (b), mid across the six full-set pooled forecasts, 614 days
\providecommand{\rhedRegBetaAbsTMax}{}\renewcommand{\rhedRegBetaAbsTMax}{\ensuremath{-}2.31}  % highest t of the beta on |r30|, (b), mid across the six full-set pooled forecasts, 614 days
\providecommand{\rhedRegBetaMoveMin}{}\renewcommand{\rhedRegBetaMoveMin}{\ensuremath{-}0.01}  % lowest beta on r30 (bp), (b), mid across the six full-set pooled forecasts, 614 days
\providecommand{\rhedRegBetaMoveMax}{}\renewcommand{\rhedRegBetaMoveMax}{0.14}  % highest beta on r30 (bp), (b), mid across the six full-set pooled forecasts, 614 days
\providecommand{\rhedRegMeanMin}{}\renewcommand{\rhedRegMeanMin}{11.5}  % lowest mean return, mid, % per day across the six full-set pooled forecasts, 614 days
\providecommand{\rhedRegMeanMax}{}\renewcommand{\rhedRegMeanMax}{14.1}  % highest mean return, mid, % per day across the six full-set pooled forecasts, 614 days
\providecommand{\rhedRegMeanXMin}{}\renewcommand{\rhedRegMeanXMin}{8.9}  % lowest mean return, bid/ask, % per day across the six full-set pooled forecasts, 614 days
\providecommand{\rhedRegMeanXMax}{}\renewcommand{\rhedRegMeanXMax}{11.4}  % highest mean return, bid/ask, % per day across the six full-set pooled forecasts, 614 days
\providecommand{\rhedRegRsqCMin}{}\renewcommand{\rhedRegRsqCMin}{0.00}  % lowest R^2 of (c), mid across the six full-set pooled forecasts, 614 days
\providecommand{\rhedRegRsqCMax}{}\renewcommand{\rhedRegRsqCMax}{0.00}  % highest R^2 of (c), mid across the six full-set pooled forecasts, 614 days
\providecommand{\rhedRegCPMaxMin}{}\renewcommand{\rhedRegCPMaxMin}{0.74}  % lowest larger fixed-b p of the two betas of (c), mid across the six full-set pooled forecasts, 614 days
\providecommand{\rhedRegCPMaxMax}{}\renewcommand{\rhedRegCPMaxMax}{0.97}  % highest larger fixed-b p of the two betas of (c), mid across the six full-set pooled forecasts, 614 days
\providecommand{\rhedRegCPMinMin}{}\renewcommand{\rhedRegCPMinMin}{0.56}  % lowest smaller fixed-b p of the two betas of (c), mid across the six full-set pooled forecasts, 614 days
\providecommand{\rhedRegCPMinMax}{}\renewcommand{\rhedRegCPMinMax}{0.82}  % highest smaller fixed-b p of the two betas of (c), mid across the six full-set pooled forecasts, 614 days
\providecommand{\rhedRegAlphaDMin}{}\renewcommand{\rhedRegAlphaDMin}{11.4}  % lowest alpha, (d) joint, mid, % per day across the six full-set pooled forecasts, 614 days
\providecommand{\rhedRegAlphaDMax}{}\renewcommand{\rhedRegAlphaDMax}{13.6}  % highest alpha, (d) joint, mid, % per day across the six full-set pooled forecasts, 614 days
\providecommand{\rhedRegAlphaDXMin}{}\renewcommand{\rhedRegAlphaDXMin}{8.7}  % lowest alpha, (d) joint, bid/ask, % per day across the six full-set pooled forecasts, 614 days
\providecommand{\rhedRegAlphaDXMax}{}\renewcommand{\rhedRegAlphaDXMax}{10.9}  % highest alpha, (d) joint, bid/ask, % per day across the six full-set pooled forecasts, 614 days
\providecommand{\rhedRegAlphaDPMin}{}\renewcommand{\rhedRegAlphaDPMin}{0.004}  % lowest fixed-b p of alpha, (d), mid across the six full-set pooled forecasts, 614 days
\providecommand{\rhedRegAlphaDPMax}{}\renewcommand{\rhedRegAlphaDPMax}{0.026}  % highest fixed-b p of alpha, (d), mid across the six full-set pooled forecasts, 614 days
\providecommand{\rhedRegAlphaDXPMin}{}\renewcommand{\rhedRegAlphaDXPMin}{0.017}  % lowest fixed-b p of alpha, (d), bid/ask across the six full-set pooled forecasts, 614 days
\providecommand{\rhedRegAlphaDXPMax}{}\renewcommand{\rhedRegAlphaDXPMax}{0.076}  % highest fixed-b p of alpha, (d), bid/ask across the six full-set pooled forecasts, 614 days
\providecommand{\rhedRegAvgAlphaM}{}\renewcommand{\rhedRegAvgAlphaM}{12.7}  % alpha of the mean of the six full-set returns, (m), mid, % per day
\providecommand{\rhedRegAvgAlphaMT}{}\renewcommand{\rhedRegAvgAlphaMT}{3.48}  % ... its t
\providecommand{\rhedRegAvgAlphaMP}{}\renewcommand{\rhedRegAvgAlphaMP}{0.002}  % ... its fixed-b p
\providecommand{\rhedRegAvgAlphaA}{}\renewcommand{\rhedRegAvgAlphaA}{12.1}  % alpha of the mean of the six full-set returns, (a), mid, % per day
\providecommand{\rhedRegAvgAlphaAT}{}\renewcommand{\rhedRegAvgAlphaAT}{2.83}  % ... its t
\providecommand{\rhedRegAvgAlphaAP}{}\renewcommand{\rhedRegAvgAlphaAP}{0.0100}  % ... its fixed-b p
\providecommand{\rhedRegAvgAlphaD}{}\renewcommand{\rhedRegAvgAlphaD}{12.5}  % alpha of the mean of the six full-set returns, (d), mid, % per day
\providecommand{\rhedRegAvgAlphaDT}{}\renewcommand{\rhedRegAvgAlphaDT}{2.86}  % ... its t
\providecommand{\rhedRegAvgAlphaDP}{}\renewcommand{\rhedRegAvgAlphaDP}{0.009}  % ... its fixed-b p
\providecommand{\rhedRegAvgAlphaMX}{}\renewcommand{\rhedRegAvgAlphaMX}{10.1}  % alpha of the mean of the six full-set returns, (m), ba, % per day
\providecommand{\rhedRegAvgAlphaMXT}{}\renewcommand{\rhedRegAvgAlphaMXT}{2.78}  % ... its t
\providecommand{\rhedRegAvgAlphaMXP}{}\renewcommand{\rhedRegAvgAlphaMXP}{0.011}  % ... its fixed-b p
\providecommand{\rhedRegAvgAlphaAX}{}\renewcommand{\rhedRegAvgAlphaAX}{9.4}  % alpha of the mean of the six full-set returns, (a), ba, % per day
\providecommand{\rhedRegAvgAlphaAXT}{}\renewcommand{\rhedRegAvgAlphaAXT}{2.21}  % ... its t
\providecommand{\rhedRegAvgAlphaAXP}{}\renewcommand{\rhedRegAvgAlphaAXP}{0.040}  % ... its fixed-b p
\providecommand{\rhedRegAvgAlphaDX}{}\renewcommand{\rhedRegAvgAlphaDX}{9.8}  % alpha of the mean of the six full-set returns, (d), ba, % per day
\providecommand{\rhedRegAvgAlphaDXT}{}\renewcommand{\rhedRegAvgAlphaDXT}{2.26}  % ... its t
\providecommand{\rhedRegAvgAlphaDXP}{}\renewcommand{\rhedRegAvgAlphaDXP}{0.036}  % ... its fixed-b p
\providecommand{\rhedRegAvgBetaShort}{}\renewcommand{\rhedRegAvgBetaShort}{0.26}  % beta of the mean full-set return on the short straddle, (a), mid
\providecommand{\rhedRegAvgBetaShortT}{}\renewcommand{\rhedRegAvgBetaShortT}{1.34}  % ... its t
\providecommand{\rhedRegAvgBetaShortP}{}\renewcommand{\rhedRegAvgBetaShortP}{0.20}  % ... its fixed-b p
\providecommand{\rhedRegAvgRsqA}{}\renewcommand{\rhedRegAvgRsqA}{0.07}  % R^2 of (a) for the mean full-set return
\providecommand{\rhedRegAvgBetaAbs}{}\renewcommand{\rhedRegAvgBetaAbs}{\ensuremath{-}1.68}  % beta of the mean full-set return on |r30| (bp), (b), mid
\providecommand{\rhedRegAvgBetaAbsT}{}\renewcommand{\rhedRegAvgBetaAbsT}{\ensuremath{-}3.26}  % ... its t
\providecommand{\rhedRegShortBetaAbs}{}\renewcommand{\rhedRegShortBetaAbs}{\ensuremath{-}3.96}  % beta of the short straddle on |r30| (bp), (b), mid
\providecommand{\rhedRegShortRsqB}{}\renewcommand{\rhedRegShortRsqB}{0.48}  % R^2 of (b) for the short straddle
\providecommand{\rhedRegAbsRatio}{}\renewcommand{\rhedRegAbsRatio}{43}  % beta on |r30| of the mean full-set return as % of the short straddle's
\providecommand{\rhedRegAvgBetaAbsJ}{}\renewcommand{\rhedRegAvgBetaAbsJ}{\ensuremath{-}1.30}  % beta of the mean full-set return on |r30|, (d) joint, mid
\providecommand{\rhedRegAvgBetaAbsJT}{}\renewcommand{\rhedRegAvgBetaAbsJT}{\ensuremath{-}1.90}  % ... its t
\providecommand{\rhedHarRegAlphaAMin}{}\renewcommand{\rhedHarRegAlphaAMin}{1.1}  % lowest alpha, (a), mid, % per day across the six HAR-benchmark pooled forecasts, 614 days
\providecommand{\rhedHarRegAlphaAMax}{}\renewcommand{\rhedHarRegAlphaAMax}{7.0}  % highest alpha, (a), mid, % per day across the six HAR-benchmark pooled forecasts, 614 days
\providecommand{\rhedHarRegAlphaAPMin}{}\renewcommand{\rhedHarRegAlphaAPMin}{0.082}  % lowest fixed-b p of alpha, (a), mid across the six HAR-benchmark pooled forecasts, 614 days
\providecommand{\rhedHarRegAlphaAPMax}{}\renewcommand{\rhedHarRegAlphaAPMax}{0.745}  % highest fixed-b p of alpha, (a), mid across the six HAR-benchmark pooled forecasts, 614 days
\providecommand{\rhedOlsRegAlphaA}{}\renewcommand{\rhedOlsRegAlphaA}{8.1}  % alpha, (a), mid, closing-interval HAR forecast, % per day
\providecommand{\rhedOlsRegAlphaAP}{}\renewcommand{\rhedOlsRegAlphaAP}{0.051}  % ... its fixed-b p
\providecommand{\rhedRegAlphaAExBestMin}{}\renewcommand{\rhedRegAlphaAExBestMin}{\ensuremath{-}0.5}  % lowest alpha, (a), mid, without the forecast's ten best days, % per day across the six full-set pooled forecasts, 614 days
\providecommand{\rhedRegAlphaAExBestMax}{}\renewcommand{\rhedRegAlphaAExBestMax}{2.1}  % highest alpha, (a), mid, without the forecast's ten best days, % per day across the six full-set pooled forecasts, 614 days
\providecommand{\rhedRegAlphaAExBestPMin}{}\renewcommand{\rhedRegAlphaAExBestPMin}{0.43}  % lowest ... its fixed-b p across the six full-set pooled forecasts, 614 days
\providecommand{\rhedRegAlphaAExBestPMax}{}\renewcommand{\rhedRegAlphaAExBestPMax}{0.97}  % highest ... its fixed-b p across the six full-set pooled forecasts, 614 days
\fi
\begin{table}[tbp]
\centering
\small
\setlength{\tabcolsep}{3.5pt}
\begin{threeparttable}
\caption{The best and worst days of the sign strategy}
\label{tab:rh_econ_bestdays}
\begin{tabular}{@{}lrrrrrrrl@{}}
\toprule
 & In & Long & \multicolumn{2}{c}{Return (\%)} & & \multicolumn{2}{c}{S\&P~500 (\%)} & \\
\cmidrule(lr){4-5}\cmidrule(lr){7-8}
Date & ten & (of 6) & Mid & Bid/ask & $\RV_t/\mathrm{IV}_t^2$ & 15:30--16:00 & Daily & Events \\
\midrule
\multicolumn{9}{@{}l}{\textit{Panel A: Best days}} \\
\addlinespace[2pt]
2021-11-22 & 6 & 6 & $1{,}032$ & $974$ & $2.79$ & $-0.83$ & $-0.32$ & -- \\
2023-03-14 & 6 & 6 & $306$ & $299$ & $1.01$ & $0.88$ & $1.63$ & CPI \\
2023-03-22 & 1 & 1 & $355$ & $347$ & $1.34$ & $-1.70$ & $-1.66$ & FOMC \\
2023-06-26 & 6 & 6 & $299$ & $292$ & $1.63$ & $-0.31$ & $-0.45$ & -- \\
2023-07-07 & 5 & 5 & $418$ & $402$ & $1.64$ & $-0.41$ & $-0.29$ & NFP \\
2023-09-22 & 6 & 6 & $333$ & $324$ & $1.49$ & $-0.52$ & $-0.23$ & -- \\
2024-01-02 & 6 & 6 & $272$ & $261$ & $1.27$ & $0.34$ & $-0.57$ & --$^{\dagger}$ \\
2024-01-17 & 6 & 6 & $269$ & $260$ & $1.34$ & $0.31$ & $-0.56$ & --$^{\dagger}$ \\
2024-02-21 & 6 & 6 & $425$ & $410$ & $1.21$ & $0.56$ & $0.13$ & --$^{\dagger}$ \\
2024-03-04 & 6 & 6 & $218$ & $212$ & $1.67$ & $-0.23$ & $-0.12$ & --$^{\dagger}$ \\
2024-03-05 & 6 & 6 & $299$ & $293$ & $1.31$ & $0.41$ & $-1.02$ & --$^{\dagger}$ \\
\addlinespace[4pt]
\multicolumn{9}{@{}l}{\textit{Panel B: Worst days}} \\
\addlinespace[2pt]
2022-11-14 & 6 & 0 & $-327$ & $-337$ & $1.15$ & $-0.72$ & $-0.90$ & -- \\
2023-01-31 & 6 & 0 & $-322$ & $-333$ & $0.88$ & $0.66$ & $1.45$ & ME \\
2023-03-22 & 5 & 1 & $-355$ & $-363$ & $1.34$ & $-1.70$ & $-1.66$ & FOMC \\
2023-04-11 & 6 & 0 & $-325$ & $-334$ & $1.28$ & $-0.32$ & $0.00$ & -- \\
2023-05-16 & 6 & 0 & $-304$ & $-321$ & $1.90$ & $-0.36$ & $-0.64$ & RS \\
2023-07-07 & 1 & 5 & $-418$ & $-435$ & $1.64$ & $-0.41$ & $-0.29$ & NFP \\
2023-08-09 & 6 & 0 & $-374$ & $-387$ & $1.79$ & $-0.47$ & $-0.71$ & -- \\
2023-11-30 & 6 & 0 & $-542$ & $-561$ & $1.89$ & $0.51$ & $0.38$ & ME$^{\dagger}$ \\
2024-02-13 & 6 & 0 & $-385$ & $-393$ & $2.08$ & $0.54$ & $-1.38$ & --$^{\dagger}$ \\
2024-03-26 & 6 & 0 & $-353$ & $-369$ & $1.29$ & $-0.35$ & $-0.28$ & --$^{\dagger}$ \\
2024-04-30 & 6 & 0 & $-381$ & $-386$ & $1.88$ & $-0.53$ & $-1.59$ & ME$^{\dagger}$ \\
\bottomrule
\end{tabular}

\medskip
\begin{tabular}{@{}lrrrrrrr@{}}
\toprule
 & \multicolumn{4}{c}{$\Delta\SR$ against the short straddle} & \multicolumn{3}{c}{$\Delta\SR$ against the HAR twin} \\
\cmidrule(lr){2-5}\cmidrule(l){6-8}
Panel C & All days & No best & No best, & Best-day & All days & No best & No best, \\
 & & & no worst & share (\%) & & & no worst \\
\midrule
Ridge & $1.87$ & $-0.18$ & $-0.25$ & $127$ & $0.77$ & $0.50$ & $0.46$ \\
Lasso & $1.87$ & $-0.25$ & $-0.32$ & $120$ & $0.77$ & $0.43$ & $0.38$ \\
Elastic net & $2.07$ & $0.00$ & $0.01$ & $110$ & $0.97$ & $0.68$ & $0.72$ \\
LightGBM & $1.78$ & $-0.34$ & $-0.35$ & $138$ & $1.06$ & $0.78$ & $1.00$ \\
XGBoost & $1.83$ & $-0.26$ & $-0.24$ & $138$ & $0.81$ & $0.75$ & $1.09$ \\
Random forest & $2.10$ & $0.13$ & $0.13$ & $107$ & $1.98$ & $1.13$ & $1.31$ \\
\bottomrule
\end{tabular}
\begin{tablenotes}[flushleft]
\footnotesize
\item \textit{Notes:} Panels A and B: the union over the six pooled intraday forecasts on the full predictor set of
each forecast's \rhedNTop{} best and \rhedNTop{} worst post-warmup days by its mid-quote return. In ten: the number of
the six forecasts with the day among their ten; Long: the number of the six that hold a long straddle. Return: of the
position of the forecasts that count the day among their ten (a long straddle bought at the ask, a short one sold at
the bid, in the bid/ask column). $\RV_t/\mathrm{IV}_t^2$: realized
over implied variance of the closing half hour. S\&P~500: the log return from the 15:30 underlying print $S_t$ to the
close $S^{\mathrm{close}}_t$, and from the previous trading day's close. Events: FOMC statement, CPI, employment report
(NFP), PPI, GDP, retail sales (RS) and trade balance (TB) release days; monthly SPX expiration (OpEx; TW, the quarterly
triple witching); month end (ME). $^{\dagger}$: after \rhedRelEndText, when the release calendar ends (the FOMC
calendar is complete). Panel C: for each forecast, the Sharpe ratio gain $\Delta\SR$ (mid-quotes) over the unconditional
short straddle and over the same model class on the HAR benchmark, on all \rheSzDays{} post-warmup days, without the
forecast's ten best days, and without its ten best and ten worst days (dropped from both series); Best-day share: the
part of the summed return gain over the short straddle earned on the ten best days.
\end{tablenotes}
\end{threeparttable}
\end{table}
     % tab:rh_econ_bestdays
%   % AUTO-GENERATED by writeup/make_rehaul_econ_decomp.py -- do not edit; re-run the script.
% Table tab:rh_econ_riskreg: risk regressions of the sign strategy's daily return
% (writeup/make_rehaul_econ_decomp.py, statistic 6).
\ifdefined\rhedNDays\else% AUTO-GENERATED by writeup/make_rehaul_econ_decomp.py -- do not edit; re-run the script.
% \rhed... macros: where the straddle strategy's profits come from (Section 7 decomposition).
\providecommand{\rhedNDays}{}\renewcommand{\rhedNDays}{614}  % post-warmup option days (= rheSzDays)
\providecommand{\rhedNFull}{}\renewcommand{\rhedNFull}{6}  % pooled intraday forecasts on the full predictor set
\providecommand{\rhedNTop}{}\renewcommand{\rhedNTop}{10}  % best / worst days per forecast
\providecommand{\rhedUncRVAboveIV}{}\renewcommand{\rhedUncRVAboveIV}{28.2}  % % of the 614 days with RV_t > IV_t^2
\providecommand{\rhedUncRVIVMean}{}\renewcommand{\rhedUncRVIVMean}{0.87}  % mean RV_t / IV_t^2, 614 days
\providecommand{\rhedUncRVIVMed}{}\renewcommand{\rhedUncRVIVMed}{0.77}  % median RV_t / IV_t^2, 614 days
\providecommand{\rhedUncShortMean}{}\renewcommand{\rhedUncShortMean}{2.4}  % mean mid return of the unconditional short straddle, % of premium
\providecommand{\rhedUncShortHit}{}\renewcommand{\rhedUncShortHit}{71.8}  % % of days on which q = -1 matches sign(RV - IV^2) (always short)
\providecommand{\rhedLongNMin}{}\renewcommand{\rhedLongNMin}{117}  % lowest number of long days across the six full-set pooled forecasts, 614 days
\providecommand{\rhedLongNMax}{}\renewcommand{\rhedLongNMax}{146}  % highest number of long days across the six full-set pooled forecasts, 614 days
\providecommand{\rhedLongPctMin}{}\renewcommand{\rhedLongPctMin}{19}  % lowest % long days across the six full-set pooled forecasts, 614 days
\providecommand{\rhedLongPctMax}{}\renewcommand{\rhedLongPctMax}{24}  % highest % long days across the six full-set pooled forecasts, 614 days
\providecommand{\rhedHarLongNMin}{}\renewcommand{\rhedHarLongNMin}{150}  % lowest number of long days across the six HAR-benchmark pooled forecasts, 614 days
\providecommand{\rhedHarLongNMax}{}\renewcommand{\rhedHarLongNMax}{193}  % highest number of long days across the six HAR-benchmark pooled forecasts, 614 days
\providecommand{\rhedLongMidMin}{}\renewcommand{\rhedLongMidMin}{19}  % lowest mean mid return on long days, % across the six full-set pooled forecasts, 614 days
\providecommand{\rhedLongMidMax}{}\renewcommand{\rhedLongMidMax}{30}  % highest mean mid return on long days, % across the six full-set pooled forecasts, 614 days
\providecommand{\rhedLongBaMin}{}\renewcommand{\rhedLongBaMin}{16}  % lowest mean bid/ask return on long days, % across the six full-set pooled forecasts, 614 days
\providecommand{\rhedLongBaMax}{}\renewcommand{\rhedLongBaMax}{26}  % highest mean bid/ask return on long days, % across the six full-set pooled forecasts, 614 days
\providecommand{\rhedShortMidMin}{}\renewcommand{\rhedShortMidMin}{9}  % lowest mean mid return on short days, % across the six full-set pooled forecasts, 614 days
\providecommand{\rhedShortMidMax}{}\renewcommand{\rhedShortMidMax}{11}  % highest mean mid return on short days, % across the six full-set pooled forecasts, 614 days
\providecommand{\rhedShortBaMin}{}\renewcommand{\rhedShortBaMin}{7}  % lowest mean bid/ask return on short days, % across the six full-set pooled forecasts, 614 days
\providecommand{\rhedShortBaMax}{}\renewcommand{\rhedShortBaMax}{8}  % highest mean bid/ask return on short days, % across the six full-set pooled forecasts, 614 days
\providecommand{\rhedHarLongMidMin}{}\renewcommand{\rhedHarLongMidMin}{0}  % lowest mean mid return on long days, % across the six HAR-benchmark pooled forecasts, 614 days
\providecommand{\rhedHarLongMidMax}{}\renewcommand{\rhedHarLongMidMax}{10}  % highest mean mid return on long days, % across the six HAR-benchmark pooled forecasts, 614 days
\providecommand{\rhedHarShortMidMin}{}\renewcommand{\rhedHarShortMidMin}{3}  % lowest mean mid return on short days, % across the six HAR-benchmark pooled forecasts, 614 days
\providecommand{\rhedHarShortMidMax}{}\renewcommand{\rhedHarShortMidMax}{7}  % highest mean mid return on short days, % across the six HAR-benchmark pooled forecasts, 614 days
\providecommand{\rhedLongHitMin}{}\renewcommand{\rhedLongHitMin}{47}  % lowest % of long days with RV > IV^2 across the six full-set pooled forecasts, 614 days
\providecommand{\rhedLongHitMax}{}\renewcommand{\rhedLongHitMax}{53}  % highest % of long days with RV > IV^2 across the six full-set pooled forecasts, 614 days
\providecommand{\rhedHarLongHitMin}{}\renewcommand{\rhedHarLongHitMin}{46}  % lowest % of long days with RV > IV^2 across the six HAR-benchmark pooled forecasts, 614 days
\providecommand{\rhedHarLongHitMax}{}\renewcommand{\rhedHarLongHitMax}{48}  % highest % of long days with RV > IV^2 across the six HAR-benchmark pooled forecasts, 614 days
\providecommand{\rhedLongPosMin}{}\renewcommand{\rhedLongPosMin}{42}  % lowest % of long days with a positive return across the six full-set pooled forecasts, 614 days
\providecommand{\rhedLongPosMax}{}\renewcommand{\rhedLongPosMax}{45}  % highest % of long days with a positive return across the six full-set pooled forecasts, 614 days
\providecommand{\rhedHitMin}{}\renewcommand{\rhedHitMin}{70.7}  % lowest % of days on which q matches sign(RV - IV^2) across the six full-set pooled forecasts, 614 days
\providecommand{\rhedHitMax}{}\renewcommand{\rhedHitMax}{73.0}  % highest % of days on which q matches sign(RV - IV^2) across the six full-set pooled forecasts, 614 days
\providecommand{\rhedPosMin}{}\renewcommand{\rhedPosMin}{58.6}  % lowest % of days with a positive mid return across the six full-set pooled forecasts, 614 days
\providecommand{\rhedPosMax}{}\renewcommand{\rhedPosMax}{60.4}  % highest % of days with a positive mid return across the six full-set pooled forecasts, 614 days
\providecommand{\rhedLongShareMin}{}\renewcommand{\rhedLongShareMin}{39}  % lowest % of the summed mid return earned on long days across the six full-set pooled forecasts, 614 days
\providecommand{\rhedLongShareMax}{}\renewcommand{\rhedLongShareMax}{41}  % highest % of the summed mid return earned on long days across the six full-set pooled forecasts, 614 days
\providecommand{\rhedLongShareXMin}{}\renewcommand{\rhedLongShareXMin}{42}  % lowest % of the summed bid/ask return earned on long days across the six full-set pooled forecasts, 614 days
\providecommand{\rhedLongShareXMax}{}\renewcommand{\rhedLongShareXMax}{45}  % highest % of the summed bid/ask return earned on long days across the six full-set pooled forecasts, 614 days
\providecommand{\rhedMeanMin}{}\renewcommand{\rhedMeanMin}{11.5}  % lowest mean mid return, % across the six full-set pooled forecasts, 614 days
\providecommand{\rhedMeanMax}{}\renewcommand{\rhedMeanMax}{14.1}  % highest mean mid return, % across the six full-set pooled forecasts, 614 days
\providecommand{\rhedGainMin}{}\renewcommand{\rhedGainMin}{9.1}  % lowest mean gain over the short straddle (2 R 1(long)), % across the six full-set pooled forecasts, 614 days
\providecommand{\rhedGainMax}{}\renewcommand{\rhedGainMax}{11.6}  % highest mean gain over the short straddle (2 R 1(long)), % across the six full-set pooled forecasts, 614 days
\providecommand{\rhedGainShareMin}{}\renewcommand{\rhedGainShareMin}{79}  % lowest gain over short / mean return, % across the six full-set pooled forecasts, 614 days
\providecommand{\rhedGainShareMax}{}\renewcommand{\rhedGainShareMax}{83}  % highest gain over short / mean return, % across the six full-set pooled forecasts, 614 days
\providecommand{\rhedHarGainMin}{}\renewcommand{\rhedHarGainMin}{0.0}  % lowest mean gain over the short straddle, % across the six HAR-benchmark pooled forecasts, 614 days
\providecommand{\rhedHarGainMax}{}\renewcommand{\rhedHarGainMax}{5.0}  % highest mean gain over the short straddle, % across the six HAR-benchmark pooled forecasts, 614 days
\providecommand{\rhedOlsLongN}{}\renewcommand{\rhedOlsLongN}{198}  % long days of the closing-interval HAR forecast (sub_ols_baseline)
\providecommand{\rhedOlsLongMid}{}\renewcommand{\rhedOlsLongMid}{9}  % ... its mean mid return on long days, %
\providecommand{\rhedOlsShortMid}{}\renewcommand{\rhedOlsShortMid}{8}  % ... its mean mid return on short days, %
\providecommand{\rhedOlsLongHit}{}\renewcommand{\rhedOlsLongHit}{46}  % ... % of its long days with RV > IV^2
\providecommand{\rhedOlsSh}{}\renewcommand{\rhedOlsSh}{1.11}  % ... its Sharpe ratio, mid (= rhrOlsSharpe)
\providecommand{\rhedCalLevelMin}{}\renewcommand{\rhedCalLevelMin}{0.95}  % lowest mean forecast / mean RV across the six full-set pooled forecasts, 614 days
\providecommand{\rhedCalLevelMax}{}\renewcommand{\rhedCalLevelMax}{1.01}  % highest mean forecast / mean RV across the six full-set pooled forecasts, 614 days
\providecommand{\rhedCalScaleMin}{}\renewcommand{\rhedCalScaleMin}{1.09}  % lowest mean of RV / forecast (QLIKE-optimal scale) across the six full-set pooled forecasts, 614 days
\providecommand{\rhedCalScaleMax}{}\renewcommand{\rhedCalScaleMax}{1.20}  % highest mean of RV / forecast (QLIKE-optimal scale) across the six full-set pooled forecasts, 614 days
\providecommand{\rhedHarCalScaleMin}{}\renewcommand{\rhedHarCalScaleMin}{1.04}  % lowest mean of RV / forecast across the six HAR-benchmark pooled forecasts, 614 days
\providecommand{\rhedHarCalScaleMax}{}\renewcommand{\rhedHarCalScaleMax}{1.10}  % highest mean of RV / forecast across the six HAR-benchmark pooled forecasts, 614 days
\providecommand{\rhedMedFcIVMin}{}\renewcommand{\rhedMedFcIVMin}{0.75}  % lowest median forecast / IV^2 across the six full-set pooled forecasts, 614 days
\providecommand{\rhedMedFcIVMax}{}\renewcommand{\rhedMedFcIVMax}{0.82}  % highest median forecast / IV^2 across the six full-set pooled forecasts, 614 days
\providecommand{\rhedBcScaleLoMin}{}\renewcommand{\rhedBcScaleLoMin}{1.06}  % lowest lowest causal scale c_t over the 614 days across the six full-set pooled forecasts, 614 days
\providecommand{\rhedBcScaleLoMax}{}\renewcommand{\rhedBcScaleLoMax}{1.15}  % highest lowest causal scale c_t over the 614 days across the six full-set pooled forecasts, 614 days
\providecommand{\rhedBcScaleHiMin}{}\renewcommand{\rhedBcScaleHiMin}{1.10}  % lowest highest causal scale c_t over the 614 days across the six full-set pooled forecasts, 614 days
\providecommand{\rhedBcScaleHiMax}{}\renewcommand{\rhedBcScaleHiMax}{1.22}  % highest highest causal scale c_t over the 614 days across the six full-set pooled forecasts, 614 days
\providecommand{\rhedBcFlipMin}{}\renewcommand{\rhedBcFlipMin}{55}  % lowest positions changed by the bias correction across the six full-set pooled forecasts, 614 days
\providecommand{\rhedBcFlipMax}{}\renewcommand{\rhedBcFlipMax}{120}  % highest positions changed by the bias correction across the six full-set pooled forecasts, 614 days
\providecommand{\rhedBcLongMin}{}\renewcommand{\rhedBcLongMin}{201}  % lowest long days under the bias correction across the six full-set pooled forecasts, 614 days
\providecommand{\rhedBcLongMax}{}\renewcommand{\rhedBcLongMax}{241}  % highest long days under the bias correction across the six full-set pooled forecasts, 614 days
\providecommand{\rhedBcFlipRMin}{}\renewcommand{\rhedBcFlipRMin}{\ensuremath{-}24}  % lowest mean mid straddle return R on the changed days, % across the six full-set pooled forecasts, 614 days
\providecommand{\rhedBcFlipRMax}{}\renewcommand{\rhedBcFlipRMax}{\ensuremath{-}8}  % highest mean mid straddle return R on the changed days, % across the six full-set pooled forecasts, 614 days
\providecommand{\rhedBcShMin}{}\renewcommand{\rhedBcShMin}{0.92}  % lowest Sharpe ratio mid under the bias correction (diagnostic) across the six full-set pooled forecasts, 614 days
\providecommand{\rhedBcShMax}{}\renewcommand{\rhedBcShMax}{1.32}  % highest Sharpe ratio mid under the bias correction (diagnostic) across the six full-set pooled forecasts, 614 days
\providecommand{\rhedBcShXMin}{}\renewcommand{\rhedBcShXMin}{0.57}  % lowest Sharpe ratio bid/ask under the bias correction across the six full-set pooled forecasts, 614 days
\providecommand{\rhedBcShXMax}{}\renewcommand{\rhedBcShXMax}{0.97}  % highest Sharpe ratio bid/ask under the bias correction across the six full-set pooled forecasts, 614 days
\providecommand{\rhedLevFlipMin}{}\renewcommand{\rhedLevFlipMin}{4}  % lowest positions changed by the level correction (sum RV / sum fc) across the six full-set pooled forecasts, 614 days
\providecommand{\rhedLevFlipMax}{}\renewcommand{\rhedLevFlipMax}{19}  % highest positions changed by the level correction (sum RV / sum fc) across the six full-set pooled forecasts, 614 days
\providecommand{\rhedLevShMin}{}\renewcommand{\rhedLevShMin}{1.61}  % lowest Sharpe ratio mid under the level correction across the six full-set pooled forecasts, 614 days
\providecommand{\rhedLevShMax}{}\renewcommand{\rhedLevShMax}{2.01}  % highest Sharpe ratio mid under the level correction across the six full-set pooled forecasts, 614 days
\providecommand{\rhedHarBcShMin}{}\renewcommand{\rhedHarBcShMin}{0.76}  % lowest Sharpe ratio mid under the bias correction across the six HAR-benchmark pooled forecasts, 614 days
\providecommand{\rhedHarBcShMax}{}\renewcommand{\rhedHarBcShMax}{1.05}  % highest Sharpe ratio mid under the bias correction across the six HAR-benchmark pooled forecasts, 614 days
\providecommand{\rhedBcShUpN}{}\renewcommand{\rhedBcShUpN}{0}  % full-set forecasts whose mid SR rises under the bias correction
\providecommand{\rhedBcShDownN}{}\renewcommand{\rhedBcShDownN}{6}  % full-set forecasts whose mid SR falls under the bias correction
\providecommand{\rhedHarBcShUpN}{}\renewcommand{\rhedHarBcShUpN}{2}  % HAR-benchmark forecasts whose mid SR rises under the bias correction
\providecommand{\rhedMoveSqIV}{}\renewcommand{\rhedMoveSqIV}{1.01}  % mean of the squared S&P 500 half-hour log return over IV_t^2, 614 days
\providecommand{\rhedMoveSqRV}{}\renewcommand{\rhedMoveSqRV}{1.15}  % mean squared S&P 500 half-hour log return over mean RV_t, 614 days
\providecommand{\rhedOlsCalScale}{}\renewcommand{\rhedOlsCalScale}{1.01}  % mean of RV / forecast, closing-interval HAR forecast
\providecommand{\rhedOlsBcSh}{}\renewcommand{\rhedOlsBcSh}{1.12}  % Sharpe ratio mid under the bias correction, closing-interval HAR
\providecommand{\rhedBcFlipRidge}{}\renewcommand{\rhedBcFlipRidge}{105}  % positions changed by the bias correction, poolpt13_ridge_all_features
\providecommand{\rhedBcShRidge}{}\renewcommand{\rhedBcShRidge}{1.17}  % Sharpe ratio mid under the bias correction, poolpt13_ridge_all_features
\providecommand{\rhedBcFlipLasso}{}\renewcommand{\rhedBcFlipLasso}{120}  % positions changed by the bias correction, poolpt13_lasso_all_features
\providecommand{\rhedBcShLasso}{}\renewcommand{\rhedBcShLasso}{1.31}  % Sharpe ratio mid under the bias correction, poolpt13_lasso_all_features
\providecommand{\rhedBcFlipEnet}{}\renewcommand{\rhedBcFlipEnet}{108}  % positions changed by the bias correction, poolpt13_enet_all_features
\providecommand{\rhedBcShEnet}{}\renewcommand{\rhedBcShEnet}{0.99}  % Sharpe ratio mid under the bias correction, poolpt13_enet_all_features
\providecommand{\rhedBcFlipLgbm}{}\renewcommand{\rhedBcFlipLgbm}{55}  % positions changed by the bias correction, pooltree_lgbm_all_features
\providecommand{\rhedBcShLgbm}{}\renewcommand{\rhedBcShLgbm}{1.32}  % Sharpe ratio mid under the bias correction, pooltree_lgbm_all_features
\providecommand{\rhedBcFlipXgb}{}\renewcommand{\rhedBcFlipXgb}{61}  % positions changed by the bias correction, pooltree_xgb_all_features
\providecommand{\rhedBcShXgb}{}\renewcommand{\rhedBcShXgb}{1.15}  % Sharpe ratio mid under the bias correction, pooltree_xgb_all_features
\providecommand{\rhedBcFlipRf}{}\renewcommand{\rhedBcFlipRf}{90}  % positions changed by the bias correction, pooltree_rf_all_features
\providecommand{\rhedBcShRf}{}\renewcommand{\rhedBcShRf}{0.92}  % Sharpe ratio mid under the bias correction, pooltree_rf_all_features
\providecommand{\rhedQMonoN}{}\renewcommand{\rhedQMonoN}{6}  % full-set forecasts whose mean RV/IV^2 rises across all five quintiles
\providecommand{\rhedHarQMonoN}{}\renewcommand{\rhedHarQMonoN}{4}  % HAR-benchmark forecasts whose mean RV/IV^2 rises across all five quintiles
\providecommand{\rhedQRPatternN}{}\renewcommand{\rhedQRPatternN}{6}  % full-set forecasts with mean R < 0 in quintiles 1-4 and > 0 in 5
\providecommand{\rhedHarQRPatternN}{}\renewcommand{\rhedHarQRPatternN}{0}  % HAR-benchmark forecasts with that pattern
\providecommand{\rhedQOneRVIV}{}\renewcommand{\rhedQOneRVIV}{0.73}  % mean RV/IV^2 in the lowest quintile, mean of the six full-set
\providecommand{\rhedQFiveRVIV}{}\renewcommand{\rhedQFiveRVIV}{1.05}  % mean RV/IV^2 in the highest quintile, mean of the six full-set
\providecommand{\rhedQFiveR}{}\renewcommand{\rhedQFiveR}{26}  % mean straddle return R in the highest quintile, %, mean of the six full-set
\providecommand{\rhedQFiveFc}{}\renewcommand{\rhedQFiveFc}{1.18}  % geometric-mean forecast / IV^2 in the highest quintile, mean of the six full-set
\providecommand{\rhedQFourFc}{}\renewcommand{\rhedQFourFc}{0.93}  % ... in the fourth quintile
\providecommand{\rhedQFiveLong}{}\renewcommand{\rhedQFiveLong}{99}  % % long days in the highest quintile, mean of the six full-set
\providecommand{\rhedQFiveHit}{}\renewcommand{\rhedQFiveHit}{52}  % hit rate in the highest quintile, %, mean of the six full-set
\providecommand{\rhedQOneHit}{}\renewcommand{\rhedQOneHit}{88}  % hit rate in the lowest quintile, %, mean of the six full-set
\providecommand{\rhedQFiveShare}{}\renewcommand{\rhedQFiveShare}{43}  % % of the summed mid return earned in the highest quintile, six full-set forecasts
\providecommand{\rhedQLowRMin}{}\renewcommand{\rhedQLowRMin}{\ensuremath{-}13}  % most negative mean straddle return R in quintiles 1-4, %, mean of the six
\providecommand{\rhedQLowRMax}{}\renewcommand{\rhedQLowRMax}{\ensuremath{-}6}  % least negative mean straddle return R in quintiles 1-4, %, mean of the six
\providecommand{\rhedHarQOneRVIV}{}\renewcommand{\rhedHarQOneRVIV}{0.74}  % mean RV/IV^2 in the lowest quintile, mean of the six Har
\providecommand{\rhedHarQFiveRVIV}{}\renewcommand{\rhedHarQFiveRVIV}{1.02}  % mean RV/IV^2 in the highest quintile, mean of the six Har
\providecommand{\rhedHarQFiveR}{}\renewcommand{\rhedHarQFiveR}{7}  % mean straddle return R in the highest quintile, %, mean of the six Har
\providecommand{\rhedHarQFiveFc}{}\renewcommand{\rhedHarQFiveFc}{1.23}  % geometric-mean forecast / IV^2 in the highest quintile, mean of the six Har
\providecommand{\rhedHarQFourFc}{}\renewcommand{\rhedHarQFourFc}{0.98}  % ... in the fourth quintile
\providecommand{\rhedHarQFiveLong}{}\renewcommand{\rhedHarQFiveLong}{100}  % % long days in the highest quintile, mean of the six Har
\providecommand{\rhedHarQFiveHit}{}\renewcommand{\rhedHarQFiveHit}{47}  % hit rate in the highest quintile, %, mean of the six Har
\providecommand{\rhedHarQOneHit}{}\renewcommand{\rhedHarQOneHit}{89}  % hit rate in the lowest quintile, %, mean of the six Har
\providecommand{\rhedHarQFiveShare}{}\renewcommand{\rhedHarQFiveShare}{23}  % % of the summed mid return earned in the highest quintile, six Har forecasts
\providecommand{\rhedHarQLowRMin}{}\renewcommand{\rhedHarQLowRMin}{\ensuremath{-}13}  % most negative mean straddle return R in quintiles 1-4, %, mean of the six
\providecommand{\rhedHarQLowRMax}{}\renewcommand{\rhedHarQLowRMax}{5}  % least negative mean straddle return R in quintiles 1-4, %, mean of the six
\providecommand{\rhedQFiveShareRngMin}{}\renewcommand{\rhedQFiveShareRngMin}{40}  % lowest % of the summed mid return in quintile 5 across the six full-set pooled forecasts, 614 days
\providecommand{\rhedQFiveShareRngMax}{}\renewcommand{\rhedQFiveShareRngMax}{48}  % highest % of the summed mid return in quintile 5 across the six full-set pooled forecasts, 614 days
\providecommand{\rhedBestUnion}{}\renewcommand{\rhedBestUnion}{11}  % days in the union of the six full-set forecasts' ten best days
\providecommand{\rhedWorstUnion}{}\renewcommand{\rhedWorstUnion}{11}  % days in the union of the six full-set forecasts' ten worst days
\providecommand{\rhedBestAllSix}{}\renewcommand{\rhedBestAllSix}{9}  % days among the ten best of all six forecasts
\providecommand{\rhedWorstAllSix}{}\renewcommand{\rhedWorstAllSix}{9}  % days among the ten worst of all six forecasts
\providecommand{\rhedTopDate}{}\renewcommand{\rhedTopDate}{22~November~2021}  % the day of the largest straddle return among the best days
\providecommand{\rhedTopR}{}\renewcommand{\rhedTopR}{1{,}032}  % ... the straddle's mid return that day, % of premium
\providecommand{\rhedTopRVIV}{}\renewcommand{\rhedTopRVIV}{2.79}  % ... RV / IV^2 that day
\providecommand{\rhedTopMove}{}\renewcommand{\rhedTopMove}{\ensuremath{-}0.83}  % ... S&P 500 15:30-16:00 log return that day, %
\providecommand{\rhedTopShareMin}{}\renewcommand{\rhedTopShareMin}{12}  % lowest % of a full-set forecast's summed mid return earned on that one day
\providecommand{\rhedTopShareMax}{}\renewcommand{\rhedTopShareMax}{15}  % highest % of a full-set forecast's summed mid return earned on that one day
\providecommand{\rhedBestFomcN}{}\renewcommand{\rhedBestFomcN}{1}  % best-day union: FOMC statement days
\providecommand{\rhedWorstFomcN}{}\renewcommand{\rhedWorstFomcN}{1}  % worst-day union: FOMC statement days
\providecommand{\rhedWorstMeN}{}\renewcommand{\rhedWorstMeN}{3}  % worst-day union: month ends
\providecommand{\rhedBestMeN}{}\renewcommand{\rhedBestMeN}{0}  % best-day union: month ends
\providecommand{\rhedBestOpexN}{}\renewcommand{\rhedBestOpexN}{0}  % best-day union: monthly expirations
\providecommand{\rhedWorstOpexN}{}\renewcommand{\rhedWorstOpexN}{0}  % worst-day union: monthly expirations
\providecommand{\rhedBestNoTagN}{}\renewcommand{\rhedBestNoTagN}{8}  % best-day union: days without any event tag (calendar as covered)
\providecommand{\rhedBestAfterRelN}{}\renewcommand{\rhedBestAfterRelN}{5}  % best-day union: days after the release calendar ends (dagger)
\providecommand{\rhedWorstAfterRelN}{}\renewcommand{\rhedWorstAfterRelN}{4}  % worst-day union: days after the release calendar ends
\providecommand{\rhedDShortAllMin}{}\renewcommand{\rhedDShortAllMin}{1.21}  % lowest Delta SR mid vs the short straddle, all days across the six full-set pooled forecasts, 614 days
\providecommand{\rhedDShortAllMax}{}\renewcommand{\rhedDShortAllMax}{1.56}  % highest Delta SR mid vs the short straddle, all days across the six full-set pooled forecasts, 614 days
\providecommand{\rhedDShortExBestMin}{}\renewcommand{\rhedDShortExBestMin}{\ensuremath{-}0.54}  % lowest Delta SR mid vs the short straddle without the forecast's ten best days across the six full-set pooled forecasts, 614 days
\providecommand{\rhedDShortExBestMax}{}\renewcommand{\rhedDShortExBestMax}{\ensuremath{-}0.12}  % highest Delta SR mid vs the short straddle without the forecast's ten best days across the six full-set pooled forecasts, 614 days
\providecommand{\rhedDShortExBothMin}{}\renewcommand{\rhedDShortExBothMin}{\ensuremath{-}0.64}  % lowest ... without its ten best and ten worst days across the six full-set pooled forecasts, 614 days
\providecommand{\rhedDShortExBothMax}{}\renewcommand{\rhedDShortExBothMax}{\ensuremath{-}0.14}  % highest ... without its ten best and ten worst days across the six full-set pooled forecasts, 614 days
\providecommand{\rhedBestShareMin}{}\renewcommand{\rhedBestShareMin}{107}  % lowest % of the summed gain over the short straddle earned on the ten best days across the six full-set pooled forecasts, 614 days
\providecommand{\rhedBestShareMax}{}\renewcommand{\rhedBestShareMax}{138}  % highest % of the summed gain over the short straddle earned on the ten best days across the six full-set pooled forecasts, 614 days
\providecommand{\rhedDHarAllMin}{}\renewcommand{\rhedDHarAllMin}{0.56}  % lowest Delta SR mid vs the HAR twin, all days across the six full-set pooled forecasts, 614 days
\providecommand{\rhedDHarAllMax}{}\renewcommand{\rhedDHarAllMax}{1.56}  % highest Delta SR mid vs the HAR twin, all days across the six full-set pooled forecasts, 614 days
\providecommand{\rhedDHarExBestMin}{}\renewcommand{\rhedDHarExBestMin}{0.30}  % lowest Delta SR mid vs the HAR twin without the forecast's ten best days across the six full-set pooled forecasts, 614 days
\providecommand{\rhedDHarExBestMax}{}\renewcommand{\rhedDHarExBestMax}{0.83}  % highest Delta SR mid vs the HAR twin without the forecast's ten best days across the six full-set pooled forecasts, 614 days
\providecommand{\rhedDHarExBothMin}{}\renewcommand{\rhedDHarExBothMin}{0.36}  % lowest ... without its ten best and ten worst days across the six full-set pooled forecasts, 614 days
\providecommand{\rhedDHarExBothMax}{}\renewcommand{\rhedDHarExBothMax}{0.97}  % highest ... without its ten best and ten worst days across the six full-set pooled forecasts, 614 days
\providecommand{\rhedHarBestShareMin}{}\renewcommand{\rhedHarBestShareMin}{47}  % lowest % of the summed gain over the HAR twin earned on the ten best days across the six full-set pooled forecasts, 614 days
\providecommand{\rhedHarBestShareMax}{}\renewcommand{\rhedHarBestShareMax}{61}  % highest % of the summed gain over the HAR twin earned on the ten best days across the six full-set pooled forecasts, 614 days
\providecommand{\rhedDShortExBestNeg}{}\renewcommand{\rhedDShortExBestNeg}{6}  % full-set forecasts whose Delta SR vs short is negative without the ten best days
\providecommand{\rhedDHarExBestNeg}{}\renewcommand{\rhedDHarExBestNeg}{0}  % full-set forecasts whose Delta SR vs the HAR twin is negative without the ten best days
\providecommand{\rhedFomcN}{}\renewcommand{\rhedFomcN}{21}  % FOMC days among the 614 post-warmup days
\providecommand{\rhedFomcR}{}\renewcommand{\rhedFomcR}{25}  % mean mid straddle return R on FOMC days, %
\providecommand{\rhedFomcRT}{}\renewcommand{\rhedFomcRT}{1.02}  % ... its t (i.i.d.; the days are weeks apart)
\providecommand{\rhedFomcLong}{}\renewcommand{\rhedFomcLong}{25}  % % of FOMC days long, mean of the six full-set forecasts
\providecommand{\rhedFomcStrat}{}\renewcommand{\rhedFomcStrat}{\ensuremath{-}6}  % mean mid strategy return on FOMC days, %, mean of the six full-set
\providecommand{\rhedFomcOtherR}{}\renewcommand{\rhedFomcOtherR}{\ensuremath{-}3}  % mean straddle R on the other (covered) days, %
\providecommand{\rhedFomcOtherStrat}{}\renewcommand{\rhedFomcOtherStrat}{13}  % mean strategy return on the other (covered) days, %
\providecommand{\rhedFomcRVIV}{}\renewcommand{\rhedFomcRVIV}{0.76}  % mean RV_t / IV_t^2 on FOMC days
\providecommand{\rhedFomcMoveSqIV}{}\renewcommand{\rhedFomcMoveSqIV}{1.43}  % mean squared S&P 500 half-hour log return / IV_t^2 on FOMC days
\providecommand{\rhedFomcShare}{}\renewcommand{\rhedFomcShare}{\ensuremath{-}2}  % % of the summed mid return earned on FOMC days, mean of the six full-set forecasts
\providecommand{\rhedCpiN}{}\renewcommand{\rhedCpiN}{25}  % CPI days among the 614 post-warmup days
\providecommand{\rhedCpiR}{}\renewcommand{\rhedCpiR}{\ensuremath{-}12}  % mean mid straddle return R on CPI days, %
\providecommand{\rhedCpiRT}{}\renewcommand{\rhedCpiRT}{\ensuremath{-}0.55}  % ... its t (i.i.d.; the days are weeks apart)
\providecommand{\rhedCpiLong}{}\renewcommand{\rhedCpiLong}{16}  % % of CPI days long, mean of the six full-set forecasts
\providecommand{\rhedCpiStrat}{}\renewcommand{\rhedCpiStrat}{22}  % mean mid strategy return on CPI days, %, mean of the six full-set
\providecommand{\rhedCpiOtherR}{}\renewcommand{\rhedCpiOtherR}{\ensuremath{-}6}  % mean straddle R on the other (covered) days, %
\providecommand{\rhedCpiOtherStrat}{}\renewcommand{\rhedCpiOtherStrat}{11}  % mean strategy return on the other (covered) days, %
\providecommand{\rhedCpiRVIV}{}\renewcommand{\rhedCpiRVIV}{0.80}  % mean RV_t / IV_t^2 on CPI days
\providecommand{\rhedCpiMoveSqIV}{}\renewcommand{\rhedCpiMoveSqIV}{0.86}  % mean squared S&P 500 half-hour log return / IV_t^2 on CPI days
\providecommand{\rhedCpiShare}{}\renewcommand{\rhedCpiShare}{7}  % % of the summed mid return earned on CPI days, mean of the six full-set forecasts
\providecommand{\rhedNfpN}{}\renewcommand{\rhedNfpN}{27}  % NFP days among the 614 post-warmup days
\providecommand{\rhedNfpR}{}\renewcommand{\rhedNfpR}{1}  % mean mid straddle return R on NFP days, %
\providecommand{\rhedNfpRT}{}\renewcommand{\rhedNfpRT}{0.06}  % ... its t (i.i.d.; the days are weeks apart)
\providecommand{\rhedNfpLong}{}\renewcommand{\rhedNfpLong}{12}  % % of NFP days long, mean of the six full-set forecasts
\providecommand{\rhedNfpStrat}{}\renewcommand{\rhedNfpStrat}{20}  % mean mid strategy return on NFP days, %, mean of the six full-set
\providecommand{\rhedNfpOtherR}{}\renewcommand{\rhedNfpOtherR}{\ensuremath{-}6}  % mean straddle R on the other (covered) days, %
\providecommand{\rhedNfpOtherStrat}{}\renewcommand{\rhedNfpOtherStrat}{11}  % mean strategy return on the other (covered) days, %
\providecommand{\rhedNfpRVIV}{}\renewcommand{\rhedNfpRVIV}{0.73}  % mean RV_t / IV_t^2 on NFP days
\providecommand{\rhedNfpMoveSqIV}{}\renewcommand{\rhedNfpMoveSqIV}{0.93}  % mean squared S&P 500 half-hour log return / IV_t^2 on NFP days
\providecommand{\rhedNfpShare}{}\renewcommand{\rhedNfpShare}{7}  % % of the summed mid return earned on NFP days, mean of the six full-set forecasts
\providecommand{\rhedOpexN}{}\renewcommand{\rhedOpexN}{33}  % OpEx days among the 614 post-warmup days
\providecommand{\rhedOpexR}{}\renewcommand{\rhedOpexR}{\ensuremath{-}45}  % mean mid straddle return R on OpEx days, %
\providecommand{\rhedOpexRT}{}\renewcommand{\rhedOpexRT}{\ensuremath{-}4.01}  % ... its t (i.i.d.; the days are weeks apart)
\providecommand{\rhedOpexLong}{}\renewcommand{\rhedOpexLong}{9}  % % of OpEx days long, mean of the six full-set forecasts
\providecommand{\rhedOpexStrat}{}\renewcommand{\rhedOpexStrat}{36}  % mean mid strategy return on OpEx days, %, mean of the six full-set
\providecommand{\rhedOpexOtherR}{}\renewcommand{\rhedOpexOtherR}{0}  % mean straddle R on the other (covered) days, %
\providecommand{\rhedOpexOtherStrat}{}\renewcommand{\rhedOpexOtherStrat}{11}  % mean strategy return on the other (covered) days, %
\providecommand{\rhedOpexRVIV}{}\renewcommand{\rhedOpexRVIV}{0.83}  % mean RV_t / IV_t^2 on OpEx days
\providecommand{\rhedOpexMoveSqIV}{}\renewcommand{\rhedOpexMoveSqIV}{0.40}  % mean squared S&P 500 half-hour log return / IV_t^2 on OpEx days
\providecommand{\rhedOpexShare}{}\renewcommand{\rhedOpexShare}{15}  % % of the summed mid return earned on OpEx days, mean of the six full-set forecasts
\providecommand{\rhedTwN}{}\renewcommand{\rhedTwN}{11}  % TW days among the 614 post-warmup days
\providecommand{\rhedTwR}{}\renewcommand{\rhedTwR}{\ensuremath{-}52}  % mean mid straddle return R on TW days, %
\providecommand{\rhedTwRT}{}\renewcommand{\rhedTwRT}{\ensuremath{-}3.98}  % ... its t (i.i.d.; the days are weeks apart)
\providecommand{\rhedTwLong}{}\renewcommand{\rhedTwLong}{8}  % % of TW days long, mean of the six full-set forecasts
\providecommand{\rhedTwStrat}{}\renewcommand{\rhedTwStrat}{37}  % mean mid strategy return on TW days, %, mean of the six full-set
\providecommand{\rhedTwOtherR}{}\renewcommand{\rhedTwOtherR}{\ensuremath{-}2}  % mean straddle R on the other (covered) days, %
\providecommand{\rhedTwOtherStrat}{}\renewcommand{\rhedTwOtherStrat}{12}  % mean strategy return on the other (covered) days, %
\providecommand{\rhedTwRVIV}{}\renewcommand{\rhedTwRVIV}{0.89}  % mean RV_t / IV_t^2 on TW days
\providecommand{\rhedTwMoveSqIV}{}\renewcommand{\rhedTwMoveSqIV}{0.31}  % mean squared S&P 500 half-hour log return / IV_t^2 on TW days
\providecommand{\rhedTwShare}{}\renewcommand{\rhedTwShare}{5}  % % of the summed mid return earned on TW days, mean of the six full-set forecasts
\providecommand{\rhedMeN}{}\renewcommand{\rhedMeN}{33}  % ME days among the 614 post-warmup days
\providecommand{\rhedMeR}{}\renewcommand{\rhedMeR}{70}  % mean mid straddle return R on ME days, %
\providecommand{\rhedMeRT}{}\renewcommand{\rhedMeRT}{2.79}  % ... its t (i.i.d.; the days are weeks apart)
\providecommand{\rhedMeLong}{}\renewcommand{\rhedMeLong}{6}  % % of ME days long, mean of the six full-set forecasts
\providecommand{\rhedMeStrat}{}\renewcommand{\rhedMeStrat}{\ensuremath{-}58}  % mean mid strategy return on ME days, %, mean of the six full-set
\providecommand{\rhedMeOtherR}{}\renewcommand{\rhedMeOtherR}{\ensuremath{-}7}  % mean straddle R on the other (covered) days, %
\providecommand{\rhedMeOtherStrat}{}\renewcommand{\rhedMeOtherStrat}{17}  % mean strategy return on the other (covered) days, %
\providecommand{\rhedMeRVIV}{}\renewcommand{\rhedMeRVIV}{1.04}  % mean RV_t / IV_t^2 on ME days
\providecommand{\rhedMeMoveSqIV}{}\renewcommand{\rhedMeMoveSqIV}{1.93}  % mean squared S&P 500 half-hour log return / IV_t^2 on ME days
\providecommand{\rhedMeShare}{}\renewcommand{\rhedMeShare}{\ensuremath{-}24}  % % of the summed mid return earned on ME days, mean of the six full-set forecasts
\providecommand{\rhedRelEndText}{}\renewcommand{\rhedRelEndText}{26~October~2023}  % earliest last flag of the six macro-release series in data/releases.parquet
\providecommand{\rhedRelFomcEndText}{}\renewcommand{\rhedRelFomcEndText}{1~November~2023}  % last FOMC flag in data/releases.parquet
\providecommand{\rhedFomcSuppN}{}\renewcommand{\rhedFomcSuppN}{3}  % FOMC statement days of the sample taken from live/close_signal/state/fomc_statement_dates.csv
\providecommand{\rhedRelAfterN}{}\renewcommand{\rhedRelAfterN}{126}  % post-warmup days after the release calendar ends
\providecommand{\rhedWorthN}{}\renewcommand{\rhedWorthN}{137}  % post-warmup days with R_t = -1 (= rheWorthlessDays)
\providecommand{\rhedWorthPct}{}\renewcommand{\rhedWorthPct}{22.3}  % ... % of the 614 days
\providecommand{\rhedWorthExpIV}{}\renewcommand{\rhedWorthExpIV}{128}  % expected R = -1 days under a normal half-hour log move with volatility IV_t
\providecommand{\rhedWorthExpIVPct}{}\renewcommand{\rhedWorthExpIVPct}{20.9}  % ... as % of the 614 days
\providecommand{\rhedWorthZ}{}\renewcommand{\rhedWorthZ}{0.90}  % (observed - expected) / sd of the count under that normal (independent days)
\providecommand{\rhedWorthExpEmp}{}\renewcommand{\rhedWorthExpEmp}{132}  % expected R = -1 days under the empirical distribution of the earlier days' moves standardized by IV
\providecommand{\rhedWorthExpRaw}{}\renewcommand{\rhedWorthExpRaw}{118}  % expected R = -1 days under the earlier days' raw half-hour moves
\providecommand{\rhedGapFiveN}{}\renewcommand{\rhedGapFiveN}{611}  % post-warmup days with K_C - K_P = 5 points
\providecommand{\rhedGapOtherN}{}\renewcommand{\rhedGapOtherN}{3}  % post-warmup days with another strike gap
\providecommand{\rhedMoveSdPts}{}\renewcommand{\rhedMoveSdPts}{9.6}  % median one-sd implied half-hour move, index points
\providecommand{\rhedStrikeDist}{}\renewcommand{\rhedStrikeDist}{1.23}  % mean distance of the close to the nearest 5-point strike, points
\providecommand{\rhedStrikeDistUnif}{}\renewcommand{\rhedStrikeDistUnif}{1.25}  % ... its mean under a uniform close
\providecommand{\rhedKsD}{}\renewcommand{\rhedKsD}{0.026}  % Kolmogorov-Smirnov D of the close's position within its 5-point strike interval against a uniform
\providecommand{\rhedKsCV}{}\renewcommand{\rhedKsCV}{0.055}  % ... asymptotic 5% critical value
\providecommand{\rhedZSd}{}\renewcommand{\rhedZSd}{1.00}  % sd of the half-hour log move standardized by IV_t, 614 days
\providecommand{\rhedZKurt}{}\renewcommand{\rhedZKurt}{2.0}  % excess kurtosis of that standardized move
\providecommand{\rhedWorthShortMin}{}\renewcommand{\rhedWorthShortMin}{110}  % lowest R = -1 days on which the forecast is short across the six full-set pooled forecasts, 614 days
\providecommand{\rhedWorthShortMax}{}\renewcommand{\rhedWorthShortMax}{115}  % highest R = -1 days on which the forecast is short across the six full-set pooled forecasts, 614 days
\providecommand{\rhedWorthOpexN}{}\renewcommand{\rhedWorthOpexN}{10}  % R = -1 days among the monthly-expiration days
\providecommand{\rhedWorthOpexExp}{}\renewcommand{\rhedWorthOpexExp}{6.1}  % ... expected under the normal with IV_t
\providecommand{\rhedWorthOtherPct}{}\renewcommand{\rhedWorthOtherPct}{21.9}  % % of the other days with R = -1
\providecommand{\rhedWorthOtherExpPct}{}\renewcommand{\rhedWorthOtherExpPct}{21.0}  % ... expected under the normal with IV_t, %
\providecommand{\rhedWorthOpexZ}{}\renewcommand{\rhedWorthOpexZ}{1.77}  % (observed - expected) / sd of the count on the monthly-expiration days
\providecommand{\rhedRegAlphaAMin}{}\renewcommand{\rhedRegAlphaAMin}{11.0}  % lowest alpha, (a) on the short straddle, mid, % per day across the six full-set pooled forecasts, 614 days
\providecommand{\rhedRegAlphaAMax}{}\renewcommand{\rhedRegAlphaAMax}{13.4}  % highest alpha, (a) on the short straddle, mid, % per day across the six full-set pooled forecasts, 614 days
\providecommand{\rhedRegAlphaAXMin}{}\renewcommand{\rhedRegAlphaAXMin}{8.3}  % lowest alpha, (a) on the short straddle, bid/ask, % per day across the six full-set pooled forecasts, 614 days
\providecommand{\rhedRegAlphaAXMax}{}\renewcommand{\rhedRegAlphaAXMax}{10.7}  % highest alpha, (a) on the short straddle, bid/ask, % per day across the six full-set pooled forecasts, 614 days
\providecommand{\rhedRegBetaShortMin}{}\renewcommand{\rhedRegBetaShortMin}{0.23}  % lowest beta on the short-straddle return, (a), mid across the six full-set pooled forecasts, 614 days
\providecommand{\rhedRegBetaShortMax}{}\renewcommand{\rhedRegBetaShortMax}{0.28}  % highest beta on the short-straddle return, (a), mid across the six full-set pooled forecasts, 614 days
\providecommand{\rhedRegAlphaAPMin}{}\renewcommand{\rhedRegAlphaAPMin}{0.005}  % lowest fixed-b p of alpha, (a), mid across the six full-set pooled forecasts, 614 days
\providecommand{\rhedRegAlphaAPMax}{}\renewcommand{\rhedRegAlphaAPMax}{0.025}  % highest fixed-b p of alpha, (a), mid across the six full-set pooled forecasts, 614 days
\providecommand{\rhedRegAlphaAXPMin}{}\renewcommand{\rhedRegAlphaAXPMin}{0.021}  % lowest fixed-b p of alpha, (a), bid/ask across the six full-set pooled forecasts, 614 days
\providecommand{\rhedRegAlphaAXPMax}{}\renewcommand{\rhedRegAlphaAXPMax}{0.077}  % highest fixed-b p of alpha, (a), bid/ask across the six full-set pooled forecasts, 614 days
\providecommand{\rhedRegBetaShortTMin}{}\renewcommand{\rhedRegBetaShortTMin}{1.22}  % lowest t of the beta on the short straddle, (a), mid across the six full-set pooled forecasts, 614 days
\providecommand{\rhedRegBetaShortTMax}{}\renewcommand{\rhedRegBetaShortTMax}{1.42}  % highest t of the beta on the short straddle, (a), mid across the six full-set pooled forecasts, 614 days
\providecommand{\rhedRegRsqAMin}{}\renewcommand{\rhedRegRsqAMin}{0.06}  % lowest R^2 of (a), mid across the six full-set pooled forecasts, 614 days
\providecommand{\rhedRegRsqAMax}{}\renewcommand{\rhedRegRsqAMax}{0.08}  % highest R^2 of (a), mid across the six full-set pooled forecasts, 614 days
\providecommand{\rhedRegBetaAbsMin}{}\renewcommand{\rhedRegBetaAbsMin}{\ensuremath{-}1.87}  % lowest beta on |r30| (bp), (b), mid, % of premium per bp across the six full-set pooled forecasts, 614 days
\providecommand{\rhedRegBetaAbsMax}{}\renewcommand{\rhedRegBetaAbsMax}{\ensuremath{-}1.40}  % highest beta on |r30| (bp), (b), mid, % of premium per bp across the six full-set pooled forecasts, 614 days
\providecommand{\rhedRegBetaAbsTMin}{}\renewcommand{\rhedRegBetaAbsTMin}{\ensuremath{-}3.49}  % lowest t of the beta on |r30|, (b), mid across the six full-set pooled forecasts, 614 days
\providecommand{\rhedRegBetaAbsTMax}{}\renewcommand{\rhedRegBetaAbsTMax}{\ensuremath{-}2.31}  % highest t of the beta on |r30|, (b), mid across the six full-set pooled forecasts, 614 days
\providecommand{\rhedRegBetaMoveMin}{}\renewcommand{\rhedRegBetaMoveMin}{\ensuremath{-}0.01}  % lowest beta on r30 (bp), (b), mid across the six full-set pooled forecasts, 614 days
\providecommand{\rhedRegBetaMoveMax}{}\renewcommand{\rhedRegBetaMoveMax}{0.14}  % highest beta on r30 (bp), (b), mid across the six full-set pooled forecasts, 614 days
\providecommand{\rhedRegMeanMin}{}\renewcommand{\rhedRegMeanMin}{11.5}  % lowest mean return, mid, % per day across the six full-set pooled forecasts, 614 days
\providecommand{\rhedRegMeanMax}{}\renewcommand{\rhedRegMeanMax}{14.1}  % highest mean return, mid, % per day across the six full-set pooled forecasts, 614 days
\providecommand{\rhedRegMeanXMin}{}\renewcommand{\rhedRegMeanXMin}{8.9}  % lowest mean return, bid/ask, % per day across the six full-set pooled forecasts, 614 days
\providecommand{\rhedRegMeanXMax}{}\renewcommand{\rhedRegMeanXMax}{11.4}  % highest mean return, bid/ask, % per day across the six full-set pooled forecasts, 614 days
\providecommand{\rhedRegRsqCMin}{}\renewcommand{\rhedRegRsqCMin}{0.00}  % lowest R^2 of (c), mid across the six full-set pooled forecasts, 614 days
\providecommand{\rhedRegRsqCMax}{}\renewcommand{\rhedRegRsqCMax}{0.00}  % highest R^2 of (c), mid across the six full-set pooled forecasts, 614 days
\providecommand{\rhedRegCPMaxMin}{}\renewcommand{\rhedRegCPMaxMin}{0.74}  % lowest larger fixed-b p of the two betas of (c), mid across the six full-set pooled forecasts, 614 days
\providecommand{\rhedRegCPMaxMax}{}\renewcommand{\rhedRegCPMaxMax}{0.97}  % highest larger fixed-b p of the two betas of (c), mid across the six full-set pooled forecasts, 614 days
\providecommand{\rhedRegCPMinMin}{}\renewcommand{\rhedRegCPMinMin}{0.56}  % lowest smaller fixed-b p of the two betas of (c), mid across the six full-set pooled forecasts, 614 days
\providecommand{\rhedRegCPMinMax}{}\renewcommand{\rhedRegCPMinMax}{0.82}  % highest smaller fixed-b p of the two betas of (c), mid across the six full-set pooled forecasts, 614 days
\providecommand{\rhedRegAlphaDMin}{}\renewcommand{\rhedRegAlphaDMin}{11.4}  % lowest alpha, (d) joint, mid, % per day across the six full-set pooled forecasts, 614 days
\providecommand{\rhedRegAlphaDMax}{}\renewcommand{\rhedRegAlphaDMax}{13.6}  % highest alpha, (d) joint, mid, % per day across the six full-set pooled forecasts, 614 days
\providecommand{\rhedRegAlphaDXMin}{}\renewcommand{\rhedRegAlphaDXMin}{8.7}  % lowest alpha, (d) joint, bid/ask, % per day across the six full-set pooled forecasts, 614 days
\providecommand{\rhedRegAlphaDXMax}{}\renewcommand{\rhedRegAlphaDXMax}{10.9}  % highest alpha, (d) joint, bid/ask, % per day across the six full-set pooled forecasts, 614 days
\providecommand{\rhedRegAlphaDPMin}{}\renewcommand{\rhedRegAlphaDPMin}{0.004}  % lowest fixed-b p of alpha, (d), mid across the six full-set pooled forecasts, 614 days
\providecommand{\rhedRegAlphaDPMax}{}\renewcommand{\rhedRegAlphaDPMax}{0.026}  % highest fixed-b p of alpha, (d), mid across the six full-set pooled forecasts, 614 days
\providecommand{\rhedRegAlphaDXPMin}{}\renewcommand{\rhedRegAlphaDXPMin}{0.017}  % lowest fixed-b p of alpha, (d), bid/ask across the six full-set pooled forecasts, 614 days
\providecommand{\rhedRegAlphaDXPMax}{}\renewcommand{\rhedRegAlphaDXPMax}{0.076}  % highest fixed-b p of alpha, (d), bid/ask across the six full-set pooled forecasts, 614 days
\providecommand{\rhedRegAvgAlphaM}{}\renewcommand{\rhedRegAvgAlphaM}{12.7}  % alpha of the mean of the six full-set returns, (m), mid, % per day
\providecommand{\rhedRegAvgAlphaMT}{}\renewcommand{\rhedRegAvgAlphaMT}{3.48}  % ... its t
\providecommand{\rhedRegAvgAlphaMP}{}\renewcommand{\rhedRegAvgAlphaMP}{0.002}  % ... its fixed-b p
\providecommand{\rhedRegAvgAlphaA}{}\renewcommand{\rhedRegAvgAlphaA}{12.1}  % alpha of the mean of the six full-set returns, (a), mid, % per day
\providecommand{\rhedRegAvgAlphaAT}{}\renewcommand{\rhedRegAvgAlphaAT}{2.83}  % ... its t
\providecommand{\rhedRegAvgAlphaAP}{}\renewcommand{\rhedRegAvgAlphaAP}{0.0100}  % ... its fixed-b p
\providecommand{\rhedRegAvgAlphaD}{}\renewcommand{\rhedRegAvgAlphaD}{12.5}  % alpha of the mean of the six full-set returns, (d), mid, % per day
\providecommand{\rhedRegAvgAlphaDT}{}\renewcommand{\rhedRegAvgAlphaDT}{2.86}  % ... its t
\providecommand{\rhedRegAvgAlphaDP}{}\renewcommand{\rhedRegAvgAlphaDP}{0.009}  % ... its fixed-b p
\providecommand{\rhedRegAvgAlphaMX}{}\renewcommand{\rhedRegAvgAlphaMX}{10.1}  % alpha of the mean of the six full-set returns, (m), ba, % per day
\providecommand{\rhedRegAvgAlphaMXT}{}\renewcommand{\rhedRegAvgAlphaMXT}{2.78}  % ... its t
\providecommand{\rhedRegAvgAlphaMXP}{}\renewcommand{\rhedRegAvgAlphaMXP}{0.011}  % ... its fixed-b p
\providecommand{\rhedRegAvgAlphaAX}{}\renewcommand{\rhedRegAvgAlphaAX}{9.4}  % alpha of the mean of the six full-set returns, (a), ba, % per day
\providecommand{\rhedRegAvgAlphaAXT}{}\renewcommand{\rhedRegAvgAlphaAXT}{2.21}  % ... its t
\providecommand{\rhedRegAvgAlphaAXP}{}\renewcommand{\rhedRegAvgAlphaAXP}{0.040}  % ... its fixed-b p
\providecommand{\rhedRegAvgAlphaDX}{}\renewcommand{\rhedRegAvgAlphaDX}{9.8}  % alpha of the mean of the six full-set returns, (d), ba, % per day
\providecommand{\rhedRegAvgAlphaDXT}{}\renewcommand{\rhedRegAvgAlphaDXT}{2.26}  % ... its t
\providecommand{\rhedRegAvgAlphaDXP}{}\renewcommand{\rhedRegAvgAlphaDXP}{0.036}  % ... its fixed-b p
\providecommand{\rhedRegAvgBetaShort}{}\renewcommand{\rhedRegAvgBetaShort}{0.26}  % beta of the mean full-set return on the short straddle, (a), mid
\providecommand{\rhedRegAvgBetaShortT}{}\renewcommand{\rhedRegAvgBetaShortT}{1.34}  % ... its t
\providecommand{\rhedRegAvgBetaShortP}{}\renewcommand{\rhedRegAvgBetaShortP}{0.20}  % ... its fixed-b p
\providecommand{\rhedRegAvgRsqA}{}\renewcommand{\rhedRegAvgRsqA}{0.07}  % R^2 of (a) for the mean full-set return
\providecommand{\rhedRegAvgBetaAbs}{}\renewcommand{\rhedRegAvgBetaAbs}{\ensuremath{-}1.68}  % beta of the mean full-set return on |r30| (bp), (b), mid
\providecommand{\rhedRegAvgBetaAbsT}{}\renewcommand{\rhedRegAvgBetaAbsT}{\ensuremath{-}3.26}  % ... its t
\providecommand{\rhedRegShortBetaAbs}{}\renewcommand{\rhedRegShortBetaAbs}{\ensuremath{-}3.96}  % beta of the short straddle on |r30| (bp), (b), mid
\providecommand{\rhedRegShortRsqB}{}\renewcommand{\rhedRegShortRsqB}{0.48}  % R^2 of (b) for the short straddle
\providecommand{\rhedRegAbsRatio}{}\renewcommand{\rhedRegAbsRatio}{43}  % beta on |r30| of the mean full-set return as % of the short straddle's
\providecommand{\rhedRegAvgBetaAbsJ}{}\renewcommand{\rhedRegAvgBetaAbsJ}{\ensuremath{-}1.30}  % beta of the mean full-set return on |r30|, (d) joint, mid
\providecommand{\rhedRegAvgBetaAbsJT}{}\renewcommand{\rhedRegAvgBetaAbsJT}{\ensuremath{-}1.90}  % ... its t
\providecommand{\rhedHarRegAlphaAMin}{}\renewcommand{\rhedHarRegAlphaAMin}{1.1}  % lowest alpha, (a), mid, % per day across the six HAR-benchmark pooled forecasts, 614 days
\providecommand{\rhedHarRegAlphaAMax}{}\renewcommand{\rhedHarRegAlphaAMax}{7.0}  % highest alpha, (a), mid, % per day across the six HAR-benchmark pooled forecasts, 614 days
\providecommand{\rhedHarRegAlphaAPMin}{}\renewcommand{\rhedHarRegAlphaAPMin}{0.082}  % lowest fixed-b p of alpha, (a), mid across the six HAR-benchmark pooled forecasts, 614 days
\providecommand{\rhedHarRegAlphaAPMax}{}\renewcommand{\rhedHarRegAlphaAPMax}{0.745}  % highest fixed-b p of alpha, (a), mid across the six HAR-benchmark pooled forecasts, 614 days
\providecommand{\rhedOlsRegAlphaA}{}\renewcommand{\rhedOlsRegAlphaA}{8.1}  % alpha, (a), mid, closing-interval HAR forecast, % per day
\providecommand{\rhedOlsRegAlphaAP}{}\renewcommand{\rhedOlsRegAlphaAP}{0.051}  % ... its fixed-b p
\providecommand{\rhedRegAlphaAExBestMin}{}\renewcommand{\rhedRegAlphaAExBestMin}{\ensuremath{-}0.5}  % lowest alpha, (a), mid, without the forecast's ten best days, % per day across the six full-set pooled forecasts, 614 days
\providecommand{\rhedRegAlphaAExBestMax}{}\renewcommand{\rhedRegAlphaAExBestMax}{2.1}  % highest alpha, (a), mid, without the forecast's ten best days, % per day across the six full-set pooled forecasts, 614 days
\providecommand{\rhedRegAlphaAExBestPMin}{}\renewcommand{\rhedRegAlphaAExBestPMin}{0.43}  % lowest ... its fixed-b p across the six full-set pooled forecasts, 614 days
\providecommand{\rhedRegAlphaAExBestPMax}{}\renewcommand{\rhedRegAlphaAExBestPMax}{0.97}  % highest ... its fixed-b p across the six full-set pooled forecasts, 614 days
\fi
\begin{table}[tbp]
\centering
\footnotesize
\setlength{\tabcolsep}{2.5pt}
\begin{threeparttable}
\caption{Risk exposures of the sign strategy}
\label{tab:rh_econ_riskreg}
\begin{tabular}{@{}lrrrrrrrr@{}}
\toprule
 & \multicolumn{6}{c}{Full predictor set} & & \\
\cmidrule(lr){2-7}
 & Ridge & Lasso & \makecell[r]{Elastic\\net} & LightGBM & XGBoost & \makecell[r]{Random\\forest} & Mean & \makecell[r]{Short\\straddle} \\
\midrule
\multicolumn{9}{@{}l}{\textit{Panel A: Returns at mid-quotes}} \\
\addlinespace[2pt]
Mean return (\%) & $13.3^{***}$ & $12.9^{***}$ & $13.9^{***}$ & $11.9^{***}$ & $11.1^{***}$ & $14.1^{***}$ & $12.9^{***}$ & $2.4^{\phantom{***}}$ \\
 & $(3.38)^{\phantom{***}}$ & $(3.31)^{\phantom{***}}$ & $(3.30)^{\phantom{***}}$ & $(3.26)^{\phantom{***}}$ & $(2.98)^{\phantom{***}}$ & $(3.80)^{\phantom{***}}$ & $(3.52)^{\phantom{***}}$ & $(0.67)^{\phantom{***}}$ \\[2pt]
\multicolumn{9}{@{}l}{(a) Short-straddle return} \\
$\alpha$ (\%) & $12.6^{**\phantom{*}}$ & $12.3^{**\phantom{*}}$ & $13.2^{**\phantom{*}}$ & $11.4^{**\phantom{*}}$ & $10.5^{**\phantom{*}}$ & $13.5^{***}$ & $12.2^{***}$ & --$^{\phantom{***}}$ \\
 & $(2.79)^{\phantom{***}}$ & $(2.79)^{\phantom{***}}$ & $(2.73)^{\phantom{***}}$ & $(2.65)^{\phantom{***}}$ & $(2.42)^{\phantom{***}}$ & $(3.15)^{\phantom{***}}$ & $(2.86)^{\phantom{***}}$ &  \\[2pt]
$\beta$, short straddle & $0.26^{\phantom{***}}$ & $0.27^{\phantom{***}}$ & $0.28^{\phantom{***}}$ & $0.22^{\phantom{***}}$ & $0.23^{\phantom{***}}$ & $0.26^{\phantom{***}}$ & $0.25^{\phantom{***}}$ & --$^{\phantom{***}}$ \\
 & $(1.33)^{\phantom{***}}$ & $(1.35)^{\phantom{***}}$ & $(1.42)^{\phantom{***}}$ & $(1.16)^{\phantom{***}}$ & $(1.22)^{\phantom{***}}$ & $(1.39)^{\phantom{***}}$ & $(1.32)^{\phantom{***}}$ &  \\[2pt]
$R^2$ & $0.07$ & $0.07$ & $0.08$ & $0.05$ & $0.05$ & $0.07$ & $0.07$ & -- \\[3pt]
\multicolumn{9}{@{}l}{(b) S\&P~500 half-hour return} \\
$\beta$, $r^{30}_t$ (bp) & $0.05^{\phantom{***}}$ & $0.06^{\phantom{***}}$ & $-0.02^{\phantom{***}}$ & $0.15^{\phantom{***}}$ & $0.12^{\phantom{***}}$ & $-0.01^{\phantom{***}}$ & $0.06^{\phantom{***}}$ & $0.05^{\phantom{***}}$ \\
 & $(0.17)^{\phantom{***}}$ & $(0.20)^{\phantom{***}}$ & $(-0.06)^{\phantom{***}}$ & $(0.55)^{\phantom{***}}$ & $(0.43)^{\phantom{***}}$ & $(-0.05)^{\phantom{***}}$ & $(0.22)^{\phantom{***}}$ & $(0.18)^{\phantom{***}}$ \\[2pt]
$\beta$, $|r^{30}_t|$ (bp) & $-1.72^{***}$ & $-1.70^{***}$ & $-1.88^{***}$ & $-1.66^{***}$ & $-1.68^{***}$ & $-1.40^{**\phantom{*}}$ & $-1.67^{***}$ & $-3.96^{***}$ \\
 & $(-3.30)^{\phantom{***}}$ & $(-3.24)^{\phantom{***}}$ & $(-3.50)^{\phantom{***}}$ & $(-3.17)^{\phantom{***}}$ & $(-3.25)^{\phantom{***}}$ & $(-2.30)^{\phantom{***}}$ & $(-3.26)^{\phantom{***}}$ & $(-6.86)^{\phantom{***}}$ \\[2pt]
$R^2$ & $0.09$ & $0.09$ & $0.11$ & $0.09$ & $0.09$ & $0.06$ & $0.10$ & $0.48$ \\[3pt]
\multicolumn{9}{@{}l}{(c) S\&P~500 daily return and VIX} \\
$\beta$, $r^{\mathrm{cc}}_t$ (\%) & $-1.35^{\phantom{***}}$ & $-0.90^{\phantom{***}}$ & $-0.10^{\phantom{***}}$ & $2.94^{\phantom{***}}$ & $2.44^{\phantom{***}}$ & $3.38^{\phantom{***}}$ & $1.07^{\phantom{***}}$ & $-0.17^{\phantom{***}}$ \\
 & $(-0.23)^{\phantom{***}}$ & $(-0.14)^{\phantom{***}}$ & $(-0.02)^{\phantom{***}}$ & $(0.51)^{\phantom{***}}$ & $(0.42)^{\phantom{***}}$ & $(0.60)^{\phantom{***}}$ & $(0.19)^{\phantom{***}}$ & $(-0.03)^{\phantom{***}}$ \\[2pt]
$\beta$, $\Delta\mathrm{VIX}_t$ & $-2.08^{\phantom{***}}$ & $-0.67^{\phantom{***}}$ & $-3.03^{\phantom{***}}$ & $-0.95^{\phantom{***}}$ & $0.68^{\phantom{***}}$ & $1.80^{\phantom{***}}$ & $-0.71^{\phantom{***}}$ & $-4.98^{\phantom{***}}$ \\
 & $(-0.40)^{\phantom{***}}$ & $(-0.12)^{\phantom{***}}$ & $(-0.58)^{\phantom{***}}$ & $(-0.18)^{\phantom{***}}$ & $(0.13)^{\phantom{***}}$ & $(0.34)^{\phantom{***}}$ & $(-0.14)^{\phantom{***}}$ & $(-1.02)^{\phantom{***}}$ \\[2pt]
$R^2$ & $0.00$ & $0.00$ & $0.00$ & $0.00$ & $0.00$ & $0.00$ & $0.00$ & $0.00$ \\[3pt]
\multicolumn{9}{@{}l}{(d) Joint} \\
$\alpha$ (\%) & $13.0^{**\phantom{*}}$ & $12.7^{**\phantom{*}}$ & $13.7^{**\phantom{*}}$ & $11.8^{**\phantom{*}}$ & $10.9^{**\phantom{*}}$ & $13.7^{***}$ & $12.6^{***}$ & --$^{\phantom{***}}$ \\
 & $(2.76)^{\phantom{***}}$ & $(2.68)^{\phantom{***}}$ & $(2.69)^{\phantom{***}}$ & $(2.62)^{\phantom{***}}$ & $(2.54)^{\phantom{***}}$ & $(3.26)^{\phantom{***}}$ & $(2.86)^{\phantom{***}}$ &  \\[2pt]
$\beta$, short straddle & $0.09^{\phantom{***}}$ & $0.11^{\phantom{***}}$ & $0.10^{\phantom{***}}$ & $0.04^{\phantom{***}}$ & $0.05^{\phantom{***}}$ & $0.18^{\phantom{***}}$ & $0.10^{\phantom{***}}$ & --$^{\phantom{***}}$ \\
 & $(0.37)^{\phantom{***}}$ & $(0.44)^{\phantom{***}}$ & $(0.40)^{\phantom{***}}$ & $(0.16)^{\phantom{***}}$ & $(0.21)^{\phantom{***}}$ & $(0.70)^{\phantom{***}}$ & $(0.38)^{\phantom{***}}$ &  \\[2pt]
$\beta$, $r^{30}_t$ (bp) & $0.10^{\phantom{***}}$ & $0.11^{\phantom{***}}$ & $0.01^{\phantom{***}}$ & $0.17^{\phantom{***}}$ & $0.15^{\phantom{***}}$ & $-0.02^{\phantom{***}}$ & $0.09^{\phantom{***}}$ & --$^{\phantom{***}}$ \\
 & $(0.35)^{\phantom{***}}$ & $(0.39)^{\phantom{***}}$ & $(0.03)^{\phantom{***}}$ & $(0.63)^{\phantom{***}}$ & $(0.56)^{\phantom{***}}$ & $(-0.09)^{\phantom{***}}$ & $(0.33)^{\phantom{***}}$ &  \\[2pt]
$\beta$, $|r^{30}_t|$ (bp) & $-1.37^{*\phantom{**}}$ & $-1.29^{*\phantom{**}}$ & $-1.49^{*\phantom{**}}$ & $-1.51^{**\phantom{*}}$ & $-1.50^{**\phantom{*}}$ & $-0.70^{\phantom{***}}$ & $-1.31^{*\phantom{**}}$ & --$^{\phantom{***}}$ \\
 & $(-2.02)^{\phantom{***}}$ & $(-1.85)^{\phantom{***}}$ & $(-2.10)^{\phantom{***}}$ & $(-2.23)^{\phantom{***}}$ & $(-2.20)^{\phantom{***}}$ & $(-0.80)^{\phantom{***}}$ & $(-1.91)^{\phantom{***}}$ &  \\[2pt]
$\beta$, $r^{\mathrm{cc}}_t$ (\%) & $-6.09^{\phantom{***}}$ & $-5.53^{\phantom{***}}$ & $-4.12^{\phantom{***}}$ & $-2.94^{\phantom{***}}$ & $-3.20^{\phantom{***}}$ & $1.80^{\phantom{***}}$ & $-3.35^{\phantom{***}}$ & --$^{\phantom{***}}$ \\
 & $(-1.08)^{\phantom{***}}$ & $(-0.93)^{\phantom{***}}$ & $(-0.80)^{\phantom{***}}$ & $(-0.56)^{\phantom{***}}$ & $(-0.57)^{\phantom{***}}$ & $(0.33)^{\phantom{***}}$ & $(-0.63)^{\phantom{***}}$ &  \\[2pt]
$\beta$, $\Delta\mathrm{VIX}_t$ & $-2.13^{\phantom{***}}$ & $-0.66^{\phantom{***}}$ & $-2.72^{\phantom{***}}$ & $-1.52^{\phantom{***}}$ & $0.23^{\phantom{***}}$ & $2.69^{\phantom{***}}$ & $-0.69^{\phantom{***}}$ & --$^{\phantom{***}}$ \\
 & $(-0.44)^{\phantom{***}}$ & $(-0.12)^{\phantom{***}}$ & $(-0.60)^{\phantom{***}}$ & $(-0.32)^{\phantom{***}}$ & $(0.04)^{\phantom{***}}$ & $(0.57)^{\phantom{***}}$ & $(-0.15)^{\phantom{***}}$ &  \\[2pt]
$R^2$ & $0.10$ & $0.10$ & $0.11$ & $0.09$ & $0.09$ & $0.08$ & $0.10$ & -- \\[3pt]
\addlinespace[4pt]
\multicolumn{9}{@{}l}{\textit{Panel B: Returns at quoted bid/ask prices}} \\
\addlinespace[2pt]
Mean return (\%) & $10.6^{**\phantom{*}}$ & $10.3^{**\phantom{*}}$ & $11.3^{**\phantom{*}}$ & $9.3^{**\phantom{*}}$ & $8.4^{**\phantom{*}}$ & $11.5^{***}$ & $10.2^{**\phantom{*}}$ & $-0.2^{\phantom{***}}$ \\
 & $(2.74)^{\phantom{***}}$ & $(2.65)^{\phantom{***}}$ & $(2.69)^{\phantom{***}}$ & $(2.57)^{\phantom{***}}$ & $(2.29)^{\phantom{***}}$ & $(3.12)^{\phantom{***}}$ & $(2.82)^{\phantom{***}}$ & $(-0.07)^{\phantom{***}}$ \\[2pt]
$\alpha$ (\%), (a) & $9.9^{**\phantom{*}}$ & $9.6^{**\phantom{*}}$ & $10.5^{**\phantom{*}}$ & $8.7^{*\phantom{**}}$ & $7.8^{*\phantom{**}}$ & $10.8^{**\phantom{*}}$ & $9.5^{**\phantom{*}}$ & --$^{\phantom{***}}$ \\
 & $(2.21)^{\phantom{***}}$ & $(2.19)^{\phantom{***}}$ & $(2.18)^{\phantom{***}}$ & $(2.04)^{\phantom{***}}$ & $(1.80)^{\phantom{***}}$ & $(2.54)^{\phantom{***}}$ & $(2.25)^{\phantom{***}}$ &  \\[2pt]
$\alpha$ (\%), (d) & $10.3^{**\phantom{*}}$ & $9.9^{**\phantom{*}}$ & $11.0^{**\phantom{*}}$ & $9.1^{*\phantom{**}}$ & $8.2^{*\phantom{**}}$ & $11.0^{**\phantom{*}}$ & $9.9^{**\phantom{*}}$ & --$^{\phantom{***}}$ \\
 & $(2.21)^{\phantom{***}}$ & $(2.12)^{\phantom{***}}$ & $(2.18)^{\phantom{***}}$ & $(2.04)^{\phantom{***}}$ & $(1.93)^{\phantom{***}}$ & $(2.64)^{\phantom{***}}$ & $(2.27)^{\phantom{***}}$ &  \\[2pt]
\bottomrule
\end{tabular}
\begin{tablenotes}[flushleft]
\footnotesize
\item \textit{Notes:} OLS regressions of the daily return of the sign strategy (in \% of the premium) on the
\rheSzDays{} post-warmup option days. Columns: the six pooled intraday forecasts on the full predictor set, their
equal-weight mean return (Mean), and the unconditional short straddle. Regressors: (a) the return of the unconditional
short straddle at mid-quotes, $-R_t$ (\%); (b) the S\&P~500 log return from the 15:30 underlying print to the close,
$r^{30}_t = \log(S^{\mathrm{close}}_t/S_t)$, and its absolute value, in basis points; (c) the S\&P~500
close-to-close log return $r^{\mathrm{cc}}_t$ (\%) and the daily change in the VIX close (points); (d) all of them.
The four non-traded regressors of (b)--(d) are demeaned, so the intercept $\alpha$ is the mean return net of the
exposure to the short straddle, a traded return; in (b) and (c) it is the mean return. In parentheses, $t$-statistics with
Newey--West standard errors (Bartlett kernel, bandwidth $S = \rheHACS$); stars from fixed-$b$ critical values
(Section~\ref{sec:rh_methods_tests}). $^{*}$, $^{**}$ and $^{***}$ denote significance at the 10\%, 5\% and 1\% levels.
\end{tablenotes}
\end{threeparttable}
\end{table}
      % tab:rh_econ_riskreg
% Text elsewhere that this draft contradicts or qualifies (not edited here):
%   07_economics.tex, Results, "Concentration of economic gains": "($R_t = -1$, corresponding to settlement pinning
%     near the strike; ...)". The close settles between the two strikes about as often as the implied distribution
%     predicts (item "Worthless expiries" below); proposed: "($R_t = -1$: the close settles between the two strikes;
%     ...)".
%   09_conclusion.tex, Maisch comparison: "ours come from timing the variance premium of the closing half hour" (parked
%     2026-10-25). Supported only in a narrow sense: the gain over the short straddle is earned on the long days the
%     forecast picks, and in this sample almost all of it on about ten of them (closing paragraph below).
%   05_methods.tex, Back-transformation: "a downward Jensen inequality bias that we leave uncorrected". It is there in
%     ratio terms and tilts the sign rule short, but a causal correction lowers every full-set Sharpe ratio (item
%     "Calibration" below); the text may say so.
%   07_economics.tex, Results, "Unconditional short-straddle benchmark": "earns no statistically detectable variance
%     risk premium". Consistent; note that the forecast target RV_t (ES futures, one-minute returns) averages
%     \rhedUncRVIVMean{} of IV_t^2 while the squared half-hour move of the cash index that settles the straddle
%     averages \rhedMoveSqIV{} of it (item "Calibration" below).

\subsection{Where the profits come from}
\begin{itemize}
\item \textbf{A short straddle that turns long on a fifth of the days.} The sign strategies of the six pooled intraday
  forecasts on the full predictor set hold a long straddle on \rhedLongNMin{} to \rhedLongNMax{} of the \rheSzDays{}
  post-warmup days (\rhedLongPctMin--\rhedLongPctMax\%; Table~\ref{tab:rh_econ_longshort}). Because
  \[
    r_t = q_t R_t = -R_t + 2R_t\,\mathbf{1}\{q_t = +1\},
  \]
  each strategy's return is that of the unconditional short straddle plus a gain $\Delta_t = 2R_t\mathbf{1}\{q_t = +1\}$
  over it, and that gain is earned on the long days only.
  \begin{itemize}
  \item The unconditional short straddle earns \rhedUncShortMean\% of the premium a day at mid-quotes; the gain on the
    long days adds \rhedGainMin--\rhedGainMax\%, which is \rhedGainShareMin--\rhedGainShareMax\% of the strategies'
    mean daily return of \rhedMeanMin--\rhedMeanMax\%.
  \item On its long days the straddle returns \rhedLongMidMin--\rhedLongMidMax\% of the premium on average at
    mid-quotes and \rhedLongBaMin--\rhedLongBaMax\% bought at the ask. On the remaining days the short straddle earns
    \rhedShortMidMin--\rhedShortMidMax\% at mid-quotes, against \rhedUncShortMean\% on all days: the long days are the
    days on which the short straddle loses.
  \item The long days account for \rhedLongShareMin--\rhedLongShareMax\% of the summed mid-quote return
    (\rhedLongShareXMin--\rhedLongShareXMax\% at quoted bid/ask prices). On the HAR benchmark the same model classes go
    long on \rhedHarLongNMin{} to \rhedHarLongNMax{} days, on which the straddle returns only
    \rhedHarLongMidMin--\rhedHarLongMidMax\% (gain over the short straddle \rhedHarGainMin--\rhedHarGainMax\% a day).
  \end{itemize}
\item \textbf{No better call of the premium's sign, but the large moves.} The position has the sign of $\RV_t - \mathrm{IV}_t^2$ on \rhedHitMin--\rhedHitMax\% of the days, against
  \rhedUncShortHit\% for the unconditional short straddle, so the forecasts are no better than always selling at
  calling the sign of the realized premium. What they add is in the long days: realized variance exceeds implied
  variance on \rhedLongHitMin--\rhedLongHitMax\% of them, against \rhedUncRVAboveIV\% of all days, yet the long
  straddle has a positive return on only \rhedLongPosMin--\rhedLongPosMax\% of its long days, so its mean return on
  those days comes from a few large moves.
\item \textbf{Forecast quintiles} (Table~\ref{tab:rh_econ_quintiles}). Sorted by the forecast variance gap
  $\log(\hat\sigma^2_t/\mathrm{IV}_t^2)$, the mean ratio of realized to implied variance rises from quintile to
  quintile for all \rhedQMonoN{} full-set forecasts, from \rhedQOneRVIV{} in the lowest to \rhedQFiveRVIV{} in the
  highest (mean of the six), and for \rhedHarQMonoN{} of the six HAR-benchmark forecasts (\rhedHarQOneRVIV{} to
  \rhedHarQFiveRVIV): the forecasts rank the closing half hours by their realized variance relative to the implied.
  \begin{itemize}
  \item The return of a long straddle is negative in each of the four lower quintiles (\rhedQLowRMin{} to
    \rhedQLowRMax\% on average) and \rhedQFiveR\% in the highest, where \rhedQFiveLong\% of the days are long days;
    for the HAR-benchmark forecasts it is \rhedHarQFiveR\%. The sign rule's threshold, $\hat\sigma^2_t =
    \mathrm{IV}_t^2$, falls between the fourth quintile (forecast ratio \rhedQFourFc) and the fifth
    (\rhedQFiveFc).
  \item The highest quintile earns \rhedQFiveShare\% of the summed mid-quote return of the six full-set strategies
    (\rhedQFiveShareRngMin--\rhedQFiveShareRngMax\% across them), against \rhedHarQFiveShare\% for the HAR benchmark.
    Across the four lower quintiles, where the strategies are short, the return is not monotone in the forecast.
  \end{itemize}
\item \textbf{About ten days carry the gain over the short straddle} (Table~\ref{tab:rh_econ_bestdays}). The six
  full-set forecasts share \rhedBestAllSix{} of their ten best days and \rhedWorstAllSix{} of their ten worst; the
  best days are long straddles on large closing moves, the worst are short straddles on large closing moves.
  \begin{itemize}
  \item Without each forecast's ten best days, the Sharpe ratio gain over the unconditional short straddle falls from
    \rhedDShortAllMin--\rhedDShortAllMax{} to \rhedDShortExBestMin{} to \rhedDShortExBestMax{} (negative for
    \rhedDShortExBestNeg{} of the six); the ten days earn \rhedBestShareMin--\rhedBestShareMax\% of the summed return
    gain over the short straddle, so the other long days lose on balance.
  \item The largest, \rhedTopDate, returned \rhedTopR\% of the premium on a closing-half-hour move of
    \rhedTopMove\% (realized variance \rhedTopRVIV{} times the implied), \rhedTopShareMin--\rhedTopShareMax\% of each
    full-set strategy's summed return.
  \item Against the HAR twin the gain is less concentrated: without the forecast's ten best days the Sharpe ratio gain
    falls from \rhedDHarAllMin--\rhedDHarAllMax{} to \rhedDHarExBestMin--\rhedDHarExBestMax, and those days carry
    \rhedHarBestShareMin--\rhedHarBestShareMax\% of the summed return gain over the twin.
  \end{itemize}
\item \textbf{The days are not scheduled announcements.} Of the \rhedBestUnion{} days in the union of the six
  forecasts' best days, \rhedBestFomcN{} is an FOMC statement day and \rhedBestNoTagN{} carry no tag in our event
  calendar, \rhedBestAfterRelN{} of them after \rhedRelEndText, when its macro-release dates end (its FOMC dates are
  complete).
  \begin{itemize}
  \item On the \rhedFomcN{} FOMC statement days the closing straddle returns \rhedFomcR\% on average
    ($t = \rhedFomcRT$, the days treated as independent); the full-set forecasts are long on \rhedFomcLong\% of them
    and earn \rhedFomcStrat\% a day, against \rhedFomcOtherStrat\% on the other days.
  \item On the \rhedMeN{} month ends the straddle returns \rhedMeR\% ($t = \rhedMeRT$) and the squared closing move
    averages \rhedMeMoveSqIV{} times the implied variance; the forecasts are long on \rhedMeLong\% of them, earn
    \rhedMeStrat\% a day, and \rhedMeShare\% of their summed return comes from these days (\rhedWorstMeN{} of the worst days
    are month ends).
  \item On the \rhedOpexN{} monthly SPX expirations the straddle returns \rhedOpexR\% ($t = \rhedOpexRT$); the
    forecasts are long on \rhedOpexLong\% of them and earn \rhedOpexShare\% of their summed return there, a gain the
    unconditional short straddle shares.
  \end{itemize}
\item \textbf{Calibration and the short tilt} (Table~\ref{tab:rh_econ_longshort}, Panel D). In levels the full-set
  forecasts are close to unbiased on the option days (mean forecast over mean realized variance
  \rhedCalLevelMin--\rhedCalLevelMax), but in ratios they are low: the mean of $\RV_t/\hat\sigma^2_t$, the scale that
  minimizes mean QLIKE, is \rhedCalScaleMin--\rhedCalScaleMax, the downward bias of squaring a level forecast
  (Section~\ref{sec:rh_methods_setup}).
  \begin{itemize}
  \item Multiplying each forecast by a causal estimate of that scale, the mean of $\RV_{t'}/\hat\sigma^2_{t'}$ over the
    earlier option days $t' < t$, changes \rhedBcFlipMin{} to \rhedBcFlipMax{} positions, all from short to long.
    The straddle's mean return on the changed days is \rhedBcFlipRMin{} to \rhedBcFlipRMax\%, and the Sharpe ratio
    falls to \rhedBcShMin--\rhedBcShMax{} at mid-quotes and \rhedBcShXMin--\rhedBcShXMax{} at quoted bid/ask prices
    (below the uncorrected ratio for \rhedBcShDownN{} of the six; the correction raises it for \rhedHarBcShUpN{} of
    the six HAR-benchmark forecasts). This is a
    diagnostic, not a strategy of the paper.
  \item A reading consistent with this: a straddle pays the absolute move of the index, so a long position is
    profitable when the expected realized volatility, not variance, exceeds the implied volatility. The sign rule
    $\hat\sigma^2_t > \mathrm{IV}_t^2$ is the same as $\hat y_t\, s_{t,c} > \mathrm{IV}_t$, a comparison of the level
    forecast with the implied volatility, which the uncorrected forecast makes.
  \item The forecast target and the straddle's settlement differ: $\RV_t$, the realized variance of the futures from
    one-minute returns, averages \rhedUncRVIVMean{} of the implied variance (median \rhedUncRVIVMed), while the squared
    log move of the cash index from $S_t$ to $S^{\mathrm{close}}_t$ averages \rhedMoveSqIV{} of it, and
    \rhedMoveSqRV{} times the mean of $\RV_t$. Measured on the move that settles the straddle, the implied variance
    carries no premium on average in our sample, which agrees with the short straddle's Sharpe ratio of
    \rheShortSh.
  \end{itemize}
\item \textbf{Worthless expiries are not pinning.} $R_t = -1$ exactly when the close settles between the two strikes,
  $K_P \le S^{\mathrm{close}}_t \le K_C$; on \rhedGapFiveN{} of the \rheSzDays{} days the two strikes are 5 index points
  apart, against a median one-standard-deviation implied move over the half hour of \rhedMoveSdPts{} points.
  \begin{itemize}
  \item The straddle expires worthless on \rhedWorthN{} days (\rhedWorthPct\%). A normal half-hour log move with the
    implied volatility $\mathrm{IV}_t$ predicts \rhedWorthExpIV{} such days (\rhedWorthExpIVPct\%; the difference is
    \rhedWorthZ{} standard deviations), the empirical distribution of the earlier days' moves scaled by
    $\mathrm{IV}_t$ predicts \rhedWorthExpEmp, and the close is uniformly placed between adjacent strikes (mean distance
    to the nearest strike \rhedStrikeDist{} points, \rhedStrikeDistUnif{} under a uniform; Kolmogorov--Smirnov
    $D = \rhedKsD$, 5\% critical value \rhedKsCV).
  \item The full-set forecasts are short on \rhedWorthShortMin--\rhedWorthShortMax{} of these days, on which the short
    straddle keeps its whole premium; these days supply \rhePoolWorthlessShareMed\% of the positive returns of the
    median pooled forecast (Section~\ref{sec:rh_econ_results}).
  \item Monthly SPX expirations are the exception: \rhedWorthOpexN{} of the \rhedOpexN{} expire worthless against
    \rhedWorthOpexExp{} expected ($\rhedWorthOpexZ$ standard deviations), and the squared closing move averages
    \rhedOpexMoveSqIV{} of the implied variance. Outside them the share is \rhedWorthOtherPct\% against
    \rhedWorthOtherExpPct\% expected.
  \end{itemize}
\item \textbf{Risk exposures} (Table~\ref{tab:rh_econ_riskreg}; $t$-statistics with Newey--West standard errors at
  bandwidth $S = \rheHACS$ and fixed-$b$ critical values, Section~\ref{sec:rh_methods_tests}).
  \begin{itemize}
  \item On the short-straddle return the full-set strategies load \rhedRegBetaShortMin--\rhedRegBetaShortMax{}
    ($t$ from \rhedRegBetaShortTMin{} to \rhedRegBetaShortTMax; $R^2$ \rhedRegRsqAMin--\rhedRegRsqAMax), with
    $\alpha = \rhedRegAlphaAMin$--\rhedRegAlphaAMax\% of the premium a day at mid-quotes (fixed-$b$ $p$ from
    \rhedRegAlphaAPMin{} to \rhedRegAlphaAPMax) and \rhedRegAlphaAXMin--\rhedRegAlphaAXMax\% at quoted bid/ask prices
    ($p$ from \rhedRegAlphaAXPMin{} to \rhedRegAlphaAXPMax). The equal-weight mean of the six returns has $\alpha = \rhedRegAvgAlphaA\%$
    ($p = \rhedRegAvgAlphaAP$) at mid-quotes and $\rhedRegAvgAlphaAX\%$ ($p = \rhedRegAvgAlphaAXP$) at quoted
    bid/ask prices.
  \item They are short volatility, less so than the short straddle: a basis point of absolute S\&P~500 move over the
    closing half hour changes their return by \rhedRegBetaAbsMin{} to \rhedRegBetaAbsMax\% of the premium ($t$ from
    \rhedRegBetaAbsTMin{} to \rhedRegBetaAbsTMax); for their mean return the coefficient, \rhedRegAvgBetaAbs\%, is
    \rhedRegAbsRatio\% of the short straddle's \rhedRegShortBetaAbs\%. The sign of the move does not matter
    (coefficients \rhedRegBetaMoveMin{} to \rhedRegBetaMoveMax).
  \item The S\&P~500 daily return and the change in the VIX explain none of their return ($R^2$
    \rhedRegRsqCMax), and with all regressors $\alpha$ is \rhedRegAlphaDMin--\rhedRegAlphaDMax\% at mid-quotes and
    \rhedRegAlphaDXMin--\rhedRegAlphaDXMax\% at quoted bid/ask prices.
  \item Without each forecast's ten best days, $\alpha$ on the short straddle is \rhedRegAlphaAExBestMin{} to
    \rhedRegAlphaAExBestMax\% ($p \ge \rhedRegAlphaAExBestPMin$). The standard errors allow for serial correlation, not
    for a return distribution whose mean rests on ten days. The HAR-benchmark strategies have
    $\alpha = \rhedHarRegAlphaAMin$--\rhedHarRegAlphaAMax\% ($p \ge \rhedHarRegAlphaAPMin$).
  \end{itemize}
\end{itemize}

\paragraph{Where the profits come from.} The strategy is a short straddle on most days, and the short straddle
by itself earns nothing detectable in our sample (\rheShortSh{} at mid-quotes, \rheShortShX{} at quoted bid/ask
prices). The profit comes from the days on which the full-set forecasts switch to a long straddle: these are the
closing half hours that the forecasts rank highest in realized relative to implied variance, the ranking is monotone
across the forecast quintiles, and on these days the straddle earns \rhedLongMidMin--\rhedLongMidMax\% of its premium
on average (\rhedQFiveR\% in the top forecast quintile, against \rhedHarQFiveR\% on the HAR benchmark). The strategy is not a disguised short-volatility
position: it has less than half of the short straddle's exposure to the size of the closing move, almost none to its
return, and none to the market or the VIX. But the payoff of the ranking is concentrated in about ten days with closing
moves several times the implied, the same days for all six forecasts, and only \rhedBestFomcN{} of the \rhedBestUnion{} an FOMC
statement day;
without them the gain over the short straddle and the regression $\alpha$ disappear. The forecasts thus time the sign
of the closing-half-hour variance premium in a narrow sense, by stepping out of the short straddle on a minority of
days, and in this sample of \rheSzDays{} days that timing paid through a handful of large closing moves that they
anticipated, not through a steady premium.
